% =====================================================================
%  Harald Høffding, Shorter Works. Third Series (1913)
%  Translated from the Danish: Mindre Arbejder. Tredie Række.  København og Kristiania: Gyldendalske Boghandel, Nordisk
%  Forlag, 1913.
%
%  TRANSLATION -- GENERATED; DO NOT EDIT BY HAND.  Built by
%    python3 ebook-method/mindre-arbejder-3/tr_parts.py build mindre-arbejder-3 --write
%  from this template and the ten parts in ebook-method/mindre-arbejder-3/tr/*.texfrag (GLOSSARY.md, TRANSLATOR-BRIEF.md).
%  Made from the collated transcription texts/hoeffding/mindre-arbejder-3/transcription.tex; [N] in the margin = the page
%  of the original.
%
%  STATUS: COMPLETE (2026-10-26).
%
%  METHOD (TRANSLATION-PLAYBOOK.md).  Translated 2026-10-26 by ten translators, one part each (a pilot, then batches of
%  five and four), cut at essay openings, from the collated transcription, with GLOSSARY.md (title, heads and the terms
%  marked there approved by Hans Halvorson) and TRANSLATOR-BRIEF.md.  tr_parts.py build checks every page marker once and
%  in order, the footnotes and paragraph breaks page by page, and the counts of italics, small capitals, quotation
%  marks and verse beside the Danish.
%  HØFFDING'S OWN ENGLISH IS THE TRANSLATION WHERE IT EXISTS (Hans Halvorson, 2026-10-26).  Three pieces were printed in
%  his English: "A Philosophical Confession" (J. Phil., Psych. and Sci. Methods 2, 1905, pp. 85-92) for *En filosofisk
%  Bekendelse*, pp. 18-27; "On the Relation between Sociology and Ethics" (American Journal of Sociology 10, 1905,
%  pp. 672-685, a synopsis of the London lecture and his reply to the discussion) for pp. 68-77; "The Influence of the
%  Conception of Evolution on Modern Philosophy" (Darwin and Modern Science, Cambridge 1909, ch. XXII) for pp. 202-228.
%  In these his sentences are used word for word where they say what the Danish says (spelling made American, OCR slips
%  of the 1905 journals mended at their page images); what the 1913 Danish adds is translated in his terms; what his
%  English has and the Danish lacks is left out; where the Danish says something else, the Danish is translated.  The
%  translators' logs (tr/anchor-*.tsv) class the 106 paragraphs of these pieces: 12 his English throughout, 67 his
%  English with Danish additions or omissions, 27 mixed, none wholly from the Danish.  Share of the English text's
%  6-word sequences found in his English: Confession 82%, Sociology 77%, Darwin 67%.  The heads of the three pieces are
%  translated from the Danish heads (his English titles differ for two).  His quotations of English authors (Darwin,
%  Spencer, James) are given in his English; elsewhere in the book, Høffding's Danish renderings of English, French,
%  German and Greek authors are translated from his Danish, and quotations in other languages stay as printed.
%  ODD SPOTS.  The translators listed 95 doubtful spots of the Danish (tr/odd-*.tsv); all were read at line crops of the
%  scan with 7 more (tr/odd-verdicts.tsv): 33 transcription faults (8 of them false paragraph breaks) and 52 misplaced dashes were mended in the Danish
%  (patches.LAST, DASH_TO_PREV; see the header of transcription.tex) and the English mended to match; the print's own
%  misprints are translated in their evident sense, foreign names and titles left as printed.
%  Not done: a second, independent reading of the English against the Danish.
% =====================================================================
\documentclass[12pt,a4paper,openany]{book}
\usepackage[utf8]{inputenc}
\usepackage[T1]{fontenc}
\usepackage{libertinus}
\usepackage{libertinust1math}
\usepackage[german,french,english]{babel}
\usepackage{textalpha}
\usepackage[protrusion=true,expansion=true]{microtype}
\usepackage{setspace}
\usepackage[margin=1.3in]{geometry}
\usepackage{fancyhdr}
\usepackage{marginnote}
\usepackage{graphicx}
\usepackage{hyperref}
\hypersetup{pdftitle={Shorter Works. Third Series}, pdfauthor={Harald Høffding}, hidelinks}
\pagestyle{fancy}
\fancyhf{}
\fancyhead[LE,RO]{\thepage}
\fancyhead[C]{\textit{H. Høffding: Shorter Works. Third Series (1913)}}
\fancypagestyle{plain}{\fancyhf{}\fancyhead[LE,RO]{\thepage}}
\setlength{\headheight}{15pt}
\setstretch{1.3}
\DeclareUnicodeCharacter{0254}{\reflectbox{c}}   % ɔ (id est)
\renewcommand{\thefootnote}{\arabic{footnote})}
\newcommand{\opage}[1]{\marginnote{\footnotesize\textit{[#1]}}}
\newcommand{\whiteline}{\par\vspace{\baselineskip}\par}
\newcommand{\essay}[2][]{\chapter*{#2}\addcontentsline{toc}{chapter}{#1}}
\newcommand{\undertitel}[1]{\begin{center}\small #1\end{center}}
\newcommand{\streg}{\begin{center}\rule{3em}{0.4pt}\end{center}}
\newcommand{\stjerner}{\begin{center}*\hspace{0.4em}*\\[-0.9ex]*\end{center}}
\newcommand{\ix}[3]{\par\noindent\makebox[2em][r]{#1}\hspace{0.6em}#2\dotfill #3}
\pdfstringdefDisableCommands{\renewcommand{\footnote}[2][]{}\renewcommand{\opage}[1]{}}

\begin{document}

\frontmatter
\begin{titlepage}\centering\vspace*{2cm}
{\large HARALD HØFFDING}\par\vspace{2cm}
{\Huge SHORTER WORKS}\par\vspace{1ex}{\large THIRD SERIES}\par\vspace{1ex}
{\small [Mindre Arbejder. Tredie Række]}\par\vfill
{\small COPENHAGEN AND KRISTIANIA}\\[0.5ex]{\small GYLDENDALSKE BOGHANDEL}\\[0.5ex]{\small NORDISK FORLAG}\\[1ex]1913\\[3em]
{\small Translated from the Danish}
\end{titlepage}
\chapter*{CONTENTS}
\begin{center}\rule{2cm}{0.4pt}\end{center}
\noindent\hfill Page
\ix{1.}{Spiritual Culture (1912)}{1}
\ix{2.}{A Philosophical Confession (1904)}{18}
\ix{3.}{The Concept of Will (1905; 1912)}{28}
\ix{4.}{Humane Ethics (1901)}{54}
\ix{5.}{Sociology and Ethics (1904)}{68}
\ix{6.}{A Case History (1901)}{78}
\ix{7.}{Plato's Symposium (1905)}{83}
\ix{8.}{Democritus and Plato (1909)}{94}
\ix{9.}{Epictetus and Christianity (1912)}{102}
\ix{10.}{Marcus Aurelius (1912)}{107}
\ix{11.}{Rousseau and the Nineteenth Century (1912)}{120}
\ix{12.}{Rousseau and Religion (1912)}{142}
\ix{13.}{Goethe and Philosophy (1907)}{162}
\ix{14.}{At the Schiller Festival (1909)}{176}
\ix{15.}{Napoleon and Philosophy (1906)}{184}
\ix{16.}{Kierkegaard and Nietzsche (1911)}{189}
\ix{17.}{At the Krøyer Commemoration (1909)}{196}
\ix{18.}{Darwin and Philosophy (1909)}{202}
\par\bigskip
\begin{center}\rule{4cm}{0.4pt}\end{center}
\par\medskip\noindent{\small The page numbers are those of the original, given in the margin as [N].  The contents and
the heads of the pieces differ in places in the original; both are translated as printed.}
\mainmatter
% --- p. 1 ---
\setcounter{footnote}{0}%
\opage{1}\essay[Spiritual Culture (1912)]{SPIRITUAL CULTURE\footnote[1]{Lectures in Kristiania, March 1912.}.}
\streg

A\textsc{ll} that in our day oppresses the human intellectual life
can be derived from a single cause, which the development of
intellectual life itself through the last centuries has brought with it. It
is of greater importance to seek to find this cause than to
complain of the externality of the age, its absorption in outward
relations and tasks, of the excessive role that technology plays
in all fields, or of the extravagant eagerness to get on in the
world at any price. For the great energy that is unfolded
outwardly need not lead to the inner life suffering harm. For my
own part I must state that on my way I constantly meet many
testimonies of spiritual sense, of an urge to immerse oneself in
the tasks of the inner life and of earnest work to satisfy this
urge, testimonies that in my eyes are of all the greater value
as the conditions of intellectual life in many ways hamper the
sense, the urge and the work alike.

The hampering cause I find, as I said, in something that
spiritual development has brought with it, and which sets our
time its most important problem. It is the intrusion of the
division of labor into the spiritual domain that conditions the
fundamental problem of spiritual culture in our day, just as its
appearance in the domain of material culture sets the social
problem in the narrower
% --- p. 2 ---
\setcounter{footnote}{0}%
\opage{2}sense. In a wider sense the problem of spiritual culture too is
a part of the social problem.

It is the dividedness within intellectual life that causes the
difficulties. Science, art and religion stand more and more as
independent, perhaps even contending powers over against one
another, divide the mind in the individual and divide society
into circles that often do not understand one another. Since
the end of the Middle Ages this opposition has become more and
more marked. In the great age of the Middle Ages thought,
imagination and feeling worked together harmoniously, and the
world of the spirit formed a whole that found its classical
expression in Dante's great work. Such a work, in which all
spiritual interests could at once find expression and
satisfaction, is unthinkable in modern times.

Rousseau was the one who first saw the problem. It lay at
the root of his watchword of returning to nature, a watchword to
which he himself has given the best interpretation in the demand
that man shall again become a unity: Rendez l’homme
un! Schiller felt the sting of the problem, as is noticeable in his
``Briefe über ästhetische Erziehung''. Romanticism believed it had
found the way back to the unity of intellectual life, in proclaiming
that the highest science, the great art and the true
religion were in their essence one, only different forms of one
and the same spiritual content. This proclamation awakened in
the various fields many germs to life and unfolding, but partly
it moved too much on the dark paths of imagination and
enthusiasm, partly it was quickly taken into the service of
political and ecclesiastical reaction\footnote[1]{Cf. on various circumstances contributing to this: \textit{Mindre Arbejder. Første Række.} p. 202---209.}. And in the course of the
nineteenth century the ways parted still more than before. Science
--- as natural science, historical research and philosophical inquiry
--- goes ever more surely and self-consciously its own ways,
puts its questions more and more and demands its answers
according to points of view that neither religion nor art knows.
And art springs more and more from an urge to a life in
contemplation and in mood, different from the life
% --- p. 3 ---
\setcounter{footnote}{0}%
\opage{3}that is led where intellectual or religious interest rules. Religion,
finally, becomes more and more clear that it owes its existence
neither to thought nor to the urge for images\footnote[1]{The Catholic Church does not acknowledge this at all. It finds the most important cause of Modernism (the tendency that would rely exclusively on psychological motives and historical investigations in the treatment of the religious problem) in this, that in one's blindness one no longer acknowledges the validity of the proofs of God's existence as the great scholastics of the Middle Ages set them forth. An interesting and substantial account from the Catholic point of view is given by \textit{Anton Gisler: Der Modernismus.} (Köln 1912).}, and since it
nevertheless dare not give up its time-honored forms and institutions,
it clings to these as an emergency aid, closing its eyes,
consciously or unconsciously, to the other sides of intellectual life.

For my own part this problem has followed me since my
youth. It was precisely in the last third of the nineteenth
century --- here in Denmark especially through Kierkegaard's
appearance ---, that the opposition between the various tendencies
of intellectual life became so strongly marked. I dare not say
that I have reached any solution that can have general
significance. But there are thoughts that have ever more pressed
forward and won confirmation for me, and which I shall here seek
to develop. And there are tracks that I consider worth following.


Rousseau already saw that the opposition between nature
and culture falls away when culture itself stands as an unfolding
of nature. Thus all forms of spiritual culture, science, art and
religion, must in the end point back to the human personality,
its urge and its capacity. From this source culture is renewed
again and again in the course of history, and there is no other
source of renewal for it. It is personal powers that work in all
three spiritual domains, powers that cannot work unless there
constantly also stirs a personal urge to release them. It is then
a standard for a tendency in spiritual culture whether it can
demonstrate its connection with the capacity and the urge of
personal life. I therefore consider the three forms of spirit from
this point of view.

% --- p. 4 ---
\setcounter{footnote}{0}%
\opage{4}In my Helsingfors lectures, which have appeared in Swedish
under the title ``Personlighetsprincipen i Filosofin'' (Stockholm
1911), I have treated some of the questions discussed in what
follows from another side, starting from the principle of personality
and showing its relation to and significance for science,
ethics and religion. I now seek in what follows, conversely, to
show how science, art and religion point back to the principle
of personality. The two modes of treatment supplement each
other, and by the nature of the case related thoughts will emerge
through them.

\begin{center}I\end{center}


Scientific culture presupposes, like all true culture,
independent, personal work with the problems of life and of
existence, the great as well as the small. It cannot come into
being through mere tradition and instruction. Strictly speaking it
is only the great pioneers of genius who possess true scientific
culture; they produce or discover the presuppositions, the
foundation, on which the following ages work. To draw conclusions
or make special applications of presuppositions that have already
been laid down is not so original a mental work as laying down
these presuppositions themselves at first hand. But it is
nonetheless always a work that demands independence and into
which personality must be put, a work that stands high above
mere appropriation and assent. This last again stands above pure
reception in trust in the correctness of tradition.

The further a piece of knowledge stands from the original
source itself, the more it stands as something foreign to
personality, perhaps even as something hostile, something that
conflicts with what personality itself in its spontaneous yearning
seeks to hold fast, and therefore as something it seeks to push
aside or to forget.

Here lies a great problem. Not all can draw knowledge
from its first sources, and yet knowledge has value only as
self-won. It is incumbent on those who are to spread the knowledge
they themselves have
% --- p. 5 ---
\setcounter{footnote}{0}%
\opage{5}acquired in a more self-active way, to awaken a work of
thought that makes the appropriation as independent as possible.
Only then does the great wonder arise that is not merely the
beginning of wisdom, and only then can the significance of science
for personal life become clear. For this significance does not rest
on the mere results, which can perhaps be communicated quite
apart from the ways by which they were won, and which are
often valued only on account of their practical significance. But it
rests first and foremost on this, that science is the human mind's
own work, which it will not in the long run be able to deny.

The whole of modern philosophy is in a certain way a great
labor to demonstrate the connection of science with personality.
The question here has first and foremost been this: what mental
powers does science lay claim to, --- what presuppositions does it
rest on, --- and what demands does it make on us? --- Kant became
the real founder of the theory of knowledge by raising these
questions. And he found that what is at work in scientific work
is nothing other than what is at work in all mental life, from the
lowest to the highest. Everywhere a striving manifests itself to
gather and order the manifold influences to which a living being
is subject. In the coming into being of the simplest sense-image
as in the highest art, in elementary perception as in the most
stringent inference, one and the same fundamental form asserts
itself, in that coherence and unity are striven after. It is the
spontaneous mode of working of intellectual life that we here
confront, and it is the philosopher's task to become clearly
conscious of it and to point out the more special forms in which
it appears. It is then precisely the task of science to work at
asserting the unity of personality over against the multiplicity
that existence presents, and the demands that science makes on
us therefore do not fall outside the domain of personality, do not
% print: denfra (misprint)
come to it from there as from an alien power. It is a significant side of
personality's own nature that is satisfied through scientific
work.

While the philosophical tendency stemming from Kant, the
so-called criticism, lays the weight especially on the forms
% --- p. 6 ---
\setcounter{footnote}{0}%
\opage{6}in which the work of thought proceeds, another
epistemological tendency, the so-called pragmatism, emphasizes the
interests and ends that at different times and in different domains
set the work of thought going. All science develops first out of
practical interests, out of a striving to find means for the
preservation and development of life. Later a purely theoretical
interest may arise, and it then determines which questions are to
be raised and how they are to be formed. This tendency too, then,
points back to personality, to its ends, its interests, and that, at
least apparently, in a more direct way than criticism.

The difference between the two tendencies is chiefly this, that
the one lays the weight on the constant forms of the life of
thought, the other on its constantly changing interests and tasks.
The one is inclined to believe that there are certain forms of
thought (``categories'') that can be found once for all and that
stand firm once for all; the other is inclined to believe that with
the interests the forms too may without further ado change.

Interest and form cannot be separated in the life of thought.
There is one great interest that sustains it, the interest in truth,
and no other interest can in the long run oppose it or lead around
it. And what truth is, the truth that a human being is to be able
to attain, --- that follows from man's own nature, from its
necessary striving after unity and connection. What man seeks,
when he seeks truth, is the greatest possible connection between
the greatest possible number of observations. The work of
knowledge therefore goes in two directions: to gather or produce
as many observations as possible, and to demonstrate as intimate
and law-governed a connection between them as possible. That is
the great ideal that has once and for all appeared before the race,
and which it cannot deny. It leads to ever renewed striving and
work. Every claim to have found truth in one field or another is
judged by the standard that this ideal contains\footnote[1]{It is on account of this line of thought, which was first set forth in my program on ``Filosofiske Problemer'' (1902), later in ``Den menneskelige Tanke'' (1910) p. 257---265, that in the Catholic work mentioned above (Anton Gisler: \textit{Der Modernismus.} 1912. p. 406 f.) I am counted among the ``Modernists''; the decisive thing is ``the fateful error'' that truth is no longer defined as the agreement of knowledge with its object. Yet I am there placed closer to pragmatism than I in reality stand. Cf. ``Den menneskelige Tanke'' p. 270---273.}.

% --- p. 7 ---
\setcounter{footnote}{0}%
\opage{7}Just as in ethical self-control we bend our momentary and
elementary impulses under the law that is given with our ethical
ideal, so in intellectual self-control we bend our favorite opinions,
fancies and wishes under the great law of truth. As a great natural
scientist has said: ``Should we not be ashamed in our inmost heart,
if we caught ourselves wanting a truth other than the real one?'' ---

But it is not thereby given that this standard can be
applied to everything that stirs in personal life. The forms with
which science can work at a given time do not contain everything
that personality experiences and learns by experience. Science
seeks a great, comprehensive connection, within which personality
too is to find its place; but the content of personality can be so
rich and manifold that that connection comes to stand as abstract
and empty over against it. Content and extension here come, as
so often, to stand in inverse relation to each other. Every
distinctive personality has its own special, individual experiences,
which no other can have, and is therefore more than an example
of a general law. There is that which we must experience and
learn by experience at first hand in order to be human at first
hand. Is there any other form of spiritual culture in which these
purely personal experiences and practical experiences can find
expression?

This question leads over to a consideration of art and of
religion.

\begin{center}II\end{center}


Science is a work of personal labor, and only in so far as it is
so is it a link in true spiritual culture. Personal life itself, out of
which science springs where it is practiced at first hand, is one of
the
% --- p. 8 ---
\setcounter{footnote}{0}%
\opage{8}objects of science, belongs within the great connection that it is
the task of science to explore; but here there is always something
that cannot be wholly converted into scientific form. The most
special and individual experience of life, the one that must
constantly be had anew, by the same human being or in a new
generation, always under new conditions, --- it stands always with
a certain independence over against the science that its own time
is able to offer. Is there then any form of spiritual culture that can
provide form and expression for it?

And besides --- even if our experiences fully enter into the
great world of truth on which science builds, there will still be an
urge for a more immediate expression of them, an expression in
which all sides of our being, hope and fear, longing and melancholy,
joy and sorrow, come to speak and are allowed to stir over against
what we have experienced.

By both paths we are led over to art. It is an essential side
of its task that is indicated by them.

A very living and intimate relation to art cannot be said to
belong to our time. Original production is not great and not
profound. Great artistic powers are not noticeably in motion. And
for the strenuously working and pleasure-seeking people of our
time art is essentially diversion and relaxation, a means, not
something valuable in itself that attracts and develops by the very
demands it makes.

In a certain sense art is indeed to be a means. It is to satisfy
the urge for images that show us life in its fullness, without the
demands on thought and will that real life makes being directed
to us. It is to present to us the purely individual moods, characters,
situations and fates that science cannot let come into their full
right. It marks a great advance when life can be enjoyed in an
image; it is a sign that the blind urge of life is not autocratic,
--- that in any case there are pauses in its working. Whether rules
or exceptions are expressed in the images presented, art itself
cannot decide. But it is precisely the sharp opposition between
rule and exception from which we are for a while freed in art.
There is nothing absolutely
% --- p. 9 ---
\setcounter{footnote}{0}%
\opage{9}ugly. Everything has its characteristic side, by which it fills a
place in the great connection. Here nothing is left out. Even the
most fleeting is held fast, even the saddest is softened, even the
extravagant is harmonized through a characteristic form.

Even more than for the beholder, art is a liberation for the
artist himself. It presupposes, in the first place, a free use of
mental powers, a release of energy, without the practical demands
of life immediately occasioning it. And it makes it possible for the
artist himself, even more than for the beholder, to rise above all
that can oppress and ensnare. Through the series of great artists'
works their personal development and struggle in life can be read.
Modern literary psychology has the merit of having brought forward
this point of view, on which, however, it perhaps lays too great
weight.

But like science, art makes its demands, and without that it
would not be a link in spiritual culture. These demands are made
on personality itself, both the beholder's and the artist's. Work
must be demanded, in order that it may be the real content of life
that is converted into the image or rediscovered in the image. It
is (as Poul Møller stated in his review of ``Extremerne'') a sign of
the dissolution of an artistic period when poetry is cultivated and
known at second hand, so that a mere spectator's life is lived,
without real experiences forming the ore that is hammered out
into the works of art. Later Kierkegaard, in what he called ``the
aesthetic stage'', described a way of taking life that makes
everything play, both life and art. And Henrik Ibsen in the great
poem ``Paa Vidderne'' --- the counterpart of ``Brand'' --- has
described how one can withdraw from life and see it as if in
bird's-eye perspective, as a picture within whose frame even the
fate of those nearest becomes merely an effective element, --- and
how a spiritual petrifaction, an abolition of all feeling and all
striving, becomes the consequence. As in ``Brand'' the unbridled
willing that merely follows the law of formal consistency leads up
into the barren mountain desert, so here it is the will-lessness
produced by losing oneself in mere contemplation that leads to
the same place.

Like science, art too is a domain that sets
% --- p. 10 ---
\setcounter{footnote}{0}%
\opage{10}us great ends, in whose service we can grow. It is to open our
eyes to what we do not like to see, --- the bitter, painful and
disharmonious in life, --- as well as to what we cannot always catch
sight of ourselves, --- the great and lofty, which we perhaps did not
believe was possible. In great art man does not run up onto the
heights to escape being part of life, but has his sense and capacity
freed, purified, widened, so that a larger view and a firmer grip on
life become possible.

But in the end everything rests on the art of life. Only when life
is strong and great can art be so. The more vigorously and earnestly
we take life, the greater possibility of poetry life will contain, a
poetry that does not come from outside but is a resonance released
by the great movement of life itself, and which we feel, if not in
the midst of the work where it is born, then at least in the times
when the after-swells of the moods of work and struggle fill the
mind.

We must ourselves become artists if we would know what art
is, just as we must ourselves think if we would know what science
is. Life-poetry is the spring from which the great artists draw, and
from which we can all draw when our mind is opened to it. It is
no doubt only genius that can shape images and symbols that can
become the property of all, but we can all --- as surely as our
personality and our life has its distinctive nuance, its particular
fate --- shape images and symbols in which our experience finds
expression, and if we are wholly bereft of this capacity, we do not
understand art. Aesthetic culture demands self-activity, like all true
culture.

But now the poets complain that they cannot express
everything, especially not the most glorious, --- that they cannot
draw life itself into their poems! There are, laments Schack
Staffeldt, two worlds that cannot be united, and between which one
must choose.
\begin{quote}
None binds the image to the thing,\\
None the ideal to its imprint:\\
Between vision and possession life hovers;\\
Choose, for only the one is given thee.\\
(Den unge Digters Klage).
\end{quote}
% --- p. 11 ---
\setcounter{footnote}{0}%
\opage{11}Guyau describes the poet's suffering (Le mal du poëte) at
reality's being richer than the poem:

\begin{quote}
Je me sens trop petit pour l’immense univers . . .\\
Je sens trop pour le dire, et pourtant le silence\\
M’oppresse comme un poids.
\end{quote}

Nietzsche places life above art, as he places art above science:
the overflowing fullness and power of life can find complete
expression neither in art nor in science!

But in such utterances lies precisely the experience that the
deepest poetry is one with life itself and is therefore not felt as an
element distinct from the conduct and course of life itself. If all
the bound force were released, life would stop at once. We need
not choose between vision and possession, when possession, life
itself, gives birth to the highest visions.

Though our circle be small and the reach of our will limited, the
conditions of human life nonetheless manifest themselves in the
small as in the great. And often the greatest images of life, of its
course and its worth, have sprung from human fates that were
outwardly inconspicuous, to the astonishment of those who did not
believe that anything great could come out of Galilee. In each of
us life stirs --- when there really is life in us --- in a way in which
it stirs nowhere else.

The master of all masters has this in view when he says
(A. Winters Tale):
\begin{quote}
Above the art\\
That, as you say, helps nature on,\\
There stands an art that nature has itself created!
\end{quote}
It is this nature's own art that stirs in, or is spontaneously
practiced by, everyone who lives life wholly and fully and in
accordance with his inmost personality. --- Can there be any
spiritual form of life that surpasses this in value and in
significance? Should it not be the case that our deepest personal
experiences and practical experiences, as well as our highest
thoughts, we can express only in poetic form? --- Before we answer
% --- p. 12 ---
\setcounter{footnote}{0}%
\opage{12}this question, we must examine the position of religion in
relation to spiritual culture.

\begin{center}III\end{center}


It is in the religious problem that the consequences of the
division of labor for spiritual life appear quite especially. The
independence that science and art have each attained in modern
times they have won at the expense of religion. For there was a
time when religion was the epitome of men's entire spiritual
culture. Either it was so that it was the one thing needful, against
which all else that, spiritually speaking, could lay claim to men's
hearts stood as worthless or even as a hindrance to be pushed
aside. Or else it was so that religion itself within its frame satisfied
all spiritual needs, --- at once satisfied the urge for knowledge and
the artistic urge. So it is no longer. Religion has got science and
art outside itself. From this comes precisely the dividedness in
modern life. If religion is not now to have life outside itself as
well, it must ever anew proceed from life, from the personal
experience of life in individuals. Just as little as one can have real
science or art at second hand, just as little, and still less, can one
have religion at second hand. Each must be his own priest, just as
his knowledge and his poetry must be self-produced. Even the most
venerable traditions are not enough. They stem, moreover, from
times when quite other conditions of spiritual life prevailed than in
our day. Religion must spring anew from the inmost of each single
personality, from the experiences and movements of life that he
alone, and no other in quite the same way, has gone through. It is
uprightness and honesty that matter first and foremost, and they
first become possible by going back to the inmost source, in
personality, of all faith and all view of life.


I take religion here in the widest sense, as covering all personal
faith or view of life. I speak of something that concerns every
human being's personal relation to life, not merely of
% --- p. 13 ---
\setcounter{footnote}{0}%
\opage{13}religion in the sense of a faith shaped in definite traditions,
dogmas and cultic customs.

But what, then, has philosophy to do with it?

Philosophy cannot produce faith, either in the wider or in the
narrower sense. Faith is itself one of the objects of philosophy.
The philosopher asks, in the first place, whether any faith or
religion can solve the riddles of thought. He asks, in the second
place, what mental powers and what mental urge come to light in
religion. And he asks, in the third place, what significance religion
can have for an ethical conduct of life.

Let us consider these questions each by itself.

\stjerner

As soon as the philosophy and science of modern times were
founded, their spokesmen declared, in the face of the attacks from
the side of religion, that religion was not science and could not
teach us science, and that from the religious side one ought
therefore to let science go its own way and by its own paths.
Utterances in this direction came from Giordano Bruno, whom they
nevertheless did not save from death at the stake, from Galilei,
whom they nevertheless did not rescue from condemnation and
enforced recantation, and from Spinoza, whom they nevertheless
did not rescue from persecution. What one would not hear then is
now repeated again and again by the defenders of religion, when
these do not belong to the Catholic Church. Only the Catholic
Church sees clearly what is dangerous, for all dogmatically shaped
religion, in such a concession.

It is not, as is often thought, the results of science that are
decisive here, however important they may be, as the names
Copernicus and Darwin will already remind us. But the decisive
thing is the manner of asking and the demands that are made on
the answer. The answers religion gives do not fit at all the
questions science asks. Properly speaking religion gives only a
single answer, no matter what questions are asked. It solves all
riddles by a reference to divine omnipotence. In his ``Dialoger''
Galilei puts
% --- p. 14 ---
\setcounter{footnote}{0}%
\opage{14}such an answer into the mouth of the one of the interlocutors
(Simplicio) who defends the old conception of the world. In the
face of the difficulty in the old world-picture, that the whole world
should revolve around the little globe of the earth, Simplicio says
that it is after all easy for God to let the great move around the
small as center. But to this his opponent remarks that in relation
to omnipotence all difference between great and small falls away,
and that it is precisely for this reason that the concept of God
cannot be used in science. If one appeals to a miracle, all
knowledge ceases; all differences and relations vanish before the
infinite.

When we go beyond the domain within which our thoughts and
our expectations can be confirmed or contradicted by strict
observations, we build on analogies that can have only figurative,
symbolic significance\footnote[1]{Cf. the essay on the concept of analogy in \textit{Mindre Arbejder. Anden Række.}}.

If the value of religion rested on its theoretical or intellectual
significance, it would soon vanish from the world of the spirit. For
those who ask clearly it gives no answer to a single question.
Therefore, too, (especially in the Protestant world) the discussion
of the religious problem has passed over into psychological and
historical research.

\stjerner

What keeps religion up, and constantly brings about new
adherence to its traditional forms as well as new attempts at the
making of religion, are motives that belong not in the domain of
knowledge but in that of feeling and will. And the deepest motive
consists in an urge to have assurance of the constant persistence
and unfolding of what makes life worth living. What is the highest
for a human being will, he hopes in his religion, also be the highest
in the world, or it will at least have its firm ground in existence
and will triumph, in us and outside us, however it may go with
% --- p. 15 ---
\setcounter{footnote}{0}%
\opage{15}all else. Such a hope, held fast through life and with clear
understanding of the course and fate of the world and of life, as
the individual knows it or at least thinks he knows it, --- is the
religious primal phenomenon, which can take countless different
forms, and which must take different forms in the different
personalities.

When one calls one religion or faith higher or lower than another,
the only possible standard is to be found partly in what it is that
man attaches so great value to that he hopes for its persistence,
but partly also in the weight and range of the experiences he has
at his disposal, and in the energy with which the hope is held fast.
--- These are precisely the points of view that the historians of
religion, without always being conscious of it, apply when they
undertake a valuation of the various religions in relation to one
another.

We saw, when we spoke of science, that there always remained
individual and special experiences and practical experiences that
were bound up with the individual's inner and outer circumstances
of life, and that could not be exhausted by scientific points of view.
To these belong now, first and foremost, the experiences that lie at
the root of a man's faith. But it is of the greatest importance that
man should have his experiences at first hand and not receive
from tradition or dogma direction as to what he shall try to
experience, which easily leads to one's forcing oneself into a
spiritual form of life that does not suit the personality, and thereby
shutting oneself out from bearing the witness about life that it is
the task of genuine religion to bear. Echo is to be avoided, in the
domain of faith as well as in that of art and of science.

\stjerner

Where religion (faith, personal view of life) has come into being
as the fruit of a life in the service of great ideals and in the work
of holding them fast through the course of fate, there we meet the
image of a spiritual power and concentration, of an inwardness
% --- p. 16 ---
\setcounter{footnote}{0}%
\opage{16}and a self-containedness that gives great religious personalities a
place among the greatest in the world of the spirit.

And here lies in my eyes the decisive problem. Can this inward
and concentrated kind of personal life be preserved on the same
conditions and in the same forms as in past times, and if it cannot,
must not new forms arise for such a life, if spiritual life is not to
lose in power and in depth?

In my view we stand only at the beginning of an unbiased,
profound and purely human treatment of these questions.
Dogmatics and criticism have each in its own way hindered such a
treatment. The powers have been taken up partly with shaping life
after traditional models, partly with working their way out of
traditions. Now life itself must first have its say and testify to its
real urge, to what one needs in order to hold the highest values
fast through the vicissitudes of fate. It is hardly the same that is
needful for all.

On a railway journey in Jutland I once heard a peasant explain
to another the difference between theology and philosophy thus:
``Theology is superstition, philosophy is unbelief.'' --- I answer this
only on behalf of philosophy.

Philosophy is for me nothing negative, but a free inquiry into
the nature and conditions of spiritual life. And it points clearly for
me to what the conditions are for the continued persistence of a
deep and vigorous spiritual life. But since the philosopher is not a
prophet, he cannot guarantee that these conditions will actually
come about. Since such a spiritual life itself belongs to what is of
highest value, it is the philosopher's personal hope that it will ever
anew unfold itself. Its greatest enemies are stubborn clinging to
forms that belong to a vanished spiritual culture, and blasé
indifference to what lies beyond the circle of the narrowest
interests\footnote[1]{What Tegnér said almost a hundred years ago still holds: \begin{center}Det murkna gamla, det omogna nya\\ med blind förbittring kämpa nu om verlden.\end{center}}. But if it proves possible to preserve the inwardness and
self-containedness of
% --- p. 17 ---
\setcounter{footnote}{0}%
\opage{17}spiritual life, it will find its expression, not in dogmas or
systems, but in the symbols of a great and genuine life-poetry. And
then the unfortunate effects of the division of labor in the spiritual
domain will be overcome. The divided powers will meet when they
each in its own domain go back to the source from which in the end
they all spring and are nourished.

% --- p. 18 ---
\setcounter{footnote}{0}%
\opage{18}\essay[A Philosophical Confession (1904)]{A PHILOSOPHICAL CONFESSION.\footnote[1]{Combined from two lectures, at Harvard University and at Wellesley College, in the autumn of 1904.}}
\streg

In the program of the philosophical department of Harvard University,
1904---1905, it is announced that one of the philosophical professors
[William James] will develop a theory of pluralism on the basis of
experience, and that in the next term his colleague [Josiah
Royce] will develop a speculative theory of the absolute; and it is
remarked that ``the student will have the opportunity of comparing
two different views regarding fundamental problems: radical
empiricism and absolutism.''

It is not very often that philosophical differences are announced
so openly and in so true a spirit of free discussion. As I had myself
discussed the fundamental problems in a little book (``Filosofiske
Problemer'', university program 1902, translated into German 1903,
later also into English), and as I, during my visit to
Harvard in October, 1904, was very much interested in the discussion
which was about to take place in American philosophy, Professor
William James proposed that I should give a lecture
for his students, in which I should describe my position on
this question. The present paper gives the essential content of
my lecture, but I have also found a place here for some parts of a
paper which I read for the Philosophical Club at Wellesley College.

% --- p. 19 ---
\setcounter{footnote}{0}%
\opage{19}\begin{center}I\end{center}


The importance of pluralism, i.e., of the tendency to accentuate
the multiplicity and the difference of phenomena, depends on its
power to raise problems. Both thought and sensations suppose
difference, contrast, variation. In the psychology of our time
Fechner's investigations especially, on the relation between physical
impression and psychical sensation, point in this direction. And our
thought starts with greatest energy when two judgments, each of
which seems well founded, contradict one another.


I believe there is special reason for accentuating this point in
the actual state of philosophy. There is much metaphysics in the air,
and we are on the way to forgetting what we have learnt from
positivism and criticism. The old English school had the mission
to keep the attention of philosophers on experience, and it
started the great movement against dogmatism. It is no
accident that the great setter of problems, David Hume, belonged
to this school. In evolutionism this school has said its last word---the
widening of the concept of experience to connote not only the
experience of the single individual, but the organized experience
of the whole species. We may hope that a new, refreshing start
will be made.

Pluralism makes the world new for us and necessitates a revision
of our categories, our principles and our methods. A dogmatic
sleep is too tempting for the human mind. We are inclined to
suppose that we can develop---or perhaps already have developed---thoughts
in which all existence can be expressed. But, as a Danish
thinker, Søren Kierkegaard, has said, we live forward, but we think
backward. Understanding comes after experience. Only when life
is closed can it be thoroughly understood. This is our tragico-comical
situation. Even a divine thinker could only understand the world
when the life of the world was finished.

\begin{center}II\end{center}


But pluralism as such brings no understanding. To
understand is to connect one fact with other facts, to
find a uniting principle. Multiplicity as such would only
% --- p. 20 ---
\setcounter{footnote}{0}%
\opage{20}make description and classification possible, and even this only under the
condition that the manifold phenomena were not only different,
but also similar. The only meaning
of ``understanding'' which pluralism can acknowledge is understanding
as mere recognition, not as derivation or explanation.

Very often we must be contented with stating a
single, isolated fact. But then this fact at once raises a problem---even
by its isolation. We have an interesting example in the botany of our
time. The Dutch botanist, Hugo de Vries, maintains that new
types can arise suddenly. Great variations---not small, as Darwin
thought---are, according to de Vries, the condition of evolution.
But, if this is so, it ought to be\footnote[1]{See \emph{Hugo de Vries}: \textit{Die Mutationstheorie.} I p. 366; II p. 662.} the task of natural science to
explain how this arising of a new type was possible,
to find the hidden conditions for that which was for our observation
a sudden appearance, to discover a continuity
between it and other appearances.

Now, it is a fact that we in many cases have found such connection
or continuity in nature. It is the ideal of knowledge to
find it in all domains of observation. Our mind can only understand by
synthesis, and the principle of continuity is therefore the presupposition,
the working hypothesis, of all science. Therefore we must also acknowledge
continuity as a characteristic of reality. We have no
right to suppose that the fact that we can not understand phenomena,
if we can find no connection or continuity, should be without
ground in reality itself. If we will build our philosophy
on experience, we ought to give full importance to connection, unity and continuity,
as well as to difference and multiplicity.
Experience shows us both, and pluralism can, therefore, not
be the last word of the philosopher. And there is an inner connection
between continuity and multiplicity. All qualities, powers
and characters which we ascribe to the single elements or
beings which pluralism acknowledges are only known through the connection of these elements
or beings with a whole order of things.
We can, for example, only ascribe energy to a being because we experience
% --- p. 21 ---
\setcounter{footnote}{0}%
\opage{21}that it actually does a certain work, that alterations in it or out of it
have their cause in it. If it were absolutely isolated,
we could not ascribe any predicate to it, we could not
know it at all.

Perhaps it is impossible to develop a metaphysical theory which
shall give both facts their full right. But this ought not to
lead us to forget the problem.

\begin{center}III\end{center}


I, for my part, call myself a monist, because connection and
continuity seem to me to be more important facts
than multiplicity; it is, as I have shown, only through
their connection one with another, that things can be understood.
Every difference and multiplicity supposes a deeper continuity,
which it is our task to discover. But as an absolute
monism can not be carried out, I call myself a critical monist,
the word ``critical'' being taken here in the same significance as when Kant
called his philosophy critical philosophy. A perfect and universal
synthesis is for us always an ideal which has to struggle for
existence. Every totality we find in nature has always a history;
it has developed through interaction of elements and
supposes differences in the nature and the tendencies of the different elements.
As a critical monist I say: If we can not carry out
our monistic ideal, the reason could be that reality
(``world'', existence) is not finished, is yet in
full development. Only if we could think this development completed
in any time, would the world be quite intelligible. There
is at least one thing which is not completed: our
thinking; and this is also an element in reality!
It is the wonderful contradiction of the great rationalistic systems---Plato's, Spinoza's
and Hegel's---that they can not explain
the striving and struggling thought whose work these
systems themselves are.

The incompleteness of thought, the necessity of always renewed
striving after truth, has, perhaps, its cause in this, that new
elements may arise in the world. It would, then,
not only be we who discover something new, but it would be the
% --- p. 22 ---
\setcounter{footnote}{0}%
\opage{22}world itself, which was renewed and had to bring its new elements
into harmony with the older. And there would then be an
inner connection between our struggling thought and the essence of
reality. This is my metaphysics, if I have one. ---

The problem also of mind and matter is perhaps conditioned by
the incompleteness of our knowledge. There is no necessity
for a choice between spiritualism and materialism, if
mind and matter are not supposed to stand in a contradictory (two-membered)
opposition one to the other. This contradictory opposition is often quite dogmatically assumed.
But existence may, as Spinoza taught, have more forms,
qualities or attributes than our experience can show us. Critical
monism points to the possibility that we have not all conditions
given for solving the problem.

\begin{center}IV\end{center}


That an absolute systematization of our knowledge is not
possible is no evil. The history of philosophy shows us that
systems have more energy, depth and freedom in their first
stage, in statu nascenti. According to that treatment of this science
which of late years is more and more general, the history of philosophy is more and more the history
of the starting, of the impulses of thought
and of the leading experiences of great thinkers, not only the description of the completed systems.
And even in its completed form a system is only a means to
understand ourselves and the world. It is an essay which will try the
strength and importance of certain ideas or certain experiences. The great systems
are projections, electric search-lights, with whose help we try
to explore the dark.

But the thought with whose help we find our way through
the world is itself only a part of the world, and we do
not know how important a part of the world it is. Can the
whole content of the world be translated into thought? Can we,
as metaphysical idealism would do, conceive the world in its
innermost ground as an expression of thought? To me
it is evident that all such idealism is founded on an
analogy which can not be verified. We can only use analogy
scientifically when we apply it to throw light on one domain of experience
with the help of another domain of experience. Here a
% --- p. 23 ---
\setcounter{footnote}{0}%
\opage{23}verification is possible through the consequences to which the analogy leads us; we can see if these consequences agree
with our progressing experience. But when we
maintain an analogy between a part of existence and existence
considered as a whole, no verification is possible. Analogy
transcends here very soon the limits of science, and has
only poetical value, if it has any value at all.
It may be sublime poetry, if the metaphysician is a man of genius---poetry of a more intellectual, though
less emotional character than the poetry of
the great popular religions, which also are founded on analogy.

I correct myself: we ought not to say ``only poetry.''
Our deepest life-experiences can only be expressed in poetical form.
There is a life-poetry which is a very serious thing,
the most concentrated expression of the experiences which
can be made of the value of life and of life's endeavors.

\begin{center}V\end{center}


Our thought is both larger and smaller than reality.
It contains forms and possibilities which never are fully realized.
It develops hypotheses between which we have to choose,
as only one of them can be true. In order to reach the aim of knowledge
we must try several means, survey the possible
ways. The progress of knowledge consists, from this point of view,
in the reduction of these possibilities, perhaps in the exclusion of
all except one. Existence or reality is in every single
case one single definite thing, is one in opposition to our many
possibilities. As Schiller says, Eng ist die Welt und das Gehirn
ist weit. --- But from another point of view the world is richer
than our thought. It contains the possibility of experiences which
have not yet been made. New discoveries and experiences lead again
and again the thought on new ways. Thought is, as already
said, only a part of the world, and there are always more things in
heaven and earth than are dreamt of in our philosophy. Here,
then, we can say the reverse of Schiller, Weit ist die
Welt, und das Gehirn ist eng.

This incommensurability, or this (in a mathematical sense)
irrational relation, between thought and reality will always
% --- p. 24 ---
\setcounter{footnote}{0}%
\opage{24}be experienced anew. Our intellectual labor will always have to progress
on the one side in simplicity and concentration, and on the other side
in complexity and extension.

\begin{center}VI\end{center}


My sympathy is more on the side of positivism and empiricism
than on the side of speculation; and in methodological---in a certain
degree also in metaphysical---respects I have a sympathy with pluralism.


I will now say a few words on the way in which my thoughts
have developed in this direction.

In my youth the influence of Søren Kierkegaard was decisive
for me. He waged a passionate war against speculation,
with strong accentuation of the conditions of thought and of the value of the real,
personal life. Later on the study of Comte and of
the old English school gave me a new start in the direction of a
philosophy of experience. And my own more independent studies have
maintained me on this way. In this last respect three
points have been of special importance for me.

My psychological studies led me to accentuate the differences
between individuals in respect of the relations between psychical
elements. Even if these elements could be said to be
``the same'' in all individuals (which only to a certain degree is
true), the manner in which they are combined, the ``timbre,''
will be different. The general psychological laws manifest themselves
in many different forms in the individual cases. There is
here a multiplicity which no analysis can exhaust. In my
Philosophy of Religion I have more specially pointed out the great importance
of individual and historical differences, and I have described the
most important types of religious life. We still miss a comparative
human psychology. The study of individual differences
is only in its beginning.

Also my ethical studies have led me in this direction. The
problems of scientific ethics are of two sorts. The
first task is to find and establish the ethical principles, the
standard we are to use in valuing human actions
and institutions. Already, here in the starting-point,
% --- p. 25 ---
\setcounter{footnote}{0}%
\opage{25}great personal and historical differences make a
scientific foundation of ethics appear as a great problem.
And even if this first difficulty is conquered, a second problem
arises: how are the general principles to be
applied in the special cases? Human individuals have
different powers and dispositions and start their development under
very different social and historical conditions. They can,
therefore, neither qualitatively nor quantitatively have the same tasks
and the same duties. How can we be sure that
the right sort of work is chosen? And how can we
be sure that sufficient work is done, that ``the quantum
satis of human will'' has been applied? Also here new experiences
are still possible, and with them new problems. ---

The third point belongs to the theory of knowledge.
I saw---what others had seen before me---that Kant has not solved
Hume's problem, and that this problem is still standing. It is a necessity for
our thought to apply the law of causality; but this does not justify the assertion that this law is universal
and metaphysically necessary. We may construct the concept of an ideal, complete and
perfect experience, and for this experience the law of causality would have universal validity.
But the experience which we really have is limited and imperfect.
New observations are still made, which are to be connected
with our earlier observations. This contrast between ideal
and real experience---and with it Hume's central problem---was
undervalued by Kant's speculative followers. At the end
of the nineteenth century the problem was taken up, not only
by philosophers, but also by men of science. It is now more and
more admitted that the importance of scientific and philosophical principles
consists in this, that they lead us in our striving after understanding.
Their truth is their validity, and their validity is experienced
in their power of leading us in our intellectual work. A principle
is true if it can be applied, if we can
work with it, i.e., gain understanding with the help of it. Truth is
a dynamic concept; it manifests itself in the working of our thought.
And it is a symbolic concept, because it only presupposes
an analogy, not an identity between thoughts and events.
This holds of the truth of our sense-qualities; they have objective
value as symbols, but can not be proved to be images
% --- p. 26 ---
\setcounter{footnote}{0}%
\opage{26}of things. It holds, too, of the truth of our formal-logical principles,
of the principle of causality, etc. We can not compare
our sensations and our principles with an absolute order of things.
Surely we have no right to regard it as a pure accident that just these
special sensations and thoughts make it possible to gain a progressing
knowledge; but neither have we the right
to regard them as direct revelations.

\begin{center}VII\end{center}


Where we can not apply the principles of our thought we end
with a problem. And, so far as pluralism can and must
be maintained, so far as the events and elements
remain isolated and unconnected, so far
we have unsolved problems before us. But pluralism would
lead us to false supposition if its meaning were that the value of a phenomenon
is diminished when it can be understood, i.e.,
brought in rational connection with other phenomena. If
we specially consider personal beings we must admit that
every personality is a little world with its inner order,
its law of development and of interaction of psychical elements.
No character without an inner continuity! And this little world
can not maintain its value without standing in interactions
with other personal beings as elements in a social totality.
Our ethical endeavor is to produce greater continuity
in the inner world of personality and in the greater
world of society, as our intellectual endeavor is to
find a still greater continuity in the world at large.

And here lies for me the great importance of irrationalism
or pluralism, so far as we must acknowledge
it. The world is not complete, not harmonious, not rational; therefore
there is a work to be done. A little American girl who
had been told that God created the world once for all,
asked her mother: ``But in what business is he then now?'' (``But in what business is he then now?''). This was a quite
philosophical question. There is perhaps in the world at large a
% --- p. 27 ---
\setcounter{footnote}{0}%
\opage{27}great work going on, through which it is developed to greater
rationality and harmony. But for us it is of the greatest importance,
that there is a work to be done by us, that our
own work in thought and action is a reality, a real
factor in a great process of evolution. There is here an analogy between
the problem of knowledge and the ethical problem, which agrees
with their standing as the two chief problems of philosophy.

% --- p. 28 ---
\setcounter{footnote}{0}%
\opage{28}\essay[The Concept of Will (1905; 1912)]{THE CONCEPT OF WILL \\ (1905).}
\streg

T\textsc{here} is a peculiar dispute about the will in the psychology of our time. On the one side it is denied that the will can be maintained as a distinctive standpoint for conscious life, a distinctive side of it, or an element in it, --- however one may wish to put it. On the other side it is asserted that the will designates precisely the most fundamental standpoint from which conscious life can be viewed; conscious life is willing from first to last, so that sensations, ideas and feelings too become intelligible only in their connection with the will. The first school goes back historically to Hume and is represented in recent psychology by Münsterberg, Ehrenfels, Ebbinghaus and Lapie; the second is proclaimed by Fichte and Schopenhauer and is represented in recent psychology by Wundt, William James, Fouillée and Jodl\footnote[1]{Cf. also \emph{Ernst Mach}: \textit{Erkenntnis und Irrtum.} Leipzig 1905. p. 21 f. (Later note).}. I side with this latter school, and in this essay I shall try to give the grounds for my position.


The problem is closely connected with the position of psychology as an empirical science. It can be put thus: can a willing be the object of a single and immediate observation? And if this question must be answered in the negative, is it then justified to
% --- p. 29 ---
\setcounter{footnote}{0}%
\opage{29}set up the will as an independent standpoint for conscious life? --- My position is that I answer the first question with No, but nevertheless answer the second question with Yes. I thus agree with the first school in its premises, but deny that it draws the right conclusion from them, because it overlooks certain fundamental facts.

When we keep to what immediate and single self-observation presents, we always find only sensations, feelings and ideas. What we call phenomena of will appears with peculiar sensations of unrest, of yearning, of a straining of force that we seem to put in, of energy that we seem to expend. They are sensations that are more restless, less in equilibrium, than sensations of color and tone, but not different in kind from these in principle. They can be taken as sensations of movement. Sensations of muscular tension. As soon as we merely think of an action we have a mind to perform or are on the point of performing, various muscles are tensed and a corresponding peculiar sensation arises. In the moment of decision this sensation is especially strong. Alongside it there appear here ideas of aims and means, of ideals and norms, of obstacles and effects; but these ideas develop quite like other ideas and do not seem able to lay claim to a special chapter in psychology. Further, we observe feelings of pleasure and pain, of hope and fear, of enthusiasm and disgust, and so forth; but all such things belong to the psychology of feeling. --- Is this analysis exhaustive, or could something have been overlooked?

In my view something very essential has been overlooked. It is correct that, if by analysis one understands the turning of attention upon single definite points, each of which can be considered in a certain isolation, we always find in psychological analysis only sensations, feelings and ideas. But we have then let abstraction begin too early. A fundamental fact has been overlooked which, to be sure, cannot be the object of a \textit{single} observation. Our sensations, feelings and ideas never occur isolated in reality. In actual experience they are like drops
% --- p. 30 ---
\setcounter{footnote}{0}%
\opage{30}in a stream, and it is the \textit{direction and speed} of the stream that determine the mutual relations of the drops, and thereby also determine, more or less decisively, the nature of each single drop. What is really given is not a chaos of isolated elements, but groups and wholes which are always, more or less markedly, directed toward a certain aim, toward a definite point, whether this is itself an object of consciousness or not. In no state of consciousness whatever is this orientation lacking, this being set in a certain direction. But a direction cannot of course be apprehended by a \textit{single} observation, only by holding several observations together. By counting the drops in a river and making a catalog of them we should not discover that the river has a definite direction. If psychological analysis is to bring out all the facts, it must not forget this most essential of all psychological facts. It was Hume's one-sidedness that he could not see the wood for the trees, and this one-sidedness flourishes among his modern successors. And yet, when we look closely, the growth and form of the single tree are understood only through its place in the wood. The analytic method must be supplemented by a genetic one, which is not content with what reflection can separate out at stages where growth is far advanced, and where the apparent independence of the single elements can therefore more easily deceive. If by a psychical element one understands everything that observation leads us to distinguish as something that cannot be derived from anything else in consciousness, then the constant direction of conscious life is just as much an element as the single sensations, feelings and ideas in their isolation. There is a constant interaction between the direction of the stream and the nature of the drops, but the direction can no more be derived from the drops than they from it. The direction is the historical element in mental life. It determines from first to last what becomes end and what becomes means, --- what is striven for and what is shunned, --- what becomes aid and what becomes obstacle. Our various sensations, feelings and ideas stand at very different distances from the central determinateness of direction, or, as we may call it, the
% --- p. 31 ---
\setcounter{footnote}{0}%
\opage{31}original and constant urge and striving; but their mutual relations, their inner nature and often their very existence are determined by their relation to it.

Able psychologists naturally cannot overlook this. But when they acknowledge it, and yet deny the independence of the element of will, they fall into decided contradiction with themselves. Thus \textit{Ebbinghaus.} As ``special elementary forms of mental life'' he assumes only sensations, feelings and ideas. Phenomena such as willing or striving are only ``peculiar combinations of sensations, feelings and ideas'' and offer ``nothing new'' in relation to these. (Grundzüge der Psychologie I, p. 168). But although he denies that acts of will are ``fundamental phenomena of mental life in the sense in which sensations or ideas are so'', he nevertheless admits that they are ``the fundamental forms of the unities in which alone sensations, ideas and feelings originally occur in reality''; although ``conceptually'' they are not something ultimate and original, they nevertheless are so with respect to time and origin (zeitlich und genetisch). ib. p. 561; 565.) --- To this I remark: 1) How can it be said that acts of will offer \textit{nothing new,} when they nevertheless consist in \textit{peculiar} combinations or unities? Whence does this peculiarity derive? 2) What difference can one make between \textit{fundamental forms and elementary forms,} or between fundamental forms and \textit{fundamental phenomena?} Must not the fundamental form in any case belong to the fundamental phenomenon, so that the latter can neither exist nor be understood without the former? And should not a fundamental form that expresses itself in everything from the simplest to the highest mental phenomena rightly be called an elementary form? 3) If sensations, feelings and ideas originally have \textit{reality} only in the combinations or constellations that we call acts of will, then those three kinds of ``elements'' are mere abstractions when we consider them outside the connection in which alone they are real! Atomistic scholasticism is put in the place of empirical psychology. An atomistic theory can have its methodical justification; modern natural science is a great example of it. But it must first be shown that
% --- p. 32 ---
\setcounter{footnote}{0}%
\opage{32}any result whatever is attained by introducing this mode of consideration into psychology, and even if it is, one must always distinguish between a mere working hypothesis and an exhaustive explanation of what is really given.

According to the conception of the will that has been put forward in the foregoing, and that even Ebbinghaus, though an opponent of ``voluntarism'', acknowledges, the will is intimately connected with the whole nature of conscious life, provided it is correct that this nature shows itself in a synthesizing activity, a synthesis, which can be traced in all sensing, all remembering, all thinking and all feeling, as well as in all willing. If the will is the most marked expression of synthesis, and if on the other hand synthesis is an activity, a function, then it is even justified to say that consciousness in its innermost essence is a willing. It is this standpoint that I have taken as the basis of my Psychology. I have tried to show that sensing, thinking and feeling are not understood without this original and constant element of activity, so that the whole of psychology --- even if in its form of presentation one follows the usual threefold division --- is in reality a psychology of will. At every point within the psychology of knowledge and of feeling we come back to a fundamental activity and a fundamental urge as the ultimate presupposition. The concept of will implies in particular that this synthesizing activity always (from its very first beginning, which is one with the beginning of individual conscious life) goes in a definite direction, inasmuch as the kind and extent of the synthesis are conditioned by a single element around which the concentration takes place. The concept of will expresses conscious life more from the real side; the concept of synthesis characterizes it more from the formal side.

An analogy, which is perhaps more than an analogy, we have in the fact that, even from a purely physical point of view, the directions of motion in nature are just as original as the forces and atoms we presuppose. At no point in the development of nature have we the right to presuppose an absolutely resting chaos. Certain directions must always be presupposed in which the forces act and the atoms move, and these directions cannot be derived from the atoms alone. In every given or presupposed
% --- p. 33 ---
\setcounter{footnote}{0}%
\opage{33}state the atoms move in definite directions and with definite velocity; this is the historical element in nature, which natural science since the time of Kant and Laplace has more and more acknowledged. And the very properties that it ascribes to the single atom we infer from the whole connection in which the atom stands, and from the way in which it is supposed to intervene in this connection.

However one conceives the relation between consciousness and the organism, it is clear that it is through the nervous system that consciousness and will can exercise an activity in the material world. The significance of nervous life lies partly in its centralizing function: through it the organic functions can cooperate in a more perfect way than would otherwise be possible; partly it lies in its capacity to store up energy, which can then later be released under definite conditions and in definite directions. Already in a purely organic way, in growth and in the first involuntary movements, there is expressed here too an original direction of movement, which can be altered by the external conditions, but is not from the first produced by them.

To me it now seems the most natural thing to assume that what for self-observation is a process of consciousness is the same as what for physical and physiological observation is a material, organic process. We must take good care not to ascribe absolute validity in nature itself to the distinctions we must make in order to see clearly, although of course we cannot take it for an accident that we must make certain definite distinctions in order to understand nature rightly. And if what we psychologically call will should, in its innermost ground, be connected with the energy that acts in the organism, the originality of the element of will would be evident. But even apart from all speculation on this question, this originality will, I maintain, be clear already on purely psychological grounds.

% --- p. 34 ---
\setcounter{footnote}{0}%
\opage{34}Psychological terminology is unfortunately still so variable that it can sometimes look as though psychological disputes were only a dispute about words. One might, e.g., say that I take the word will in a wider sense than various other psychologists, and that the settlement of the whole question rests on a definition. It is clear that the fewer marks a concept has, the greater the totality of phenomena for which it can hold. Content and extension stand in an inverse relation to each other. If, in order to assume a willing, we demand no more than what can be shown even in the most elementary phenomena of consciousness, it is a matter of course that the concept of will acquires a very great extension. About words I will of course not dispute. But the point is that the narrower, and still more the narrowest, sense of the word will entails arbitrary and unnatural distinctions in a domain where exact observation shows a continuity that reaches ever deeper. It shall now be my task to show this, beginning with will in the narrowest sense of the word and passing step by step to the more elementary forms.

Some psychologists --- in the most extreme form \textit{Paul Lapie} in his ``Logique de la volonté'' (Paris 1902) --- will speak of will only when an action is preceded by two consciously formed judgments, one of which declares an aim worth attaining, and the other that the means of attaining it can be procured. I will something when I am clear that I \textit{should} do it and that I can do it, or, as Lapie expresses it: Une action est volontaire, quand elle est préjugée bonne et possible. The nature of the will depends entirely on the processes that precede it, and Lapie draws from this the consequence that the will is really only a special form of the intellect. Even psychologists who do not go as far in the intellectualist direction as Lapie would restrict the concept of will to such actions as are performed with clear consciousness of aims and means. Thus A. F. \textit{Shand} in his excellent essay ``Types of Will'' (Mind 1897) and \textit{Pierre Janet} in his work ``Les obsessions et la psychasthénie'' (Paris 1903), epoch-making for psychology and psychiatry.

% --- p. 35 ---
\setcounter{footnote}{0}%
\opage{35}Now there are not many actions that are performed with full clarity about aims and means. The horizon of the will is always limited, and clarity in any case has many degrees. Both the range of the experiences that come to influence which aims one sets and which means one finds, and the degree of strength with which these experiences assert themselves, vary infinitely. But what I lay special weight on is that it does not always make an essential difference whether what asserts itself in us and determines our aims and our means appears in the form of conscious judgments or not. We are inclined afterwards to express in the form of a conscious judgment what in the moment itself appeared as an estimate or as an immediate perception (an intuition). That an explicit judgment is formed means only that a given content becomes the object of clear and distinct consciousness; it is the conversion of a mental content into a new mental form; --- but no change in principle of the content necessarily follows from this. And besides there are --- as I have shown in my essay ``Det psykologiske Grundlag for logiske Domme''\footnote[1]{Det kgl. danske Videnskabernes Selskabs Skrifter. 6. Række. Hist. og filos. Afdeling IV, 6. (1899).} --- intermediate forms between immediate perceiving (sensing, remembering, fantasy) and judging; especially to be stressed is what I have called articulating perceiving, where the bringing out of a single element leads to a shift or transformation of the perceptual image without the setting in of the dissolution and analysis that is the condition of judgment. Such perceiving can have greater value and practical significance than an explicit judgment, in which what is gained in clarity is often lost in energy. We really pass an explicit judgment only when we can no longer find vent in mere perceiving, or when there is a doubt to overcome, or when an urge to communicate with others stirs.

It is in any case not necessary that two judgments be passed. When I feel the value of something, a striving to acquire or produce it will arise, and when the means lie ready to hand, or when the recognition of the
% --- p. 36 ---
\setcounter{footnote}{0}%
\opage{36}valuable involuntarily (as by a kind of instinct) releases the movement in its direction, one judgment, the judgment of value, will be enough; the other (the one we may call the judgment of possibility) will not be necessary.

The next step will be to maintain that not even the judgment of value is necessary. It is not necessary that the recognition of the value or the setting up of the end take place with full consciousness. The valuable can stand before me in a great image or a thought that becomes the star I steer by, and that immediately, without any analysis and any judgment, gains power over my striving. In such cases the Danish word Drift applies. (The German ``Trieb'' is ambiguous, whereas the French désir and the English desire, as well as Spinoza's cupiditas, correspond to the phenomenon in question).

We see how many nuances there are, and how unnatural it is to point to a single point on this whole scale (two judgments --- one judgment and an involuntary striving --- the end given as an image) as the one where a wholly new concept should come into force.

But we are not yet at the end. The concept of impulse implies that an end appears before us, even if it is not formulated by an explicit judgment of value. Logically speaking, the concept of value does indeed underlie the concept of end, since we make something an end only because it has value; but psychologically speaking this positing of ends precedes the judgment of value, just as practice precedes theory. That we ascribe value to something we discover from the first only by having made it an end. There is, moreover, a still more elementary form of striving for something that proves to be valuable than impulse. There need be no consciousness of the end at all, but an obscure urge leads in a definite direction, inasmuch as it does not allow a halt before an aim not known beforehand has been reached. Through a series of involuntary releases of energy the individual is carried forward to a more or less valuable result. Such a union of urge and capacity without consciousness of an end we have in instinct. (Where the urge occurs without the capacity, it perhaps leads to action, but it then becomes more a matter of chance
% --- p. 37 ---
\setcounter{footnote}{0}%
\opage{37}whether anything valuable is attained, that is, something that, when consciousness awakens, can be set as an end). Such an urge (Spinoza's appetitus, Fouillée's tendance) would be the most elementary form of what we, in the widest sense of the word, could call will. I do not see that in this whole series --- from the urge to the decision determined by the judgments of value and of possibility --- there is reason to set up absolute boundary posts.

\whiteline

Of the greatest psychological interest, but at the same time most enigmatic, is the transition from urge to impulse, or from the involuntary to the voluntary (if by voluntary action one understands an action that is conditioned by the idea of an aim). Here again there are many degrees of clarity as regards the idea of the aim, right down to the threshold. The more dim and obscure the idea of the end becomes, the more we pass over from impulse to mere urge. It can often be difficult to decide whether an idea of an end has been present at all, and great problems thereby arise for self-examination and for moral self-scrutiny. Read, e.g., in Søren Kierkegaard's journals (for 1849) his repeated deliberation, swinging back and forth, on the question whether his manifold activity as an author was preceded by the idea of a definite ethical aim that he meant to reach through it. But apart from this difficulty a psychological problem appears: how it comes about that something becomes an end for us, that is, how the transition from the involuntary to the voluntary takes place. This transition must of course take place involuntarily: the first end is not derived from any other end. Only later can it happen that we make something an end because it is part of an end set up earlier, or a means to its attainment. But the first end --- what Dante has called ``the first thought of the will'', --- cannot be derived in this way. It is at this point that we approach the original direction or tendency in our mental life pointed out in the foregoing, the ultimate foundation of the will. If our being, before
% --- p. 38 ---
\setcounter{footnote}{0}%
\opage{38}the first idea of an end emerges, were entirely indifferent or neutral, it would be entirely impossible to understand how we could set ourselves ends. It is otherwise if mental life, even before it has the capacity to form ideas, has an original direction, a yearning or striving, that can be promoted or inhibited by the experiences it undergoes. The condition for anything at all being able to acquire value for us is that we will something even before we know what we will. Thus we must conceive the state of animals during instinctive actions (insofar as we can imagine it at all). The state need not be wholly unconscious. The urge can be felt even if what it aims at is not felt. Goethe goes too far when he says: ``Der gute Mensch in seinem dunklen Drange ist sich des rechten Weges wohl bewusst.'' It is really a self-contradiction. What Goethe meant to express is the security with which one can abandon oneself to an obscure urge, without knowing whither it leads in the involuntary selection of means and ways. Helen Keller (the blind and deaf-mute American student), on the other hand, has aptly described her state before it dawned on her that other people had a means of communication other than the finger language: ``Before I knew that a deaf child could learn to speak, I felt dissatisfaction with the means of communication I already possessed. One who is confined to the finger language alone always feels a constraint, a limitation. This feeling began to make me uneasy; there came a painful, goading sensation of a gap that had to be filled. My thoughts often rose and would work their way up like birds against the wind, and I constantly used my lips and my voice'' (Story of my Life, p. 58 f.). Two elements in the urge are brought out here in an interesting way: an indefinite feeling of a lack, and an involuntary, instinctive release of movements. Had these conditions not been present, the news that the method of speech could be applied would not have struck her like lightning, as she later relates. In an instant the end emerged into consciousness. When Schiller rightly says that man grows with his greater ends, it must be added that man must have grown
% --- p. 39 ---
\setcounter{footnote}{0}%
\opage{39}in order to be able to set himself greater ends, indeed, in order to be able to set himself ends at all. It is not only at the first stages of the volitional life that we are compelled to go back to the obscure world of the involuntary, if not of the unconscious. It may also become necessary in the midst of the development of the clear volitional life. When deliberation leads to a point where a choice must be made between two ends, each of which has developed and taken root in the soul, the issue can often be settled only by one of them having a stronger hold in the obscure urge that lies behind all values and all ends. If the struggle between the conscious ends has been fought out without a settlement, a new power must step forth and throw its weight into the scale. Perhaps man thereby makes progress in self-knowledge, as he now discovers elements in his being that he had not known before, or that only now had grown so strong that they could have a say. (Cf. on this my Psychology VII B, 1; 4---5). New formations can occur in the domain of the will, just as in the domain of feeling and of fantasy.

Philosophers have laid different weight on the transition from urge to impulse, from the involuntary to the voluntary. Some, such as \textit{Spinoza and Schopenhauer,} lay no weight at all on this transition. ``Whether man is conscious of his urge (appetitus) or not'', says Spinoza (Eth. III, def., expl.), ``the urge remains one and the same''. And according to Schopenhauer ``the will to live'' is unalterable and the same in men and animals, although man can become conscious of it and the animal cannot. This conflicts with experience, since it is often a significant turning point when the direction of the urge is recognized, and when aims and means can be represented. Extension or limitation, alteration or consolidation thereby becomes possible; perhaps even a whole metamorphosis of character is initiated. It is inconsistent of Spinoza and Schopenhauer to deny the significance of that transition, since they both, each in his own way, lay so much weight on the metamorphosis, the spiritual deliverance, that can take place precisely through the influence of the conscious life of thought and feeling. In sharp contrast to them J. G. \textit{Fichte} has ascribed an absolute
% --- p. 40 ---
\setcounter{footnote}{0}%
\opage{40}significance to the transition from the involuntary to the voluntary. It is a miracle, a transcendence, when the first end is set. ``The urge (Trieb) itself is not my product, but the product of nature\dots{} But the urge comes to consciousness, and what it is to effect in this domain lies in my power\dots{} Here lies the transition of the rational being to self-activity; here lies the definite, sharp boundary between necessity and freedom\dots{} The form as urge \textit{with consciousness} is due to the free act of reflection; but that an urge exists \textit{at all} is due to nature'' (Sittenlehre 1798, p. 159 ff.) With what right the transition from urge to impulse is called ``a free act of reflection'' is not easy to see; before this transition there is, after all, no reflection present at all\footnote[1]{Fichte states (Sittenlehre p. 35 f.): ``Der schlechthin angefangene Zustand wird nicht schlechthin an Nichts angeknüpft . . . Nur wird er nicht an ein anderes Sein angeknüpft, sondern an ein Denken.'' --- But surely there exists no thinking before the first thought has arisen! And if there existed a thinking, it would surely have to concern something definite, and so be an actual ``being''!}. The transition takes place involuntarily, as Aristotle and Dante already saw. The enigma does not thereby disappear, but one does not solve the riddle with the catchword of a free act, and one has no right, at this point of the development of the will, to posit a gulf between two domains as opposed as ``necessity'' and ``freedom'' are supposed to be according to Fichte. There is no reason to assume a break in continuity here. The enigma that remains concerns not only the psychology of the will, but also, and in quite as high a degree, the psychology of ideas.

\stjerner

It would be wrong to divide the phenomena of will into two mutually exclusive kinds, involuntary and voluntary. For not only does the transition from the involuntary to the voluntary take place involuntarily, but in the so-called voluntary actions a greater or lesser fullness of involuntary moments always co-operates. It is not least here that the intellectualist theory of the will shows itself in all its insufficiency. With
% --- p. 41 ---
\setcounter{footnote}{0}%
\opage{41}full clarity we never act, and what conditions the outcomes of our will can never be exhausted in definitely formulated judgments. From this derive the many problems of the volitional life, which will constantly give both the thinker and the poet new work. Even in our most conscious actions, then, the urge to release stored-up energy co-operates. Without a surplus of energy no action is possible, and the direction in which the release takes place is determined partly by sense impressions, as in reflex movements and instinctive actions, partly by ideas, as in impulses, partly by judgments, as in deliberate decisions. And all these forms of will act together in our conscious actions, with all possible degrees of clarity and strength, all possible relations between the central and the peripheral in us. The involuntary can exist without the voluntary, but not conversely. The development from the elementary, involuntary willing to the willing founded on deep self-understanding often takes place amid great conflicts and fluctuations, and it takes place in different forms in different individuals. There are many dangers on this road, many possibilities of morbid and one-sided development, but on the other hand there also shows itself here the great wealth of variety of human mental life, which is still so little acknowledged --- or known --- by our ordinary morality. Only a comparative psychology of will, when we once get one, will be able to do full justice to the actual volitional life. To one point only shall I refer here: the inverse relation between strength, certainty, concentration on the one side, and fullness of content and a comprehensive horizon of will on the other. Every time the horizon is widened, the possibility of a crisis arises at this point. Here the volitional life stands in the closest interplay with the other sides of mental life.

\stjerner

The will stands in intimate connection with the feeling of pleasure and pain. The relation is not only that pleasure and pain determine the direction of the will, become its ``motives'', but also the converse, that we feel pleasure and pain at what
% --- p. 42 ---
\setcounter{footnote}{0}%
\opage{42}respectively promotes or inhibits the fundamental direction of our life. And this latter relation is the most primitive. It gives us the only possibility of understanding in some measure the biological significance of the feelings. Just as sensations of sense have their biological significance as signals of what is going on in the surrounding world or in our organism, so feelings of pleasure and pain are symptoms of whether our striving is advancing or receding, in the simplest case a striving to live, to persist, the one Hobbes called conatus primus. Through our feelings we thus come to know our innermost willing. Tell me what you rejoice over or grieve over, and I will tell you what you will! An evolutionary mode of consideration leads us to divine, if not fully to understand, the significance of the feelings. In the most elementary beings neither symptoms nor signals are available. For them it is a matter of being or not being, of living or dying, without their feeling beforehand in which of the two directions their life is moving. The rise of the feeling of pleasure and pain marks a higher type of life, since the symptoms make it possible to alter the direction, or to hold to it and quicken the pace, according as the demands of life require. As an English investigator has put it: substitution of pleasure and pain for life and death as the sanctions of conduct! \textit{(Hobhouse:} Mind in Evolution. London 1901. p. 388). This consideration can, by way of analogy, be applied to all levels of life, the highest as well as the lowest, the most ideal as well as the most material. Everywhere the presupposition of mental life is a fundamental activity in a definite direction, and with it the possibility of furtherance or hindrance. What can become food for the mind is determined, just as much as what can give us bodily nourishment, by the urge and direction of our nature, which is not the same at every level, but whose history is the real history of our life. Spinoza has expressed this most clearly of all when he says: ``We do not strive for, will, desire or wish for anything because we deem it good, but we deem it good because we strive for, will, desire and wish for it'' (Eth. III, 9. Schol). But it is an old truth, which Aristotle and Augustine already knew.

% --- p. 43 ---
\setcounter{footnote}{0}%
\opage{43}To what degree the original direction of our life (our temperament, our talent, our character) can be shaped and altered when a consciousness of aims and means has been awakened in us, that is the great question that leads beyond psychology to ethics and social hygiene. But every theory must admit that we begin at a level where we have not placed ourselves, and with qualitative and quantitative presuppositions that we have not given ourselves. Our first willing is to such a degree one with ourselves that it cannot be regarded as our own work in the sense in which a later willing can. And this first willing can be traced, through many shifts and metamorphoses, right up into the highest forms of will. It will give these a distinctive stamp, determine their timbre. Fichte maintained a psychological impossibility when he admitted indeed that ``the point at which we find ourselves when we are first able to use freedom does not depend on us'', but added that ``the series which we describe from this point depends, for all eternity, in its whole extent, entirely on us''. (Wissenschaftslehre. 2 Aufl. p. 266.) In so external a way one cannot keep the involuntary and the voluntary, the given and the self-produced, apart from each other. We come into being as active beings, and even in our highest activity the conditions of our coming into being co-operate in determining it.

\stjerner

On the conception I have put forward here, it is a long series of different phenomena that come to stand as the object of the psychology of the will. But the question then naturally arises whether a general concept can be set up for such different phenomena that can hold for them all.

Every concept that is to hold for a series of different phenomena rests on analogy. From example to example the relations vary, the more exactly we examine them, to such a degree that it will not be possible to ``extract'' marks that recur everywhere in quite the same way. But although all single features can vary, a similarity can be present in the relation between them, and it is such a
% --- p. 44 ---
\setcounter{footnote}{0}%
\opage{44}similarity of relation that makes the formation of concepts possible\footnote[1]{Cf. my essay \textit{On analogy and its philosophical importance} (Mind. April 1905) p. 202). (In Danish in Mindre Arbejder II).}. In the phenomena I have brought forward in the foregoing, degrees of consciousness and forms of consciousness vary as much as is at all possible. Even psychologists who narrow the concept of will in advance more than I consider possible and necessary despair of forming a typical concept of will\footnote[2]{Thus \textit{Shand} says (\textit{Types of Will.} Mind. 1897, p. 325): The more closely the typical forms of will are studied, the more we shall appreciate the difficulty of embracing them in any one supreme type . . . The profoundest introspection will not show us the universal character of will. --- It is of interest to me to see that for Shand the difficulty consists not only in the diversity of the phenomena of will, but also in the way in which unconscious and involuntary processes co-operate, even at the higher levels of the will: Those deep forces within us which work for the most part unseen, their tendencies unforeseen, whose object only rise into clear thought at times, . . . seem to us the real and abiding will, and their conscious expression an accident or momentary phase, the more play of thought upon their upmost surfaces., ib. p. 306.)}. I believe, however, that it is possible to find a point of similarity among them which at the same time makes it possible to characterize the element of will in its difference from the element of knowledge and the element of feeling. From the involuntary urge to activity and urge to self-preservation up to the willing that chooses with clear consciousness, there are two features that constantly recur under different conditions, namely 1) that the direction of activity is determined by preferring\footnote[3]{Already \textit{Charles Bonnet} has this definition: Un être qui préfère un état à un autre état, et qui agit conséquemment à cette préférence, est un être qui a une volonté. (\textit{Essai analytique sur l’âme.} Copenhague 1760, p. 114.) Bonnet adds (p. 115) that a being that throughout its whole life had the same sensation, and in the same degree, would have no capacity to will, would have no will at all. --- In most recent times the same definition is found in \textit{Fouillée}: La force inhérente à tous les états de conscience a sa dernière raison dans l’indissolubilité de ces deux fonctions fondamentales: le discernement, d’où naît l’intelligence, et la préférence, d’oú naît la volonté. (\textit{Psychologie des idées-forces.} I. p. IX). Les discernements en apparence indifférents sont un résultat ultérieur (p. X).}, and 2) that what is preferred is essentially determined by the individual's \textit{own inner nature.}
% --- p. 45 ---
\setcounter{footnote}{0}%
\opage{45}On this rests the analogy among all phenomena of will. The equilibrium must be abolished, and a difference of direction, determined by the relation to the individual's nature, must assert itself. Since now, as long as life lasts, there is never perfect equilibrium in us or around us, an uninterrupted willing stirs in us, and it is only where the differences become of a more intensive kind that our preferring takes place with consciousness\footnote[1]{An involuntary preferring takes place in every instinctive action, as when an animal rushes at the definite object and leaves others alone. Likewise in involuntary attention, as when the eye, on a many-colored surface, dwells on certain definite places but passes lightly over others. Here distinguishing and preferring are one, as Fouillée says. (1912).}. Urge to activity, reflex and instinct all presuppose a difference in inner and outer conditions, and consist in an overflow, an expansion in one direction or the other. In impulse the difference is conditioned not only by impression and sensation, but by an idea of an aim and by the opposition between this aim and the present state, an opposition that can be felt quite immediately. In decision there occurs a conflict between two values, each of which may be formulated in judgments, and by this conflict the outcome is determined, inasmuch as the value that has its strongest hold in our nature is preferred, and precisely thereby comes to stand as the greatest value.

As preferring, as practical positing of differences or active differentiation, the will comes under the law of relativity, the chief law that holds for all sides of conscious life, and that is evidently connected with the fact that we live in a world full of differences and oppositions. An existence raised above all opposition and resistance could neither will, feel, think nor sense. The chief law of conscious life points back in turn to the fundamental form of consciousness, synthesis, the synthesizing activity: there can be distinguishing and preferring only when two elements can assert themselves at once and be held together as closely as possible; only then does the difference, and perhaps in the same moment the preference, spring forth. And, as already
% --- p. 46 ---
\setcounter{footnote}{0}%
\opage{46}remarked, since this fundamental form is an activity, the will is more akin to it than any of the other sides of conscious life. We are the more active, the more the causes of our states and our actions lie in ourselves, in our own nature, the original nature and the nature acquired through experience and action. The standard for the development of the volitional life is thereby given. ---

\whiteline

Addendum (1912). If by a psychological element one understands a side or a quality of mental life that can \textit{immediately} be observed as something single, in a certain difference from other such sides or qualities, then there are no psychological elements other than sensations, ideas and feelings of pleasure or pain. But it is unnecessary and unjustified to stop at such a definition. And it is more natural to let the word element in psychology designate a side or a quality of mental life that asserts itself in all mental states, whether it can be the object of immediate observation or not. And something of this kind is, for every coherent consideration of mental life, precisely what the word will (or better: willing) in the widest sense is an apt designation for.

It was established long ago that will is not an object of immediate observation in the way the other sides of conscious life are. \textsc{Berkeley} is the one who first saw this. Only through its results is our activity known, and we have no simple idea of what is designated by the word will, he teaches (Principles of Knowledge § 27). \textsc{Hume} extends this consideration to hold for every causal relation, within mental life as within the material world; nowhere can causality be apprehended by immediate observation. And \textsc{Fichte}, that pronounced philosopher of will, says expressly that one was quite right in denying that freedom could be an object for consciousness. ``Ich habe Kausalität, heisst bekanntermassen: das, was ich mir als Zweck vorsetzte, trifft in der Erfahrung ein'' (Sittenlehre p. 272 cf. p. 175). ``Das Ich'', he says further (Wissenschaftslehre. Neue Ausg. 1802. p. 413), ``wird seines Handelen unmittelbar sich nie
% --- p. 47 ---
\setcounter{footnote}{0}%
\opage{47}bewusst, sondern nur in dem Produkte, und vermittelst des Produkts''.

This clear demonstration was not adhered to by the psychologists. \textsc{Schopenhauer} made willing one with urge and yearning, with emotion, and thereby obliterated the difference between feeling and will; this made it possible for him to assert that willing could be observed immediately. And \textsc{Maine de Biran}, the author who next to Schopenhauer has had the greatest influence on the psychology of will of the nineteenth century, found willing immediately given in the sensation of effort or straining, a sensation of which he then gave a spiritualistic interpretation\footnote[1]{See on Schopenhauer and Biran ``\textit{Den nyere Filosofis Historie}'' II. p. 221---223 and p. 300---302.}. \textsc{Søren Kierkegaard}, too, wanted to demonstrate the will, the ``leap'', as he called it on his indeterminist presuppositions, as a distinctive object of psychological observation. But he is too clear a thinker not to see that the act of will itself is not a state, but expresses itself through the relation between two states, and so in ``Begrebet Angst'' he embarks on the impossible task of showing how that which is itself not a state and does not fill a moment, but for every sharp analysis always lies between two states or two moments, should be able to become an object of observation. As if, besides what we call the cause and what we call the effect, one should also be able to observe the causal relation itself!\footnote[2]{Cf. my book ``Kierkegaard som Filosof''. p. 70---81.}

A closer analysis of the actual states of mental life proves Berkeley and Fichte right. And they are right also in this, that the circumstance that willing cannot be observed immediately by no means excludes that will can be a constant and original element in mental life.

When we examine more closely under what conditions our sensations, ideas and feelings arise, and what nature they acquire in the different states owing to the difference of the conditions, it turns out that it is a great and
% --- p. 48 ---
\setcounter{footnote}{0}%
\opage{48}misleading abstraction to speak of sensations, ideas and feelings as if they were the same, quite alike, under all circumstances. Not only is it of decisive importance with which other elements they appear together, but it is also of quite particular importance what direction the whole course of our sensations, ideas and feelings has in a given state. A sensation of sense acquires a different character according as it stands as occasion and introduction to a reflex movement or an instinctive movement or not. If it corresponds to an influence that releases or tends to release such a movement, it stands otherwise in consciousness than when it corresponds to an influence that has no essential motor effect; in the latter case it becomes more resting, in the former it acquires, figuratively speaking, a direction forward, an unsteady character, which is removed only when sensations of the released movement and of its effects appear. This is even clearer with ideas. An idea can have a resting character, inasmuch as its content holds the interest in and for itself. But this is not the original state, and hardly the usual one. The first ideas, the ``first thoughts'', are ideas of ends, ideas of something that has proved to bring satisfaction and that therefore becomes the object of impulse, or at least of expectation. Expectation is probably secondary in relation to impulse; from the first no ideas arise except as motives, as occasions and introductions to movement and action. Expectation is already a more resting, contemplative state. If impulse meets obstacles, it either falls away or leads to the idea of means. But an idea acquires a different character when its content is a means than when its content is an end. And here there is an order that cannot be altered in the single case. One cannot simply let the idea of the means and the idea of the aim change places within the given state. The idea of the aim conditions the idea of the means. Later, during the action, the idea of the means precedes the idea of the aim. In both situations, each taken by itself, the two ideas stand in a definite
% --- p. 49 ---
\setcounter{footnote}{0}%
\opage{49}relation to each other. The relation of direction comes out clearly here. ``One and the same'' idea thus acquires a different character according as it appears as a member of a motive or not, and where it is a motive (a member of a motive), it acquires a different character according as it concerns means or end. The relation between motive on the one side, and impulse, wish, intention, purpose or decision on the other, cannot be better designated than as a relation of direction. And the word direction is here at bottom no more figurative than all psychological expressions whatever. --- As for feeling, finally, it is originally, and at bottom always (since absolutely resting feelings hardly occur), either a symptom that an unconscious movement is favored or inhibited, or a member of a motive that forms the occasion and introduction to action determined by an end. And feeling acquires a different character according as it is a symptom or is a motive or is of a more resting nature, and further according as it is attached to the idea of an end or to the idea of a means.

The difference between motive and aim (or satisfaction) appears, as Sibbern already remarks, only in voluntary (i.e. idea-determined) action. Involuntary action runs its course so continuously that only by way of analogy can we apply that difference in describing it. It is precisely in involuntary action that the greatest difficulties for description and analysis lie. There is no transition here from one clearly appearing member to another; rather there occurs a series of pulsations, a progressing, more and more concentrated process.

Characteristic of the states in which the stamp of direction is strongest is a peculiar condensation or concentration, for which the word attention is a good expression, although it stresses the intellectual element too much. Taught by experience, we can judge from the concentrated character of a state that it points in a certain direction, will give rise to consequences of importance, even if we perhaps cannot at once make out in which direction it points. By analogy with this we must also
% --- p. 50 ---
\setcounter{footnote}{0}%
\opage{50}conceive the state of an animal when an instinctive movement is about to be released. We then speak of an urge asserting itself. The concentration can itself be motivated, as when we exert ourselves to call up a memory or to solve a problem. It is in this sense that one can say that the will can act on itself, inasmuch as preceding motives (concentrations) determine later motives (concentrations).

The more elementary mental life is, the more it holds that every sensation, idea and feeling is determined by direction, that is to say, has the character of occasion and introduction. Animal psychologists lay quite particular weight on this. It is in any case difficult to demonstrate free ideas in animals, and with them also ideas of ends proper. But in connection with this it is also maintained, with full right, on the side of animal psychology that it is an entirely mistaken conception of mental life to think of it as built up of independent elements of sensation, feeling and idea, as one still all too often finds it, even in authors who believe themselves raised above the one-sidednesses of the old ``association psychology''. In the development of mental life, \textsc{Thorndike} maintains (Animal Intelligence. New York 1911. p. 154; 240), there is not a progress from the small and simple to the large and complicated, but from direct to indirect connection, from undifferentiated states to distinguishing or generalizing, and we must therefore turn our usual human psychology upside down and make what was mentioned only in a special chapter on the will into a general psychological law.
--- It is this conception that I put forward long ago in the fourth chapter of my Psychology, and that is entirely decisive for all talk of the psychological reality of the will. Pure sensations, ideas and feelings are abstractions, nothing else.

Comparative human psychology leads to the same result. The mental life of more primitive men does not permit the clear difference between different mental elements that that of the adult European does. \textsc{Lévy-Bruhl} has recently illuminated this relation in an interesting way in his work ``Les fonctions mentales
% --- p. 51 ---
\setcounter{footnote}{0}%
\opage{51}dans les sociétés inférieures'' (Paris 1910). He even goes so far as to ascribe to primitive peoples a logic quite different from ours, inasmuch as their ideas are from the first so infiltrated into one another, and so interwoven with emotional and motor elements, that they do not admit of clear comparison. The ideas are from the first members of connections (participations), from which they can only slowly be separated out so as to acquire a truly intellectual character (p. 442). And in his masterly characterization of the involuntarily formed ``community'' of earlier times, in contrast to the ``society'' of later times, shaped by conscious agreement, \textsc{Tønnies} (Gemeinschaft und Gesellschaft. 2. Ausg. 1912, p. 109 ff) describes what he calls ``essential will'', an involuntary life, a striving in which there is an immediate transition from state to action, in contrast to the voluntary action guided by free ideas. The former is characteristic of the primitive ``community'', the latter of modern ``society''. --- Sociology thus warns us, just as much as animal psychology does, against operating with abstractions taken from a more differentiated stage of the history of mental life as if they were eternal and original principles that could be found again at all stages. It must even be maintained that differentiation is never so far advanced that those abstractions express the full truth. They always hang in the air.

Neither experimental, descriptive nor analytic psychology can do without the genetic standpoint. Mental life has, in the single individual and in the whole race, a history without which it is not understood.

From pure elements of sensation, feeling and idea one can never in the world derive the living, definite, practical character that the actual states of mental life always present in a lower or higher degree. But conversely one can, starting from these states, hand in hand with experience, show how mental elements can step by step lose their original, practical, motivating and motor character and approach a relative state of rest, which is the one that many psychologists find easiest to observe, and which language, which
% --- p. 52 ---
\setcounter{footnote}{0}%
\opage{52}always has difficulty in expressing life, the dynamic, misleads us into regarding as what is really given.

I take the occasion to remark that the conception I have put forward here is already presented in the first edition of my Psychology (1882) (see sections IV, 7 and VII B, 4). As early as when the first German edition (1887, corresponding to the second Danish edition) had appeared, Christoph Sigwart expressed in a letter to me his interest in my regarding the will so decidedly as the fundamental element of consciousness. In later editions and in other writings I have developed this conception further and made it more precise. But from the very first my presentation of psychology was laid out in such a way that I first tried, within the psychology of knowledge and of feeling, to show that neither of them could be understood without presupposing elements of will, and then, in the special section on the psychology of the will, collected and ordered these constantly presupposed elements of will and thereby found the ultimate foundation for the understanding of mental life.

Recently \textsc{Schmied-Kowargik} (Analytische Psychologie. Leipzig 1912. p. 192) has stated that when I assume that the direction of mental life can be felt, I thereby at bottom admit that the will can be observed. But direction can everywhere be apprehended only through comparison of different observations, not by a single observation. Direction is always apprehended by a kind of construction, which is carried out more or less consciously. It is only the \textit{immediate} apprehension of willing that I deny. A direction in mental life can assert itself already in and through the very growth of mental life, in which unconscious and conscious elements work together. Thus when instincts awaken, e.g. at the transition between childhood and adulthood, or when, already earlier, the instinct of self-preservation takes on special forms, or when habits develop. In all such cases an involuntary willing expresses itself; later, when free ideas have arisen, which from the first also happens through involuntary growth, a voluntary willing, determined by ideas and ends, can arise and take on the various forms that descriptive psychology distinguishes as impulse, wish, intention, purpose and decision. But at all these stages it is the definite succession of states and
% --- p. 53 ---
\setcounter{footnote}{0}%
\opage{53}elements, and the impossibility of interchanging the members of the series, that furnish the mark of the reality of the will. There is here, as I have shown in my Psychology, an analogy between the conception of the will and the conception of reality, which is of great interest. No more than the will does ``reality'' correspond to a single element or a mere sum of elements. The will is the inner reality within which in the end all the elements of mental life find their definite place, determined directly or indirectly by the energy and direction of the whole mental striving.

% --- p. 54 ---
\setcounter{footnote}{0}%
\opage{54}\essay[Humane Ethics (1901)]{HUMANE ETHICS.\footnote[1]{An outline for use in lectures at the Folkeuniversitet (1901).}}
\streg

M\textsc{en} at all times pass judgments on good and evil
(ethical judgments), which are called forth by the feelings and ideas
awakened in them by the contemplation of human actions, traits of
character, endeavors and orders of life (institutions). \textit{Ethics}
inquires whether such judgments can be justified, and whether there is
a standard by whose help a distinction can be drawn between valid and
invalid ethical judgments. When the justification and the standard are
sought in human nature and in the conditions of human life, we have an
attempt at \textit{humane} ethics, in contrast to such attempts as start
from authorities (religious, social or political) as the ultimate source.

There is an inner connection between human actions, traits of
character, endeavors and orders of life, so that they cannot be
discussed exhaustively each by itself. Here, however, we shall keep
chiefly to the traits of character. We begin by setting up the
hypothesis that one single definite trait of character is the most
valuable, and by giving a \textit{description} of it (I). Next we try to
show how this trait of character can develop in man (the
\textit{psychological possibility)} (II), and on what its \textit{ethical}
value is grounded (III). The attempt to give such a justification leads
naturally to a discussion of \textit{the conditions and limits to which
every scientific ethics is subject} (IV): We then compare the ethical
point of view with other
% --- p. 55 ---
\setcounter{footnote}{0}%
\opage{55}points of view on human actions and traits of character, and
compare at the same time the different main types of ethics which
history presents \textit{(comparative ethics)} (V). Finally we cast a
glance at \textit{the most important oppositions which in our time assert
themselves in the ethical domain} (VI).

\begin{center}I\end{center}


Among the traits of character which have at different times been
regarded as valuable (as virtues), \textit{justice} has as a rule held a
prominent place. We shall here try to show that from an ethical point of
view it is the most valuable of all qualities, and we begin, then, by
describing it.

In civil life he is just who gives everyone his own, keeps his
promises and does not encroach upon the property and freedom of others
(as property and freedom are ordered at a given time and in a given
place). Where the individual will not act in this way, he is compelled
to do so by the power of the state or by public opinion. --- A certain
likeness to this kind of justice is borne by the consideration for
others and the self-limitation and self-control which decorum and
courtesy demand.

Characteristic of these forms of justice \textit{(elementary justice)}
is, 1) that \textit{limitation} (the negative) is more prominent than the
free unfolding of personality, 2) that \textit{outward} action or omission
plays a greater role than disposition and the inner will, and 3) that
\textit{regularity,} uniformity in the manner of acting, is demanded
without regard to the individual peculiarities and circumstances of the
different personalities.

Elementary justice proves insufficient for two reasons in particular.
--- \textit{The inner and outer circumstances of the individual} may be so
peculiar that the same manner of acting cannot be demanded of him as of
others. One then applies ``mercy'' or ``equity''; but if these are not to
be unjust, they must be the expression of a higher justice. --- Where it
is not merely a matter of maintaining an existing order of life and the
distribution of rights and duties that follows from it, but where
\textit{a new order of life} is to be founded, a new distribution made,
one must
% --- p. 56 ---
\setcounter{footnote}{0}%
\opage{56}go back to the free unfolding of personality and to the inner
disposition. The new springs from the inner life of personality.

Through such experiences there develops the thought of \textit{an ideal
justice,} which demands not merely a limitation of the individual's
self-assertion out of regard for the self-assertion of others, but
demands that the individual's development shall positively further that
of the others, and conversely. This comes about through self-assertion
and self-surrender (sympathy, altruism) entering into a harmonious
relation to one another. Just as the individual's self-assertion
benefits the others, so his self-surrender benefits himself by
developing faculties and impulses which would otherwise lie fallow.
While elementary justice on the whole regards the individual as a means
for the persistence of society, according to ideal justice no personal
being comes to stand as a mere means, not even where limits are set to
it and burdens laid upon it. The external relation between means and end
here falls away. Elementary justice abolishes only the outward, literal
struggle between men, and secures against brutality and encroachment by
compelling men to observe certain limits and forms, however remote from
one another they may stand in disposition. Ideal justice lets
individuals feel themselves as organs of a living whole, free and
dependent at once. They draw nourishment from the common life, and
through their peculiar conception of life and conduct of life they in
turn react upon it and determine its further development. Self-surrender
is not extorted, and self-assertion is not seized by force.

\begin{center}II\end{center}


Psychologically considered, justice can develop predominantly either
from the urge to self-assertion or from the urge to self-surrender. We
briefly describe these two types of development.

1) Just as, according to the law of inertia, every material point
remains in its state of rest or of rectilinear motion where no external
cause acts upon it, so every living being has \textit{a striving to keep
itself up, to assert itself, to act and to spread itself in existence.}
This striving is necessary for the persistence of life, but if it ruled
unconditionally and in elementary form
% --- p. 57 ---
\setcounter{footnote}{0}%
\opage{57}in all, it would make life a struggle between egoisms. In
reality it is broken, damped and modified more or less, and on this
rests the possibility of a development in the direction of justice.

a) The difference between itself and others the individual learns to
know by experience, and the very idea of itself it forms under the
influence of others. It learns from others to say ``I'' of itself, and
involuntarily sees itself in the light of their praising and blaming
judgments. It has a strong urge to imitate others and a great
receptivity to suggestions from others. In different natures
\textit{imitation and receptivity} will assert themselves in different
degrees, but always sufficiently for an altogether isolated individual
to be an impossibility. Both by social heredity and by physical heredity
the individual is bound to the species. When the urge to self-assertion
becomes fully conscious of itself, it has already been bent in definite
directions under the influence of the social surroundings. And this
influence goes on acting in the course of life. Just as involuntarily as
music awakens our moods, our intellectual, emotional and volitional life
is influenced by the models, traditions and orders of life under whose
dominion we live. There are those whose whole view of life and conduct
of life are predominantly determined by the involuntary way in which
they have thus lived themselves into the species.

b) But the urge to self-assertion can grow so strong and conscious
that the individual makes his own enjoyment and power, in general his
own isolated welfare (wherever he places it), his chief end, perhaps his
only end. Other beings and other endeavors then become for him only
means or hindrances. Since the others now go on existing and asserting
themselves in their own way, it comes to \textit{a trial of strength and
cunning between the individual and society.}

Even where this struggle ends in a victory for the isolated
individual, unconditional self-assertion will nonetheless undergo
\textit{refractions and modifications.} It has had to go by roundabout
ways, to use cunning and to sacrifice strength. Scars remain after the
struggle. --- But the struggle may also end with the individual learning
to limit its unbridled urge to self-assertion. It finds its account in
% --- p. 58 ---
\setcounter{footnote}{0}%
\opage{58}\textit{recognizing certain limits once for all,} so that it can
then move all the more securely as it pleases within these limits.
Elementary justice has then come into being. --- There is perhaps even a
possibility that the other personalities and endeavors come to stand as
something other and more than means or hindrances, --- that they acquire
immediate value. Partly they often act in the same direction as the
individual himself, so that a union and mutual completion becomes
possible. Partly a kinship in nature and fate is experienced which makes
it easy for the individual to put himself in the place of the others and
to view life from their point of view. Partly the individual's strength
may be so great and overflowing that it is able not only to bear his own
life, but gives rise to an urge to perform a more comprehensive work by
acting also for other personalities and endeavors. By different ways the
individual may thus come \textit{to stand as a member of a whole whose
welfare he more or less consciously makes his own} welfare. Within this
whole, this realm, each single personality stands as a peculiar end, and
each single endeavor unfolds according to its nature and according to
its significance for the whole. Elementary justice has then been widened
into ideal justice.

2) The urge to self-surrender too can at first stand in opposition to
justice and yet be developed into justice. --- This urge can express
itself toward persons (in family feeling, friendship, love of mankind)
and toward endeavors (art, science, charity, politics). In both cases
that to which one surrenders oneself stands as immediately valuable, not
merely as a means for one's own self-assertion.

For several different reasons the urge to self-surrender needs a
development through which it can develop into an urge to justice. --- a)
It expresses itself most strongly in narrow circles, toward a single
being or a single endeavor, perhaps even toward a single moment or a
single side. By its limitation and its blindness it can hamper the
welfare of its very object, or it can injure human life in wider circles
no less than unbridled self-assertion can. \textit{The removal of the
limitation and the blindness} takes place through
% --- p. 59 ---
\setcounter{footnote}{0}%
\opage{59}experience, since, the more the connection and the kinship
between the different personalities and endeavors dawn upon
consciousness, the more it comes to stand as a self-contradiction to set
absolute limits to self-surrender, even though its kind and degree must
be different in the different relations. In this way love of mankind has
historically developed from narrower to wider circles. --- But it is not
merely a matter of the eye being opened through experience to
\textit{what is common.} It must also be sharpened for \textit{the
peculiarity of each single object (personality or endeavor),} so that
the task can become to work at once for the single object according to
the demands of its particular nature and for a harmony among all
personalities and endeavors. --- Justice thus arises here through a
cooperation of clear understanding and deep sympathy. But only he can
have justice in this sense who is himself a strong and marked
personality. The urge to self-surrender would come into contradiction
with itself if it led the individual to make himself nothing before his
object. For then there would be nothing to surrender! \textit{Only a
strong self-assertion makes self-surrender fruitful and comprehensive.}
Self-surrender must, where it has developed into justice, stand as the
expression of an overflowing strength which is able to bear more than
its own burden and to take full account of the peculiarity of other
beings and endeavors. Otherwise justice toward others becomes injustice
toward oneself. ---

Two types of development have here been depicted, between which
there lie many nuances. --- Although the goal of the two processes of
development may be said to be the same, the difference of the ways will
nonetheless influence the character of what is attained. \textit{Justice
will acquire a different stamp, a different timbre in the two types,}
just as a tone sounds different when it is produced on different
instruments.

\begin{center}III\end{center}


Every value presupposes an urge that is satisfied. The urge need not
be felt beforehand: the satisfaction may surprise us (like an
unexpectedly kind word or a beautiful view)
% --- p. 60 ---
\setcounter{footnote}{0}%
\opage{60}so that only from the joy we feel do we infer that there must
have been an urge. But the urge can also stir with consciousness, often
with longing or with pain. And whether it is bound up with consciousness
or not, it can lead to an involuntary striving. --- What satisfies the
urge has \textit{immediate} value. Where this can be attained only
through work, it is made an \textit{end,} and what helps us to attain the
end then acquires \textit{mediate} value. --- For one and the same man
there may be many values, and a coherent, consistent valuation will
therefore be possible only if there is a \textit{fundamental value,}
which conditions all other values, and on which it depends what ends man
sets himself in the particular case.

The recognition of justice as the most valuable trait of character
presupposes as fundamental value \textit{the welfare of society or of the
species.} The welfare of a society is the greater, the more inwardly its
members are bound together, while at the same time each of them develops
his powers and impulses as freely, peculiarly and many-sidedly as
possible. Even elementary justice makes this possible to a certain
degree, though in a mediate and external way. In ideal justice, on the
other hand, the persistence of society and the free self-unfolding of
individuals go hand in hand in the highest degree. Each individual
stands at once as end and as means, both in relation to the other
individuals and in relation to society as a whole.

If we presuppose a fundamental value other than the welfare of society
or of the species, the valuation too becomes another, and it is then not
certain that justice comes to stand as the most valuable trait of
character. For the egoist his own isolated welfare (wherever he places
it) is the fundamental value; if he can attain it, it must be indifferent
to him how things go with the others and with society as a whole. It is
similar with those whose fundamental values are limited by family,
class, sect, nation or race. History shows us a struggle between these
different fundamental values. It is a struggle between as many ethical
systems. The struggle may lead to a change of the fundamental values,
and the thought of justice has arisen during this struggle, especially
at
% --- p. 61 ---
\setcounter{footnote}{0}%
\opage{61}such times when a clearer understanding and a deeper feeling of
humanity became possible. For European humanity, Greek antiquity,
Christianity and the Enlightenment and humanity of the 18th century mark
turning-points in this respect. ---

To point out the ideal, and with it the standard, that corresponds
to a certain fundamental value, has not merely the theoretical interest
which clarity about the relation between premises and consequences
always has, but gives us at the same time a guiding thought (a
principle) by which development in the individual and in society can be
valued and regulated. Involuntary development proceeds blindly,
one-sidedly, scatteredly and slowly. A clear consciousness of ideal and
standard can remedy this. And in particular the involuntary ethical
judgments can be tested with regard to their validity. A conduct of life
or an order of life is the more perfect, the more self-assertion and
self-surrender become one, and the more the single personality in its
relation to society comes to stand both as end and as means. (The ethics
of economic life, of the relation between the sexes and of the penal
system may furnish illuminating examples of this.)

\begin{center}IV\end{center}


Ethics as a science is subject to a twofold limitation. The one has
already been indicated: the dependence of all valuation on the
fundamental value. Besides this limitation, which rests on the
\textit{foundation,} a limitation will always also assert itself with
regard to the \textit{content:} no theory will be able to exhaust the
wealth of life and in strictly scientific form justify tasks and duties
for the wholly special and individual circumstances. Work of thought
must constantly be applied to reduce this twofold limitation to the
least possible; but it is important always to keep it in view.

1) One cannot refute him who rejects the very foundation on which one
builds one's reasoning. And in ethics the foundation is the fundamental
value which determines the ends one sets oneself or recognizes at all,
and with them in turn the laws and norms that are assumed to be valid.
The fundamental value cannot be changed except when the character, the
whole stage of life, is changed. Every living
% --- p. 62 ---
\setcounter{footnote}{0}%
\opage{62}being's pleasure and pain, joy and sorrow, love and hate are
conditioned by the stage of life on which it stands. The caterpillar
does not feel pleasure and pain at the same things as the butterfly. The
savage does not recognize the values of the European, and a Greek of
antiquity would not recognize the values of modern man.

The different fundamental values which succeed one another in the
individual or in the species naturally stand in inner connection with
one another. It is a growth, a historical development, that leads over
from the one stage of life to the other. It is the business of
psychology and of history to describe and explain this development. It
is the \textit{whole} man and the \textit{whole} of human life that
undergoes this development, not merely the life of thought. It is
therefore a process of life, not merely a process of thought, that leads
from a certain stage of life, and with it from a certain fundamental
value, to other stages of life and fundamental values. The justification
of the development can be fully established only when it is completed.
Ethical development in history is therefore a constant struggle, not
merely a struggle between thoughts, but also between feelings and
passions. \textit{Every single attempt at scientific ethics corresponds
to a definite stage in the great struggle for life of the fundamental
values.} Its scientific character rests on the degree of clarity and
consciousness with which the foundation is stated, on the consistency
with which conclusions are drawn from this foundation, and on the wealth
of individual and social circumstances that are illuminated by these
conclusions.

Immanuel Kant, the greatest modern philosopher, held that ethics must
be a pure science of reason. He built on the principle that in our
action we should follow a rule which under the same circumstances can
hold for every other spiritual being. This is, to be sure, a great and
significant proposition; but it is not recognized by standpoints for
which one's own self, one's own family, one's own class, one's own sect
or one's own people is the highest. Its recognition has become possible
only through a long historical development (as Kant himself has also
shown in his philosophy of history). And even where it is recognized, it
is motivated differently (e.g. now more from the urge to self-assertion,
now more from the urge to self-surrender).

2) That the historical circumstances to which ethical propositions
% --- p. 63 ---
\setcounter{footnote}{0}%
\opage{63}are to be applied constantly vary and are difficult to survey is
as a rule willingly admitted. One is less inclined to admit that the
single individuals come with very different dispositions and
possibilities in ethical respect, and that therefore the task, as regards
both kind and degree, must be set differently for each individual.
\textit{When men have different capacity and strength, the same cannot
without injustice be demanded of all.} A trait of character which the
one displays without needing to summon up special effort or
self-control, the other can perhaps acquire only through strenuous work.
The strength which the one spends on producing works that benefit the
whole species is perhaps consumed by the other in the labor of keeping
his own head above water. Ethical judgments are so difficult to pass
because they presuppose a deep knowledge of the individual's inner
presuppositions, of his whole character. One does not notice much that
men are conscious of this difficulty.

This twofold difficulty (which, moreover, in different form and degree
holds for every science) is an expression of the great multiplicity and
ramified development of life, which means that man must find his way
practically in life without being able at every point to support himself
on a theory. The art of life precedes the science of life. The
individual does not, however, practice this art entirely on his own. He
lives himself, as we have seen (II), involuntarily into the experiences
and traditions of the species. How far he gets beyond what he has thus
involuntarily lived himself into will depend on the strength and
peculiarity he possesses and on the particular experiences he has. He
who breaks new paths must venture. The earlier experiences of the
species and the ethical theories can support him, but not lead him
completely, since he is precisely engaged in producing the foundation
for new experience and new theory.

\begin{center}V\end{center}


Ethical valuation is not the only way in which we can view human
qualities, actions and orders of life. --- We view them also
psychologically and
% --- p. 64 ---
\setcounter{footnote}{0}%
\opage{64}historically: describe them and try to understand them in their
connection with the characters, the peculiarities of race and the
conditions of life. Such a description and understanding is possible
without any valuation (at any rate conscious valuation) being made, and
it can be set in sharp opposition to ethical valuation. But in reality it
is so far from conflicting with the latter that it is precisely a
necessary condition of thorough valuation. What one would value, one
must first know and understand. With better understanding, valuation
will gain in thoroughness. The decisive difference between the
psychological-historical way of viewing and the ethical consists in
this, that the former turns toward the past, toward what lies before us
completed and can now be surveyed in its development, while the latter
turns toward the future, since the ultimate purpose and proper
justification of valuation lies in its becoming a means for actions and
orders of life to come in the future into greater agreement with the
fundamental values. --- We view actions and orders of life
\textit{aesthetically} or \textit{artistically} when we merely see them as
characteristic pictures, abbreviated and clarified expressions of what
life can bring. In a picture we enjoy even what we turn away from in
life. Art does not value (at any rate not consciously). It can give only
the single picture, the individual situation or fate; but whether it has
thereby given an example or an exception it cannot teach us, since it can
set up no general propositions. The living picture has its immediate
value whether it contains an exception or an example, and the aesthetic
or artistic way of viewing has also a mediate value of an ethical kind,
since it teaches us to view the scenes of life without egoistic
preoccupation.
--- The \textit{religious} way of viewing sees in the development and
fate of human life an expression of the struggle of values to persist in
the world of reality. The religious feeling can be described as man's
fundamental mood after the experiences he has had of the relation between
value and reality.

These different ways of viewing stand in constant interaction with one
another: science, art, religion and morality influence one another
mutually. In every ethics this
% --- p. 65 ---
\setcounter{footnote}{0}%
\opage{65}interaction will appear in some form or degree. Chiefly,
however, the peculiarity of a particular ethical conception will rest on
the relation it establishes between self-assertion and self-surrender,
and on the relation in which foundation (fundamental value) and content
(the special goods, tasks and virtues) stand to one another. ---
Historically there appear here (for European humanity) three main types.

\textit{In Greek ethics} the chief task was the limitation and
harmonizing of the unbridled promptings of nature. Hence great weight was
laid on moderation, self-control and self-knowledge. The single elements
in human nature were not to be cowed or extirpated, but to enter as
members into a living whole. Each single free man was to be a little
work of art, and the state a great work of art. The relation between
these two works of art (the individual and the political), and with it
the relation between self-assertion and self-surrender, was in the end
left unclear. The foundation was predominantly intellectual, and the
perfection of knowledge became for the thinkers the highest virtue and
the highest good.

\textit{Christian ethics} emphasizes strongly and inwardly the universal
love of mankind, although obedience to God's revealed will becomes the
highest virtue and the ultimate foundation. The relation between
self-assertion and self-surrender is not made clear here either, since
obedience (faith) is grounded in the wish for one's own salvation. The
intellectual and the aesthetic are altogether subordinate, in so far as
they are not regarded as sin, and there is a lack of interest in, and of
positive motives for, work on the development of human life through
spiritual and material culture.

\textit{Modern ethics} (from the time of the Renaissance on) demands a
development of human life such that all elements come into their own. In
particular it strives for a harmony between self-assertion and
self-surrender and demands room for both within social life, since only
so can all powers be released and work together. Evil here is neither
unbridledness (as in Greek ethics) nor pride (as in Christian ethics),
but sloth, which hinders unfolding and expansion. The demand for justice
as the condition for all personal powers being able to develop in
% --- p. 66 ---
\setcounter{footnote}{0}%
\opage{66}harmonious cooperation sounds ever more strongly both in
practice and in theory. At the same time a clearer consciousness of the
relation between foundation and content asserts itself in scientific
ethics.
--- In all these respects modern ethics has learned much both from
Christian and from Greek ethics, although modern European social life
and the experiences gained within it have exerted the most decisive
influence.

\begin{center}VI\end{center}


In our day significant oppositions assert themselves in the ethical
domain which more or less directly point to the problems that have been
discussed in the foregoing. These oppositions can be brought under three
groups.

1) The old opposition between \textit{theological and humane ethics}
constantly asserts itself, since the possibility of an ethical view of
life and conduct of life on a purely human foundation is often denied.
This denial, however, is not supported by new reasons, although it
appears in a sharper form than in the preceding age. There is the further
peculiarity that the opposition concerns rather what is set up as the
foundation of ethics (the fundamental value) than the content (the single
tasks, goods and virtues).

2) In literature the opposition between \textit{self-assertion and
self-surrender} is sharply prominent, since two brilliant writers (Fr.
Nietzsche and Leo Tolstòy) energetically maintain each one of these
tendencies in its strictest form. It is a symptom that the age needs
both ideals, which it is naturally easier to unite in pure theory than
in the practice of life. The task of ethical theory can only be to point
to the direction in which the practical solution must lie.

3) Within scientific ethics (where both the oppositions mentioned can
naturally also be traced) the chief thing to note is the difference
between a tendency which --- like the English doctrine of utility
(utilitarianism) --- mainly emphasizes the results of \textit{actions} in
their influence on men's weal and woe, and another tendency ---
essentially represented by German thinkers ---, which mainly emphasizes
\textit{the inner foundation of actions} in disposition and will. At the
beginning of the 19th century
% --- p. 67 ---
\setcounter{footnote}{0}%
\opage{67}Bentham and Kant stood as the chief representatives of these
two tendencies, and the discussion between them extends through the whole
century. A similar opposition appears in all scientific domains, since in
every science it is a matter both of objective confirmation by experience
(verification) and of subjective work of thought (in the setting up of
principles and in the founding of new points of view). Man lives both in
an outer and in an inner world, and the decisive proof of how ably he
lives in the inner world he can give only through the effects his inner
life leaves outside him.

% --- p. 68 ---
\setcounter{footnote}{0}%
\opage{68}\essay[Sociology and Ethics (1904)]{ON THE RELATION BETWEEN SOCIOLOGY \\ AND ETHICS\footnote[1]{Lecture to the Sociological Society in London, 11 Nov. 1904.}.}
\streg

1. The task of sociology is to study social life in all its manifold
forms of manifestation. Ethical ideals and ethical endeavors, therefore,
are objects of sociological research. They are working factors in social
development, while they are themselves symptoms of social conditions or
results of social development. They have their roots in the inner world
of individuals; but this inner world itself is not indifferent for
sociology, which traces the interaction of individual and society in all
its finest ramifications. The inner world does not develop itself
independently of the outer. Social conditions determine directly or
indirectly that which the individual conscience adopts as ideal or as
true. Very often the character and the direction of ethical life are
determined by physiological or social heredity. And even the fact that
individual initiative is at work sets no definite limit to sociological
research, any more than biology gives up the right to investigate the
organic variations which are the prerequisites of all natural selection.

From this point of view, sociology is a more comprehensive science
than ethics, which is a more special and limited science. Sociology
stands in a similar relation to ethics as does psychology. Ethical ideals
and endeavors are not only sociological, but also psychological
phenomena; they are, therefore, objects for
% --- p. 69 ---
\setcounter{footnote}{0}%
\opage{69}psychology as well as for sociology, and psychology is, in its
turn, a more comprehensive science than ethics.

2. Sociology is also a necessary foundation for ethics. The ethically
right must be sociologically possible, must be consistent with the
conditions and laws of social development. Ethics is not a system of
castles in the air, but a doctrine of the means and ways for developing
human life, individual and social, to greater richness and to greater
harmony. Ethics ought then to be founded on a study of the nature and the
conditions of the actual social development. Only with the help of this
study can new values be discovered and old values be maintained.


Since ethics is dependent on sociology, it follows that it cannot be
the same at all times, but must vary as its historical foundation varies
under different historical conditions. Ideals and motives, aims and
means, are subject to change. The historical development may lead to
ethical turning-points or to ethical dilemmas without any breach of its
continuity. There may be ethical discontinuity, though there is
historical continuity. I will give some examples.

Very often ethical development consists in this only, that what has
been done involuntarily, or even unconsciously, under the influence of
social heredity, is afterwards done with clear consciousness or as a
result of deliberate choice. Aristotle says that involuntary working
precedes voluntary and conscious working. It is, he says, by playing on
the zither that we become good players, and it is by acting justly that
we become just; in this way, the young are through education and tuition
introduced to ways of acting and thinking which later on can be followed
out with consciousness and free choice. When the time of involuntary
imitation and exercise is over, it remains to be seen whether the same
direction will be followed with full consciousness. But even if the young
generation breaks off and adopts quite other ideals and endeavors, yet
the first direction remains of great influence, both directly, as an
element, and indirectly, through an effect of contrast. The sociological
continuity is not broken, though new ideals and aims are acknowledged.
The Aristotelian principle (so I
% --- p. 70 ---
\setcounter{footnote}{0}%
\opage{70}name the law here spoken of) shows us a deeper connection, while
from a purely ethical point of view there might seem to be only
disharmony and opposition.

But social development may also in a more positive way be the
condition for ethical development. A thing which at first only had value
as a means may later on acquire immediate value --- value as an end in
itself. Social conditions make certain actions necessary, which the
individual would not undertake if they were not favorable to his own
interests; for instance, he must respect the liberty of others, if he
will have his own liberty respected. But later on he may adopt as his
immediate end the liberty of all. There is here going on what I will
call a subjective transformation of value. What is value at first only
as a means has become value as an end. An objective transformation of
value is going on when an end is so closely connected with the whole
order of things that it cannot be isolated. The original end may then
become a means to the acknowledgment of greater and higher ends, and
perhaps it may itself be swallowed up by them.

This concatenation of means with ends and of ends with other ends is
in great measure due to the social life and its institutions. By means
of these the individual who wills something is drawn into a great
process through which a whole series of transformations may arise. When
such a transformation has taken place, ethical ideals and aims are not
the same as before. New ethical formations have been produced through the
continual social processes. It is the task of sociology to show us the
conditions to which these transformations are due, and it is the task of
ethics to give form to the new ideals and to develop the consequences of
the new standpoint. All ethics has so far an empirical character. Every
step of evolution has its own ethics, and the ethics of one period cannot
be deduced from that of the preceding one without knowledge of the whole
intermediate social evolution.

The relation of sociology to ethics is here again similar to the
relation of psychology to ethics. The ethically right must be
psychologically possible, and there may be psychological continuity,
though from an ethical point of view it must be admitted that a wholly
new movement has been initiated.

% --- p. 71 ---
\setcounter{footnote}{0}%
\opage{71}3. Though ethics is more specialized than sociology, and though
it is in essential points dependent on sociology, yet it is an
independent science. This independence manifests itself in the fact that
social data, the results of historical development, are the objects for
ethical valuation. In ethics these data and results are examined in their
relation to an ideal, a standard; and their value depends on their
harmonic or inharmonic relation to this ideal. The main difficulty of all
ethics consists in the establishment and the justification of the ideal
which is to give the standard of valuation. The great struggle between
ethical systems concerns this establishment. Ethics has here problems
which sociology, as such, as a purely descriptive and causal science,
does not know. Sociology looks backward to discover the course and the
laws of social evolution. Sociology comes after reality. A social process
must be to a certain degree finished, before its law can be determined.
We may say of sociology, what Hegel has said unjustly of philosophy in
general, that the bird of Athene first begins its flight in the
twilight, when the labors and the struggles of the day are over. But in
ethical thought the goddess herself appears on the field of battle. She
does not forget what she has learned of the struggles and labors of
former times, but she applies this knowledge in her war to lead the
ideal she has won to victory.

Historical data are the foundation on which ethical development is
based. The great question is, in what manner and in what direction the
conscious continuation of what history has won is to proceed. Here
individual stands up against individual, individual against society,
society against individual, and society against society. We have here all
the four sorts of war which Hugo Grotius long ago distinguished. Ethical
valuation must always be the work of individuals who start from certain
definite social and mental presuppositions. How, then, can a universal
ideal be formed, and how can the validity of a definite standard be
justified? Here we have the real sting of the ethical problem. It is the
greatest war in history, which here is carried on in the quiet world of
thought. This war has also interest for the sociologist; but he is here
only a spectator. The moral philosopher takes his place in the battle
itself.
% --- p. 72 ---
\setcounter{footnote}{0}%
\opage{72}The sociologist examines what is going on, and how it is going
on; but the moral philosopher asks on which side the highest value is to
be found, and how he can get a standard to test this value. The moral
philosopher can as little set himself above sociological laws, when such
laws really exist, as the agriculturalist can set himself above chemical
laws. But as the agriculturalist can make use of chemical laws in order
to make the earth produce the profit he looks for, so the moral
philosopher can make use of sociological laws in order to produce
ethically valuable results.

The independence of ethics manifests itself in the selection of ends
and means within the manifold possibilities which sociology presents.
The farther we advance on our way from sociology to ethics, the more the
field of possibilities becomes narrowed. There are more possibilities in
the marble block than in the sculptor's artistic plan.

The relation of sociology to ethics is here again similar to the
relation of psychology to ethics. Psychological possibilities present to
ethics the same problem as was presented to it by the sociological
possibilities. In both respects what matters is to find a determining
principle for the choice.

4. But not only is there a difference between sociology and ethics,
but there may be a sharp contrast between them, and it is important to
lay stress on this contrast --- the contrast between valuation on the one
side, and description and explanation on the other.

If we efface the distinction between the sociological and the ethical
points of view, we are led either to regard the results of development
as valuable in and for themselves, or to presuppose the full realization
of the ethical ideals as given. In the first case we lower the ethical
demand; in the second we give an idealistic misinterpretation of the
actual. In the first case, sociology swallows up ethics; in the second,
ethics swallows up sociology.

But it is quite possible to maintain the difference between ethics and
sociology, and yet to assert the intimate connection
% --- p. 73 ---
\setcounter{footnote}{0}%
\opage{73}between them. Sociology leads us on to ethics by the application
of the comparative method. The comparison of social forms or social
states naturally leads us to characterize some as ``higher,'' others as
``lower.'' This characterization necessarily presupposes a certain
standard. We call a form of society higher than another, if it more than
this other makes it possible to attain two ends at once, namely, the free
and rich development of individual peculiarities and differences, and
the realization of unity and totality in social life. In sociology, as in
biology, the standard is this: the intimate connection of differentiation
and concentration. This has led to a comparison between society and an
organism, and great scientific profit has been expected from this
comparison\footnote[1]{Cf. the essay on the concept of analogy in \textit{Mindre Arbejder. Anden Række} p. 40.}.

There is certainly also an analogy present, when we call organisms
higher and lower, and when we call societies higher and lower. Every
science whose objects present both unity and multiplicity must in its
comparisons make use of such a standard as sociology and biology exhibit.
So is it, for instance, also with psychology. We call a personal life
high if it exhibits at once richness (of endeavor, emotion, and ideas)
and firm concentration (character); we call it low if it is poor and
incoherent. The task of biology, psychology, and sociology is only
description and explanation of facts. Comparison is here only a
methodological means. Though comparison presupposes a standard, and this
standard can be constructed as an ideal, so that beings or species are
called higher or lower according to their relation to this ideal, yet
such a manner of view has in those three sciences only the character of a
working hypothesis.


But in ethics this manner of view has not merely methodological value.
Here the standards and the ideals have also a positive and practical
significance. The most intimate connection
% --- p. 74 ---
\setcounter{footnote}{0}%
\opage{74}between unity and multiplicity in the life of the single
individual and in the life of society is in ethics an end which ought to
be reached, and it is the task of ethics to find means and ways which can
make an approximation to this end possible. The highest aim of social
ethics would be an empire of humanity in which there is the greatest
possible multiplicity in the development of personal life and, even by
this very means, the greatest intimate union of personal beings. If the
single individual in developing himself in his own peculiar way gives the
best possible contribution to the whole life of society, and if, on the
other side, society is organized in such a manner that a free and full
development is possible for all individuals, then we are approaching to
the ethical ideal. Ethical imperatives are only logical and psychological
consequences of the acknowledgment of this ideal\footnote[1]{See on this ideal my \textit{Etik} XIII, 6., and on the possibility of proving its validity III, 8---15. --- Cf. also the essay ``Humane Ethics'' above.}.

The so-called social problem is also an ethical problem, and this
gives it its own particular sting. A social problem arises when
multiplicity --- by progressing division of labor, for instance ---
prevails in such a degree that the free development of individuals is
checked. There would not be any sting in this problem, if our ideal did
not claim that every man shall be treated, not only as a means, or as a
part of a machine, but also as an end in himself. If this ethical
principle is not presupposed, social factors may be interesting objects
for science; but such science will be sociology, not ethics.

5. I have already said that the difference between ethics and
sociology cannot be said to consist in the fact that ethics has for its
object the development of single individuals; sociology, that of
society. There is nothing in the life of individuals which may not be of
interest for sociology; and, on the other hand, ethics is not only
individual, but also social. But it cannot be denied that the point of
view of ethics causes it to accentuate, in a higher degree than sociology
has occasion to do, that which is going on in the inner world of the
individual. Even if a system of ethics maintains the social
% --- p. 75 ---
\setcounter{footnote}{0}%
\opage{75}point of view, it cannot forget that society consists of
individuals, and that the welfare of society is the welfare of these
individuals. Only in the consciousness of individuals can the value of
life be experienced. The concept of society does not, therefore, lose its
importance. The importance of this concept is similar to the importance
of the conception of potential energy in physics. It notes conditions and
possibilities for the unfolding of human life which are more
comprehensive than the horizon of any individual --- forms and germs of
life by which many generations can profit. It contains the heredity of
the past, organized results of the experiences of former generations ---
and at the same time dispositions and possibilities for the future. It
expresses that which persists and continues, in opposition to the
shifting interests of single individuals, groups, and generations. But
these potential values can be actualized only when they are appropriated
and worked out by particular individuals. ---

Not all ethical theories give this point of view its full right.
For Hegel the main point was what he called ``the ethical substance,''
which works itself out in the different forms of society, most typically
in the state. In comparison with this social substance, the existence of
the single individual is indifferent. The perfection of the individuals
is to live and breathe in the great whole of society; but the essential
in the ethical world is something which transcends the consciousness of
any single individual. Though Wundt accentuates the importance of
individual will more than Hegel, a similar view is found in his ethics.
For Wundt the total will, ``der gesammte Wille,'' which manifests itself
in the existence of society, is the mightiest of all facts. It is
imperishable, and it is always right. The single individuals are
perishable; with all their endeavor and all their capacity for happiness,
they are only drops in the ocean!

Such theories overvalue potentiality at the expense of actuality. It
is true that actuality always lives on potentiality, and that individuals
live and work in virtue of social conditions. But that potentiality
exists is known only from the fact that it can be transformed into
actuality, and on this depends its value. The value of Hegel's ``social
substance'' or of Wundt's
% --- p. 76 ---
\setcounter{footnote}{0}%
\opage{76}``total will'' consists in their making a rich personal life
possible in great groups of individuals, and only through the study of
what is going on in such groups is the construction of these concepts
possible.

For Hegel and Wundt sociology and ethics are ultimately identical --- a
consequence of the absolute subordination of the individual point of
view to the social. The same may be said of August Comte in the first
period of his philosophizing. In his ``Cours de Philosophie positive'' he
did not regard ethics as an independent science which has its special
place in the series of the sciences; ethical ideas are here to be found
partly in the biological, partly in the sociological, chapters of his
work. But later in life Comte regarded ethics as an independent science,
the seventh and last in the series of the fundamental sciences. The
principle on which Comte arranged the sciences in a series was that the
following science shall always be more concrete in its object and more
inductive in its method than the preceding. When, then, in his later work
he puts ethics after sociology, he presupposes that ethics is the more
concrete and inductive science of the two. Now, Comte's reason for this
view is that in sociology the individual motives and tendencies
neutralize one another; it is the average results which are of
sociological importance; while in ethics the character and the tendency
of the inner life and the individual realities in their multiplicity have
the first place.

I do not myself believe that Comte deduced all the consequences of the
position which he finally assigned to ethics, nor that he gave a complete
concept of the relation between the individual and the social points of
view in ethics; but he points in the right direction. It is the strength,
but it is also the weakness, of ethics that it is the most concrete of
all sciences. It stands almost at the boundary between science and art.

To conduct life ethically is the greatest of all arts. And, like all
arts, it develops itself spontaneously. All that theory can do at first
is to learn from this spontaneous development, to find its moving forces,
and to formulate the thoughts which it presupposes. Later on, the art can
be exercised with greater
% --- p. 77 ---
\setcounter{footnote}{0}%
\opage{77}consciousness, and there can then be an interaction between
thought and life. And such an interaction cannot be established if ethics
does not --- all the differences in point of view and in methods
notwithstanding --- remain in indebtedness both to sociology and to
psychology\footnote[1]{Since this essay appeared, the relation between sociology and ethics has been the subject of much discussion, occasioned especially by the position of the new French sociological school on this question. Of great interest here are especially \emph{Durckheim}'s statements in ``Bulletin de la société française de Philosophie'' (Vol. VI) in a discussion of the concept ``le fait moral'' (1906), and \emph{Levy Brühl}'s book ``La Morale et la Science des Moeurs'' (4. éd. 1910). On the basis of my conception of the ethical problem I have treated the question anew in ``Den menneskelige Tanke'' (1910) p. 350---369. (1912).}.

% --- p. 78 ---
\setcounter{footnote}{0}%
\opage{78}\essay[A Case History (1901)]{A CASE HISTORY.}
\undertitel{\emph{Hannah Joël}: \textit{Hinsides}. En psykologisk Redegørelse. Copenhagen 1900\footnote[1]{Hannah Joël was the pseudonym of Helga Johansen (the sister of the painter Viggo Johansen). († 1912).}.}
\streg

I\textsc{t} is neither a novel nor a religious treatise that lies before
us here, though from the title page one might be tempted to expect either
the one or the other. It is a suffering human being who tells of her
struggle with mental illness, of her defeats and of her victory. Rarely,
no doubt, does one, outside the circle of psychiatrists, form any living
idea of what inner struggles must be fought by those whose mental life
has come out to the border of mental illness. The dark moods, the violent
effect of all sense impressions, the rising hallucinations and the sudden
impulses to movement must constantly be held down by the explicit idea,
consciously held fast, of where reality lies and what reality demands.
What ordinary, healthy or normal people learn to know only exceptionally
from their own experience, the difficulty of distinguishing between the
world of dreams and the world of reality, becomes for such people a
constant task, and a task that makes it impossible to solve other tasks;
indeed, that very task they are perhaps not able to solve.

In the present book such a case history is told in great detail and
with no small psychological and literary talent. On a first reading the
book perhaps puts the reader's patience to the test; but the second time
he will
% --- p. 79 ---
\setcounter{footnote}{0}%
\opage{79}admire the wealth of single observations and the delicacy of
description which the author commands. The book will, better than any
scientific treatise, be able to further the understanding of the sick
life of consciousness, and that is something that is needed in a high
degree. For even if it should not be true that mental illness is
increasing, it is nevertheless more frequent than had formerly been
believed. And moreover the book teaches not only about things that
concern the mentally ill. All can know from their own experience the
power of melancholy, of cares and of illusions, and it is therefore no
wholly strange world into which the book leads.

Hannah Joël did not regard herself as mentally ill when, on the advice
of those around her, she took refuge in the mental hospital. Her
condition she characterizes especially by the violent way in which all
sense impressions act, both the impressions from the outer world and
those from her own body. She had to struggle to collect her thoughts. She
wished only to be able to close her eyes and ears to everything and only
``look inward,'' so as to be able to have done with what was working in
her thoughts. Even when all is quiet around her, the heart still beats
and destroys the peace. To get rest she has used opium, and she praises
this remedy because it does not, like wine, call forth ``a stream of the
feeling of life,'' but merely damps the painfulness, so that thought may
be allowed to live its life in peace. From various hints one can infer
that she has pursued a study, and that it is in order to be able to
concentrate on this that she has taken refuge in the dangerous remedy.
Now it no longer helps her, and although it is not said in plain words,
it is nevertheless an opium withdrawal cure that she undergoes at the
hospital.

At the hospital she wages a long double struggle, --- a struggle against
the inner disturbances and a struggle against the systematic treatment
which doctors and attendants would force her under. The inner chaos
threatens to overcome her, and from the other side she meets the firm
will which demands that a definite order be observed. In her delirium
she understands neither the one nor the other; both seem to her equally
to threaten her personality, and against both she appears with a
% --- p. 80 ---
\setcounter{footnote}{0}%
\opage{80}feverish self-assertion, from which she sinks back again and
again into torpor, fits of rage or hallucinations.

She feels an aversion to being placed, as a patient, under a definite
heading. Her case and condition have for her such a peculiarity that she
resents the ease with which they manage to classify her. She thinks that
the doctors make the personalities conform to the headings, instead of,
conversely, making the headings conform to the personalities. It is a
struggle that is no doubt fought again and again. The patient feels
wronged at having merely to stand as an example of a generic concept
which has been formed beforehand, and which perhaps is not even now
enriched by a new nuance. A rebellious mood, a constant irritability and
irony were the consequence.

No less hard was the other struggle, the struggle against the uneasy
moods, against hallucinations and impulses. Attempts at suicide and fits
of rage alternate with restless agitation and chaotic fantasizing. The
description of the fits of rage surely belongs to the best in the book.
The impulse comes so quickly and overwhelmingly that the patient first
discovers what she is doing by seeing her hands clasped around the throat
of another person. There is no particular ill will toward the one it
falls upon; during the fit the victim does not exist at all as an object
for the patient's consciousness. --- Excellent, however, is also the
description of the hallucinations. There are voices that speak; ---
could these be hallucinations? she asks herself, but rejects this
possibility because the sound is so distinct, although she also heard
voices of people who were demonstrably not present. Sometimes it is a
single word or a turn of phrase that has been used to her and that now
goes on echoing in constant variations, accompanied by hallucinations of
sight and hearing. The hardest struggles are fought against a host of
figures that press forward from behind. It is then a matter of having
one's back free, up against the wall or the stove, and of getting the
figures out in front of one in order to fight them better. The agitation
and restlessness of the patient arise from her having at any price to get
the better of ``what is behind.''

The turning-point in the illness comes suddenly and
% --- p. 81 ---
\setcounter{footnote}{0}%
\opage{81}surprisingly. Whether this is grounded in the natural course of
the illness, or whether it is connected with an imperfection in the
presentation, I cannot judge. Just after severe fits the conviction of
being well emerges, and a walk is used to flee from the hospital. Hannah
Joël feels that she must now take her fate into her own hands and tear
herself out of the irresponsible condition which life in the hospital,
this ``living grave of self-abandonment,'' brings with it. Neither the
doctors nor her nearest are to bear responsibility for the step she
takes: ``I must bear the guilt alone. The step that must be taken, and
that the others do not dare to let me take when I merely ask, I must
quickly take myself, myself take and give myself freedom.'' --- But a hard
struggle still remains after she has returned home. Through long
sleepless nights there must be struggle with shadows and voices. And even
by day it is only ``in jerks'' that the full reality of the surroundings
can be held fast. Yet it goes forward foot by foot. On a lonely Christmas
Eve the last struggle is fought to keep herself up. The bright points of
life and its dark points can each be gathered into stripes; healthy
people see bright stripes; in mental illness the dark points gather into
ruling stripes; --- it is then a matter of holding the bright stripes
together, although the dark points threaten to dissolve them. ``I choose
to try life in the light again, and yet remember what I learned beyond!
Whether I can manage to get through, I do not know. I must only try
humbly, let go all former arrogance, venture forward dot by dot.'' Now
there stirs also the urge to give and receive love, to feel oneself in
living relation to people. Thus she captures ``one by one the elements of
life,'' builds her house anew, the house that had collapsed through
sorrow and illness. ---

Perhaps, as already remarked, the transition from illness to
health, from ``the beyond'' to present reality, from the dark to the
bright ``stripes,'' should have been depicted in greater detail. But the
author has not wished to give transitions and explanations where her own
observations showed none. And this is also what one can especially learn
from this book, that the will to hold life together must be born, come
into being
% --- p. 82 ---
\setcounter{footnote}{0}%
\opage{82}through a process of life which for a great part runs its course
in hiding, in the unconscious. That what builds the world for us, --- the
world which the healthy and normal readily take for an all too
self-evident thing, --- lies in the end in our own inner self, is
experienced in earnest by him in whom the foundation tottered and sank,
so that it had laboriously to be joined together again, stone by stone.
And this whole process, both the dissolution and the joining together,
could only be depicted as it has been done here by one who was a thinker
and an observer, and who after the dark experiences had the ability and
the urge to give new birth to what had been experienced in a recollection
full of mood and understanding.

% --- p. 83 ---
\setcounter{footnote}{0}%
\opage{83}\essay[Plato's Symposium (1905)]{PLATO'S SYMPOSIUM}
\undertitel{Lecture at the first meeting of the ``Græsk Selskab'' [Greek Society].}
\streg

W\textsc{ith} great misgiving I take the floor here this evening, and for a twofold reason. How can one in a worthy manner open the first gathering of a circle that meets in order to be strengthened and renewed in understanding and admiration of some of the greatest works the human mind has brought forth? And how in particular can one speak worthily of one of the most brilliant of these works, a work in which the splendor of poetry, the plastic power of the portrayal of character and the profundity of thought have entered into a unique union? The art of interpreting the works of art, and of making them vivid through interpretation, stands high in this country; I need only mention names such as Julius Lange and Georg Brandes. I know that I do not command this art. If I have chosen to speak about Plato's Symposium, it is because this work is a classical expression of Greek thinking and can at the same time serve to illuminate the kinship and the opposition between ancient and modern intellectual life. The subject is thus in itself well suited to open our gatherings. It is to be the task of our Society, after all, to immerse itself in antiquity in such a way that its lasting significance through the ages, beyond its national and historical limitation, is also brought out. --- It is also a debt of gratitude that I am paying off in honoring, in my own way, this work of antiquity. It is now 45 years since I first read it, but it has accompanied me on my way, and I have returned to it at very \opage{84}different points in my life. Often it has called up in me the memory of the dark, low-ceilinged schoolroom in which, in my last school year, weighed down by reading for examinations, I worked my way through it. It was the bright spot of that winter. I was far from understanding it fully; but there were thoughts and impulses in it that have accompanied me through the years, and that I have understood more clearly with each new reading. If only I may now succeed in making clear to others what the work has been for me!

First of all I must attempt an account of the content of the work, of the frame in which the development of the thought is set.

\whiteline

A young Athenian, an ardent follower of Socrates, tells of a banquet that he had attended several years before at the house of the tragic poet Agathon, and in which Socrates, the comic poet Aristophanes, the physician Eryximachus and several others took part; toward the end Alcibiades joined them. Agathon had won a victory in a dramatic contest, and on the day after the great feast he had given for the actors, he gathered a few friends for a smaller celebration. Several of the guests had been present the day before and were feeling its after-effects; it was therefore easily agreed not to hold any real drinking bout, but to discuss a subject, and it was arranged that each in turn should deliver an encomium on Eros.

The rhetor Phaedrus, from whom the idea of this subject came, was the first speaker. He recalls that Eros, according to the account of some poets, was one of the oldest gods, and then dwells on the inspiring influence of the urge to love and to friendship. The Greeks knew, between men too, a relationship of love of a passionate and intimate kind, and it is above all the significance of this relationship that Phaedrus stresses. The urge to win honor and to avoid degradation in the eyes of the beloved is the source of the greatest things in public and in private life. No other feeling is able to such a degree to create enthusiasm for great deeds. No man is so wretched that Eros could not inspire him to action and to measure himself against the best, even should it cost him his life.

% --- p. 85 ---
\setcounter{footnote}{0}%
\opage{85}Pausanias, one of the following speakers, maintains that there is both a heavenly and an earthly Eros, just as there is a heavenly and an earthly Aphrodite. Praise can apply only to the noble, the heavenly Eros, which is not ensnared in the sensuous but is able to call forth lofty thoughts and to found an ideal friendship that can last throughout life, because it is first and foremost built on admiration for character.

This distinction between a noble and an ignoble Eros is approved by the physician Eryximachus. But he adds that they contend not only in the souls of men but in all that lives, indeed in the whole of nature. It is the task of the art of medicine to search out the noble Eros; it is only with its help that the art of medicine can do its work: to bring agreement and harmony between the different elements of the body, just as music seeks to tune the elements of the soul to harmony. In great nature too there is found now a harmonious, now a disharmonious relation between the elements; here too the noble and the ignoble Eros contend: the former sends fertility and good seasons, the latter bad years and failed harvests.

Aristophanes has this objection to make against the preceding speakers, that they have not explained how the erotic urge came into being. His explanation is given in a myth: We human beings are not now constituted as we originally were. Originally men were twice as large as now, but their mighty figures rose in arrogance against the gods and threatened to storm Olympus. The gods were loath to exterminate men --- for then there would, after all, be no one to pray to them and bring them sacrifices. So they hit upon the expedient of splitting them lengthwise --- as one cuts eggs through. Since that time the halves now go about in longing and search for one another, and great is the rejoicing when they meet. Eros, then, is the urge to completion, to become a whole again. But we must be pious, lest the gods divide us anew, so that we become like figures in relief! There is good reason, then, to praise Eros, who takes away our unrest and stills our want!

After the comic poet comes the tragedian Agathon, who thinks that the preceding speakers have rather called men happy for the gifts of the god than portrayed the god himself. And so \opage{86}he praises in flowery words and rhetorical turns the glory and beauty of the god.

Now at last the turn comes to Socrates. He turns at once to Agathon and asks whether he is sure that Eros is so glorious and beautiful, indeed, that he is a god at all. Eros means, after all, urge or impulse, and we have need of what we do not have. Eros, then, is not beautiful and is not a god, but is the urge toward the beautiful and the divine. Thus the intellectual urge expresses itself neither in those who have wisdom nor in the wholly ignorant, but in those who stand between. No god philosophizes, i.e. strives after knowledge, but men philosophize as soon as they feel their ignorance with pain. Eros is the son of Poros and Penia, of plenty and want; hence his two-sided nature. He moves between the two opposites, but coincides with neither of them, and always begins his course anew from the one to the other. This holds of all urge and impulse. Not only in love and friendship and in sensuous passion does Eros express himself, but in all human striving. And he knows no final limit. He wants to reach his goal not merely for a time but eternally; therefore the urge is born ever anew. No single object can satisfy it. In order to understand his nature rightly, we must consider that a finite being cannot become eternal. How, then, can the urge toward the eternal be satisfied in a finite being? Through the production of a work in which its life is continued. In the passion that expresses itself in all living beings in propagation and in the care for offspring, the urge toward the eternal is at work: if the single being cannot go on persisting, it can still leave behind another being of its kind. Even the single being, after all, persists only under a constant change of all its elements; this holds of the thoughts of the soul as well as of the substances of the body. But in nobler natures the urge toward the eternal does not express itself in this way only. It is at work also in the urge for honor and for producing works in art, in wisdom of life, in statecraft, that can endure beyond the limits of the individual life. And in its highest form it expresses itself in thought, when thought rises above everything particular and sensuous and forms great ideas in which long series of observations are summed up, ideas that also stand \opage{87}as exalted models of what observation always shows us only in an imperfect way. Eros as the urge of love and as the urge for honor is only a preparatory school for this highest urge, which step by step leads us up into the world of ideas and ideals, to thoughts that cannot be expressed in images or sensuous forms. The possibility of this ascent is what gives life its highest value.

Just as Socrates has reached this summit, a knocking is heard at the street door, and the sound of voices is heard. It is Alcibiades with a following of drinking companions, who is coming from another banquet. He is very drunk and shouts to know where Agathon is; he wants to get in to Agathon to crown him on the occasion of his victory. He is then received among the guests and lies down at table, but starts up at once when he discovers Socrates. After some jesting words he crowns him too: Socrates was, after all, a greater victor than Agathon. He is then set the task of delivering an encomium on Socrates. But how is one to portray Socrates, who resembles no other human being? Why, one can compare him to the cabinets that the sculptors have in their workshops, which are shaped in the form of satyrs or sileni and are used to keep images of the gods in. This fits Socrates exactly, whose outward form is as ugly and bizarre as his inner self is noble and glorious, and whose seemingly insignificant talk contains, on closer consideration, the deepest thoughts. And just as the satyr Marsyas enchanted men with his playing, so Socrates does with his speech. Alcibiades himself can bear witness to this: ``When I hear him, my heart beats harder than the Bacchants', and my tears flow. I feel disgust at my accustomed ways, and since I nevertheless cannot desist from them, I shun Socrates, indeed sometimes wish that he did not exist, and yet that would be the greatest grief of all to me! What am I to do with this man, the only one before whom I feel shame!'' --- Not least does Socrates' satyr character show itself in his irony, especially in his relation to the young. It looks as if he were taken with them for their beauty or other outward advantages, and yet he feels contempt for all such things. It looks as if he were jesting and playing, and yet there is \opage{88}the deepest earnestness in his thought. Alcibiades himself makes an unreserved confession that he had thought he could ensnare Socrates, but testifies that it was just then that the man's purity and loftiness appeared in the clearest light. Then he truly caught sight of the images of the gods within him in all their wonderful beauty. And this loftiness and self-mastery Socrates showed in all the circumstances of life. Alcibiades dwells especially on scenes from wars in which Socrates' hardiness and fearlessness showed themselves in an admirable way.

When Alcibiades had finished his speech, more drinking companions arrive. The feast now takes on a more chaotic character. Soon some leave, while others fall asleep, among them the narrator himself, who does not wake until toward morning. The only ones awake now are Socrates, Agathon and Aristophanes, who are discussing the relation between comedy and tragedy. Socrates maintains the claim that comedy and tragedy are akin, and that it lies within the power of the same man to create works of both kinds. Agathon and Aristophanes soon fall asleep, but Socrates settles them comfortably and then goes out to take his bath, whereupon he goes to his day's work.

\whiteline

It was of Eros that there was to be speaking at the banquet, and then it was revealed, when Alcibiades, inspired by the wine-god, arrived, that Eros was in their very midst in the form of the great seeker after truth, who was more worthy of the wreath than the poet to whom the whole people had given the prize. It was precisely through his lofty urge toward truth and greatness that Socrates had become the model for the young, the object of their love and their longing. To these young men Plato himself had belonged. He, like Alcibiades, has seen the images of the gods within his great teacher, and from that moment on he has dedicated to him all the best fruit of his thought, has ascribed it to him, however much it was owed to his own inquiry. In almost all his books he lets Socrates do the speaking. Never has one who was himself a master so made himself the instrument for the glorification of another great mind. But it was also a new ideal that in Socrates' personality dawned upon the \opage{89}world. Alcibiades was right that one could not compare him with any known figure. He stands by himself --- and precisely thereby he became a model through millennia for everyone who puts his personality into lofty striving and has his life in the constant yearning from the finite toward the infinite and in the endeavor to awaken the same striving in others. To be sure, very little is known of Socrates' own thinking and teaching. It has gone with him as with an angel in Dante's Paradise, who is hidden by the light he himself sheds. Plato has caught this light and sent it to us in a multiplicity of rays (the richest here in the Symposium), and even if all we knew of Socrates were that it was he who inspired Plato to the portrayals he has given in the ``Symposium'' and in other dialogues, we should really know enough.

And yet there is one thing that can with certainty be ascribed to Socrates, a discovery that was ever anew the object of Plato's enthusiasm. It was the quite simple thing that general concepts can be formed, in which, as in a collected expression, the essential content of long series of observations is comprised. It was the possibility of gathering observations into concepts that gave Plato the impetus to his distinctive philosophy, the so-called doctrine of ideas. Thus the ``Symposium'', after all, treats of what is common to all urge, all striving, all Eros, from their lowest to their highest forms, and the highest form of striving became for Plato precisely the one that leads up to the idea from the single observations.

The ideas were for Plato not merely thoughts that we form in order to orient ourselves in the world; they are themselves the true reality, the constant, the fixed, the abiding, in contrast to the constant change of all sensuous phenomena. The world of ideas was the true reality, and the world of sense only an imperfect copy of it.

It was owing partly to Plato's poetic nature, partly to the connection of incipient thinking with mysticism and mythology, that Plato conceives the ideas as fixed, eternal beings, unchangeable, raised above the world of experience. Properly speaking, he here came into conflict with his own doctrine of Eros: for if constant \opage{90}seeking is the highest, if the urge can never be fully stilled, how then can the final goal be placed in an eternal and unchangeable being such as the one ascribed to the ideas? Is that striving, then, only illusion? In the doctrine of Eros and in the praise of Socrates we hear a constant Excelsior! --- but the world of ideas stands eternal and unchangeable, like the old Olympus. It is characteristic of Greek intellectual life that it so quickly comes to rest in plastic, calmly reposing forms, and ascribes higher value to the unchangeable than to the changeable. ---

Both the doctrine of Eros and the doctrine of ideas have their counterparts in modern times, but it is of great interest to see what changes they have undergone.

Eros, the constant urge, the eternal striving, we find again in modern times especially in such periods as those in which the horizon is widened and the barriers to development are cleared away. And it then appears freed from the ancient limitation, but often also with a knowledge of resistance and pain, a tragic element that accorded well with Plato's theory of the inner connection between tragedy and comedy, but for which he found no place in his doctrine of Eros itself. On the ground of the Italian Renaissance Giordano Bruno expresses the thought of the eternal striving, now in purely philosophical form, now in poetry. For him it has its value in itself, even should it not lead to the goal: ``Even if the goal is not reached, and even if the soul is consumed by the intense yearning, it is enough for it that so noble a fire has been kindled in it''. Heaven is farther off than it was in antiquity ---, but this has not quelled the urge to ascend. On the ground of the German Renaissance at the close of the eighteenth century we find the great Eros in philosophical form in Lessing and Kant, in poetic form in Goethe's ``Faust''. Lessing's choice of the constant striving after truth, even with the possibility of going astray, if only the powers of the soul are constantly enlarged, contains, like Bruno's words, a consciousness of the possibility of disappointment and of the seriousness of the struggle, which Plato's young, sanguine Eros does not yet have. And Kant concludes his mole's work in the investigation of the nature of human thinking with the result, which in him calls up the memory of Plato, that from the nature of our thought there arises an ideal that no experience \opage{91}can satisfy, but that for that very reason leads us on to ever new and enlarged experience. Goethe's ``Faust'' finds happiness in constant striving, although no single moment brings full satisfaction, and it is this constant striving that makes it possible for the heavenly hosts to save him.

A corresponding change has come over the doctrine of ideas. In place of the ``idea'' there enters in modern times --- already in Baco and Bruno --- the law. It is no longer a resting concept, a fixed type, that is to contain our experiences. Galilei, the founder of modern science, expressly maintains that it is unjustified to ascribe greater value to the unchangeable than to the changeable. The word law means the rule for the changes in the world; it concerns what is to be expected when certain conditions are present. The idea leads only to classification, to the ordering of unchangeable phenomena, but the law gives an explanation, a How and a Why, not merely a What. The task of seeking the constant, the eternal, in the midst of the world of change has become greater, but it is nevertheless a continuation of the task that Greek thinking set itself.

Of both of Plato's principal thoughts it may be said, then, that they have the fate that they shall not die, even though they must undergo thorough metamorphoses.

But there is yet another thought that comes forward in the ``Symposium'', and that still lives. Plato shows how the one Eros, the one urge or impulse, can be transformed into the other: the urge of love, the urge for honor, the poetic urge and the urge for knowledge --- they form for him what with a modern expression one might call psychical equivalents. It is displacements of motive that he describes in his doctrine of the great Eros. There is an ideality in every urge or impulse, insofar as it points beyond the given state, beyond what has been attained, and makes a demand. And in this Plato sees the possibility of a metamorphosis from one impulse into another. Eros can here become the great teacher, in that the impulse, once it has been awakened, can be transformed into new and perhaps higher forms. Thus spiritual development becomes possible, not through asceticism or laborious self-compulsion, but through an \opage{92}inner transformation. The earthly Eros can pass over into the heavenly. They are not such complete opposites as some of the speakers thought.

It is surely a personal trait of Plato's that reveals itself to us in the ease with which this metamorphosis of impulse takes place in his portrayal. It is not in all his dialogues that the metamorphosis is so easy; but it is perhaps Plato at his height, the real Plato, that we have before us here. It seems to me that Gomperz is right in maintaining this. In the dialogues that precede, a one-sidedly intellectualistic conception prevails, and in later dialogues a more mystical and dualistic conception prevails. In the ``Symposium'' he does not dwell on the struggles and sufferings, the crises and catastrophes, that can arise from the conflict of the impulses or from the power of single impulses. It was, after all, only after the portrayal of Eros was concluded and the banquet over that he had Socrates speak of the connection between tragedy and comedy; the tragic element he did not, as I have already remarked, take up into his doctrine of Eros. For Plato himself the intellectual Eros, the urge to ascend to comprehensive thoughts, was surely the highest. This Eros had been awakened in him by Socrates in his young days, and it accompanied him throughout his long life. It cast a radiance over life even when life darkened before the old thinker's eyes as he looked at the course of things in the world. Never has anyone been able to give so brilliant and so profound an expression of what for him gave life value.

The idea of psychical equivalents plays a significant role in modern psychology as well as in psychiatry. In human development feelings and impulses may arise and disappear; but on closer investigation we may find that they arise as a new form of preceding manifestations of life, and that they live on in new, apparently quite different manifestations of life. Plato himself remarked, after all, that a man persists under a constant change of the elements of the soul as well as of the parts of the body. It is in any case, then, a task that presents itself, whether the new stages can be explained from the earlier ones through \opage{93}displacement and metamorphosis.\footnote[1]{Cf. my ``Mindre Arbejder''. I. p. 226---230.} What significance this acquires for pedagogy and criminology, when they want to be something other and more than instruction in outward drilling, I need not specify more closely. Plato has shown here that, however much he wandered in the ideal world of thought and art, he nevertheless had his footing in the world of reality.

% --- p. 94 ---
\setcounter{footnote}{0}%
\opage{94}\essay[Democritus and Plato (1909)]{SOME REMARKS ON THE RELATION \\ BETWEEN PLATO AND DEMOCRITUS}
\streg

A\textsc{ttention} has been drawn to this relation by Mrs.\ \textsc{Hammer-Jensen's} doctoral dissertation \textit{Den ældste Atomlære}. The author seems to have succeeded in showing that the conflict that appears in \textsc{Plato's} \textit{Timaeus} between teleological and mechanical explanation of nature is due to an influence of \textsc{Democritus'} writings, an influence that makes itself felt in the dialogue from a definite passage onward (p. 45), breaking the connection in a strange way. The detailed investigation of this is the contribution that the dissertation mentioned makes to the illumination of the relation between the two Greek philosophers. One may agree with the author on this point, as well as in her great admiration for Democritus, even if one, as I for my part must do, keeps a very great distance from her account and her assertions on other points.

It is by no means the case that --- as is asserted in the dissertation --- philosophers should be inclined to regard Democritus merely as a physicist. In all the accounts of the history of philosophy that I know, he is regarded as a philosopher, and with full justice, even if by a philosopher one understands a metaphysician, a thinker who wants to solve the riddle of existence, to explain to us ``what holds the world together in its inmost being''. For that is precisely what Democritus wanted, and believed he had done. He by no means presents his atomic theory merely as a method or as a hypothesis, but as an exhaustive \opage{95}explanation of the world. He believes he has found the real in contrast to the sensuous, and the truth in contrast to mere opinion. He is a metaphysician, a systematizer as much as anyone. And his system has a special interest in that it is the first consistent materialistic system ever set up. In the older natural philosophers animistic ideas constantly played a part; these Democritus has separated out, and he has thereby attained a greater purity of execution than his predecessors. If we had had his writings, whose loss is one of the greatest losses of ancient literature, it would have been a great intellectual pleasure to see how his train of thought was worked out in detail; he is, after all, highly praised as a writer in antiquity and is even placed on a level with Plato in this respect. As regards the train of thought too, there is --- precisely by reason of the great opposition --- reason to put him together with Plato: they had worked out their systems at about the same time, the one in a decidedly idealistic, the other in an equally decidedly materialistic direction. For Plato the idea, thought, for Democritus the atom (and space) is the absolute. They had each seized one side of existence and built their world-view on it. Two of the great possible principal systems thereby appeared for the first time in European history. --- \textsc{Windelband} has therefore rightly placed Democritus parallel to Plato in his \textit{Geschichte der alten Philosophie} (p. 206), as both the analogy and the chronology justify. He thereby gets a more correct place than in earlier accounts, in which he appears among the ``pre-Socratic'' philosophers, as a rule as a superseded standpoint in relation to the line Socrates---Plato---Aristotle.

That Democritus was a metaphysician does not, now, exclude his having also been a physicist. On the contrary. All physicists, right up to the most recent times, have been both. The physical explanations were regarded as definitive. Beyond that, people arranged matters in various ways. Either one took matter (the atoms) for that which lay at the ground of everything, and so paid homage to materialism, or one assumed that God had created the atoms. In those who held this latter view (e.g. Gassendi and Newton) there can then be shown a conflict between a mechanical and \opage{96}a teleological conception of nature similar to the one the author mocks in Plato. There is in this respect only a difference of degree (to be sure an extraordinarily significant difference of degree) between Plato's ``Timaeus'' and Newton's ``Principia''.

Already in antiquity analogies were found between Plato and Democritus, without any historical connection between them being assumed. --- \textsc{Aristotle} puts their atomic theory side by side (Phys. I p. 525 b). --- And in the theory of knowledge a likeness was found between them in that they both took only what is accessible to thought (τὰ νοητά) as true. \textsc{Sextus Empiricus}, in whom this remark is found (adv. \textit{Mathem.} VIII 6), adds that this was a natural consequence of the fact that no sense qualities could hold for Democritus' atoms, and that for Plato only the unchangeable, not the changeable that the senses show us, could be the true. This remark points to the fact that the two philosophers, so different in their results, were agreed in this, that in the end it is thinking that must decide what is true and what is not. They are both rationalists. It is another matter whether it is as consistent for the materialist as for the idealist to be a rationalist. --- In the field of ethics too an analogy was found between them in that they both place blessedness in ``the soul'', and not in sensuous enjoyment. Both Democritus and Plato maintain the intellectual joys as the highest, in contrast to the sensuous joys, which are built on unstable ground (\textsc{Stobaeus}, after Didymus Arius, see \textsc{Diels}: \textit{Fragmente der Vorsokratiker.} 55 A 167). ---

In modern times one has not stopped at mere comparison, but has assumed not only an influence of Democritus on Plato, but even an express reference to Democritus on Plato's part. It is an entirely incorrect assertion, made in the dissertation, that no one hitherto should have supposed that Plato had let himself be taught by Democritus.

After Hirzel (in his work on Cicero's philosophical writings) had first sought to show that Plato's ethics is influenced by Democritus', \textsc{Natorp} has taken up this question for detailed investigation \textit{(Die Ethica des Demokritos.} Marburg 1893. p. 155---175. Cf. already his Forschungen zur Geschichte des \opage{97}Erkenntnisproblemes im Altertum. 1884. p. 200 ff. and his article in \textit{Archiv für Geschichte der Philosophie.} III p. 521 fl.). Two passages are decisive. In the \textit{Philebus} (p. 44 B) mention is made of some who have a great name in physics (μάλα δεινοὶ λεγόμενοι τὰ περὶ φύσιν), and who have combated the hedonists, and who thus stand here as Plato's allies. In the \textit{Republic} (IX p. 583 B ff.) Plato, as he seeks to establish that there are joys that are more than the relief of a physical urge, has Socrates say that he has heard from one of the wise that only the joy of thought is pure. In these two passages, which concern the same subject, hardly anyone other than Democritus can be meant. Natorp further finds references to Democritus in the \textit{Protagoras, Gorgias and Phaedo,} in that, besides the fundamental thoughts, he also points to a series of expressions that recall Democritus' ethical fragments, e.g. γαλήνη (calm) and analogous expressions for the peace of the soul. The famous passage in the ``Republic'' (496 B) about the one who is born with a great soul in a small state Natorp also interprets as referring to Democritus. Whether he is right in everything must be left to closer investigation. But the Philebus passage one will hardly be able to get away from. ---

As regards the \textit{Timaeus} itself, too, an influence of Democritus was found long ago. Already \textsc{Stallbaum} says \textit{(Prolegomena ad Timæum} p. 37) that Plato has made use, among other things, of atomisticorum decreta, but complains that the relation between the various sources is not clear. \textsc{Windelband} says in his \textit{Geschichte der alten Philosophie} (p. 234), in discussing the ``Timaeus'': ``Es ist gut, auch hier die Parallele zu Demokrit im Auge zu behalten'', and he expresses himself still more definitely in the book on Plato that he has written for ``Frommanns philosophische Klassiker'', when he says in connection with the ``Timaeus'': ``Hier sind die Lehren der jüngeren Physiologen, besonders auch diejenigen Demokrit's sehr sorgfältig benutzt.'' And \textsc{Natorp} finds it (in the works mentioned as well as in \textit{Platons Ideenlehre} 1903) ``beyond doubt and evident'' that Plato took account of Democritus' physics.

% --- p. 98 ---
\setcounter{footnote}{0}%
\opage{98}A point that in my eyes is very interesting, and that has not come into its own in the dissertation, is the influence that Democritus' physics came to have on Plato's theory of knowledge in the ``Timaeus''. Moreover, Plato, without being aware of it, has come to criticize the atomic theory in the dogmatic form in which it was set up by Democritus and, with some modification, adopted by Plato himself. And it is a criticism that would correspond entirely to the one that would be made nowadays from the standpoint of the theory of knowledge.

Democritus has two kinds of knowledge (Diels 55 B 11): the obscure (γνώμη σκοτίη), to which all the senses (sight, hearing etc.) belong, and the genuine (γνώμη γνησίη), which is entirely separated (ἀποκεκριμένη) from the former, and which, by virtue of its finer organ of thought (ὄργανον τοῦ νῶσαι λεπτότερον), can conduct a finer investigation where the former must stop. It is with genuine knowledge that we discover the qualityless atoms, while obscure knowledge stops short at the sense qualities (color, sound etc.), which have only apparent (νόμῳ), not real (ἐτεῇ) validity (cf. Diels 55 B 9 and 125). In reality only the atoms and the void exist\footnote[1]{Democritus' theory of knowledge is very naive and unclear. Genuine knowledge is supposed to be wholly different from obscure knowledge, and yet the difference is said to be only relative; the former can perceive more finely! Evidently Democritus imagined that if our senses were somewhat sharper, we would be able to sense the atoms. It would be interesting to know what they would then look like. (Even in Dalton, the founder of the modern chemical atomic theory, there occur, moreover, statements to the effect that with more perfect microscopes we would be able to perceive the atoms.) --- In agreement with this, the view is ascribed to Democritus that thinking and sensation are the same, that they spring from one and the same faculty (Diels 55 A 105). --- Plato has seen more sharply here than Democritus.}.

Plato in his other dialogues also has two kinds of knowledge: sensation, which apprehends the qualities in their diversity and in their change, and thinking, which from the sensuous and changing qualities ascends to the general concepts (the ideas). What again and again awakens his admiration and enthusiasm is that it is possible to ascend from the world of the senses to the world of the ideas. His whole philosophy is such an ascent. Mathematics, \opage{99}especially geometry, to which he attached great weight, is a station on the way.

But now --- in Timaeus p. 52 A and B (not yet 27 D!)
--- he is compelled by the adoption of an atomic theory to set up a third kind of knowledge. For the atoms can neither be ``sensed'' nor ``thought'': they have, after all, no qualities, and yet they are single, determinate! Plato now distinguishes, then, between three different objects of knowledge and corresponding modes of knowing. ``One is that which always keeps in the same state, ungenerated and imperishable, never receiving anything into itself from elsewhere, nor itself passing into anything else, invisible and altogether insensible (ἀναίσθητον); it is that which it is the task of thinking (νόησις) to investigate. Another is that which bears its name and resembles it, and which is sensible, in constant motion, appearing in some place and vanishing from it again; it is that which can be grasped by representation together with sensation (δόξῃ μετ' αἰσθήσεως περιληπτόν). A third is everlasting space, which is not subject to destruction, and which provides a place for everything that has a becoming; it must be grasped without the use of sensation, by means of a kind of spurious thinking (μετ' ἀναισθησίας ἁπτὸν λογισμῷ τινὶ νόθῳ), but is hard to believe in.''

What Democritus in his naive materialism found no difficulty in, and what most natural scientists right up to the most recent times have not found any problem in either, is where Plato stops short. Hardly has he accepted the thought of atoms before he asks: of what nature, really, is the knowledge we have of the atoms? He is clear that atoms cannot be sensed. And we do not think them --- he holds --- in the same way as when, by comparing mutually similar phenomena, we form a general concept. By ``spurious thinking'' Plato means (as is seen from the end of 52 B) inference by analogy. That third, he says, is the same as what we mean when we assume that everything is in a definite place and fills a space. It is clear that what Plato calls spurious thinking corresponds precisely to Democritus' genuine thinking, just as his sensation corresponds to Democritus' obscure thinking.

% --- p. 100 ---
\setcounter{footnote}{0}%
\opage{100}It is a stroke of genius on Plato's part to have seen that the atomic theory rests on an analogy. Nowadays one would say roughly this: We conceive the material phenomena as if they consisted merely in motions of small particles in space. In that way we can find definite laws for their changes. But then we disregard all other properties than mobility and divisibility, thus all qualities in the narrower sense. We regard all changes in color, sound etc. by analogy with motions in space. It is a model that we form. We cannot, of course, avoid imagining the atoms with certain colors etc., but we take no account of that.

Now on closer reflection it will appear that Plato need not have changed his theory of knowledge for the sake of the atoms. For the concept of motion and of movable particles is, after all, a general concept (an idea)! When Plato himself speaks of ``the idea of likeness'', he demands that we disregard all other properties of things than precisely this, that they resemble one another. Such an abstraction is required by all general concepts, and it is required by the concept of the atom too. That concept, too, is formed by thinking on the basis of sensation. But because of its clarity we use the analogy with it in all the fields of sensation.

\whiteline

I will only make a few further remarks on the significance of Plato and Democritus for the science of modern times. Democritus' idea of atoms was taken up in modern physics and chemistry because it was suited to serve as a basis for exact conception and calculation. After Bruno and Basso had taken it up in the period of the Renaissance, it was especially Gassendi who brought it en vogue. From Gassendi it came to Newton, and from Newton to Dalton. But it is a question whether Plato's fundamental thought, that definite logical and mathematical relations prevail in existence, has not been fully as fruitful. Both Galilei and Kepler were Platonists. It was in the confidence that \opage{101}definite mathematical relations must hold for the orbits and motions of the heavenly bodies that Kepler carried through the great work that led to the establishment of the laws that bear his name. ---

The world of thought is great, and opposing forces can find work within it, indeed, are necessary for the work to go forward.

% --- p. 102 ---
\setcounter{footnote}{0}%
\opage{102}\essay[Epictetus and Christianity (1912)]{EPICTETUS AND CHRISTIANITY}
\streg

\undertitel{ADOLF BONHÖFFER: \textit{Epiktet und das Neue Testament.} Giessen 1911,\\ Töpelmann. XII + 412 S.}

T\textsc{he} author, who is known through several fine works
on Epictetus, especially ``Die Ethik des Stoikers Epiktet'' (1894),
continues in the present work his study of the Stoic intellectual
movement and its significance through a many-sided and penetrating
comparison between it and the New Testament view of life.
The same qualities that made his earlier works so attractive,
the great thoroughness, the open and free outlook and the great
enthusiasm for the noble figure around whom his studies have so
long gathered, will be found again in the present book as well.
And they stand out here all the more strongly as he has here
arrived at an investigation of so great an interest for intellectual
history as that which a juxtaposition of Stoicism and the New
Testament conception of life presents.
--- The author rightly distinguishes between what belongs to the New
Testament and what is Christian; the latter is a more comprehensive
concept. Underlying this distinction in the author there are not
merely purely historical reasons, but one also notices a certain
influence of Kierkegaard, who precisely in the author's homeland,
Würtemberg, has exercised a quite special influence.

The author begins by criticizing the attempts that have been made
from the theological side to demonstrate Epictetus's dependence on
the New Testament, either in such a way that he is supposed to have
taken up thoughts and impulses from it (Zahn), or that he is supposed
to place himself in conscious opposition to Christianity (Kuiper).
The author shows that
% --- p. 103 ---
\setcounter{footnote}{0}%
\opage{103}the thoughts that are Epictetus's fundamental thoughts had
already come forward before the New Testament writings were
composed, and that Epictetus himself, who alludes to the Christians
in only a single place, had no real knowledge of Christianity and
its source writings. Epictetus has his root in Stoicism, as it had
long since been founded and had, in the first century A.D., been
renewed, especially as a practical direction of life. His own
thoughts he develops from few but fruitful principles, in the
application of which a rich experience of life, psychological depth,
a lively mode of expression and, first and foremost, a mighty
personality come to light. Neither adherence nor opposition nor
rivalry with respect to quite different directions of life and
thought are here of decisive importance.

While from the theological side one has wished to find New Testament
influence, of a positive or a negative kind, in Epictetus, from
another side it has been thought that Stoic elements could be found
in New Testament writings. The interest here attaches especially to
Paul, of whom it has often been supposed that he must have had a
certain Greek cultivation. In the genuine Pauline epistles the author
finds, as far as vocabulary and style are concerned, nothing that
necessarily presupposes Stoic influence. Paul draws, like Epictetus
later, on the Hellenistic common language (κοινή), though in such a
way that Paul's Greek is also influenced by the Septuagint. When Paul
and Epictetus come to related subjects, similarities in style and
manner of argument can then quite naturally arise. There is, however,
no reason to assume that Paul had a thorough knowledge of Greek
popular philosophy. --- A special interest is presented by the famous
passage in the Epistle to the Romans (II, 14 ff.), where Paul says
that the Gentiles do by nature (φύσει) the works of the law. The
author thinks that this utterance comes forward in that particular
context in order to shame the Jews; it is a sacrifice that Paul, in
spite of his religious theory, makes to sound human understanding.
But the sentence is not Stoic; for according to Stoic ethics only
insight (λόγος), not natural instinct (φύσις), can lead to truly good
action. ---

Concerning the speech that the ``Acts of the Apostles'' (XVII)
attributes to Paul in Athens, the author agrees with
% --- p. 104 ---
\setcounter{footnote}{0}%
\opage{104}Wilamowitz-Moellendorff's remark that with whoever can assume
that this speech is by Paul, ``mit dem ist nicht zu reden'', and he
adds: ``Im übrigen finde ich in der Rede nichts spezifisch Hellenisches
oder Stoisches, was den Philosophen sonderlich imponieren
konnte''. --- As regards the concept of the Logos in the Prologue of
the Gospel of John, the author strongly emphasizes the difference
between this concept and the Stoic concept of the Logos, and joins
the view that maintains the Jewish origin of the Johannine concept of
the Logos and requires λόγος to be translated by ``Word'', not by
``Thought'' or ``Reason'', although he concedes that Hellenic thoughts
have had an influence on Philo and through him on the Prologue of the
Gospel of John. ---

In the book's second main section the relation between Epictetus
and the New Testament is investigated partly through a consideration
of the vocabulary, partly through a systematic comparison of the two
worlds of thought.

The section on the vocabulary, and a section attached to it on
particular expressions and thoughts, contain many interesting
contributions to clarity about Stoic terminology and good examples of
shifts of meaning, just as they also show how the two different
directions of thought determined a difference in vocabulary. But it
is also shown how on several points an analogy nonetheless holds
between Epictetus's ideas and those of the New Testament. ---

The most interesting section of the book is the last, which gives
a systematic valuation of the two conceptions of life. The author's
great learning here works together with his open sense and his
personal experience of life.

According to the author there is a certain spiritual kinship between
the two conceptions of life. What the Christian possesses in his
faith that God can be moved by prayer and penance to save his soul,
the Stoic has in his conviction that man, through the working of the
deity, has once and for all been so constituted and endowed that he
has the source of strength and happiness within himself, a source
that can ever anew well forth, even if it seems to have dried up.
--- But Stoicism falls short of Christianity in being philosophy, not
religion; therefore its influence
% --- p. 105 ---
\setcounter{footnote}{0}%
\opage{105}was limited to a small circle, and it did not have the
gripping and captivating power that a positive religion can exercise.
--- On the other hand it stands in a more positive relation to human
culture, while the New Testament religion, in which the expectation
of the kingdom beyond was all-dominating, took an ascetic attitude
toward everything earthly. It is eschatology that here sets primitive
Christianity apart from Stoicism. --- The Stoic concept of God was
indeed more abstract than the Christian one, but in return it did not
present the difficulty of how an absolute and unconditioned being can
have ethical attributes when it has no inner and outer resistance to
overcome, and thus no work to perform.
--- Finally, Stoicism has the merit that, even before Christianity, it
maintained the unity and belonging-together of the human race.


An objection against both conceptions of life is that they lack the
idea of a great, progressive improvement of human conditions and do
not see the great importance that outer reforms can have for
spiritual development. Stoicism worked for the abolition of slavery
as little as Christianity did, although it, like Christianity, taught
that all men were free in their inner selves. In neither of the two
conceptions of life was the criticism of slavery grounded in the idea
of social justice. ---

In the author's conviction the influence of Stoicism is not over
even today. Anyone can still experience the strengthening and
consoling power of Epictetan idealism. If Stoicism fell in the
ancient world, it was because there was no longer any soul brave and
proud enough to feel attracted by the stern and lofty greatness of
its teaching! ---

I believe that I have herewith rendered the train of thought of
this interesting book. I permit myself only to mention further that I
myself at one time (see Mindre Arbejder I p. 158---179) expressed a
similar view of Stoicism and its relation to Christianity. But I could
not make use of so comprehensive a learning as that which is at the
author's disposal, and I had besides already gained much instruction
from the author's
% --- p. 106 ---
\setcounter{footnote}{0}%
\opage{106}earlier studies of Epictetus. When the author therefore refers
to my treatise, he refers to a work that stands in debt to himself.
--- The author ends his book by joining in the remark in my Ethics
that our ethics is Greek ethics, deepened and enriched by all later
development in the various spiritual and social spheres.

But whatever view one may oneself hold concerning these things, one
will gain rich instruction in Bonhöffer's book about one of the most
important episodes in intellectual history.

% --- p. 107 ---
\setcounter{footnote}{0}%
\opage{107}\essay[Marcus Aurelius (1912)]{MARCUS AURELIUS\footnote[1]{A lecture in the ``Græsk Selskab'', 22 Oct. 1912.}.}
\streg

G\textsc{reece} succumbed to the Roman legions, but in return
conquered Rome through its culture, its art, its science, its
philosophy. The Greek culture that thus came to rule the world was,
however, not the one that flourished when Greek political life was
free and vigorous. It was a culture that had burst the national
barriers and led to the development of a general consciousness of
humanity, at the same time as it stressed the inner, peculiar life of
the individual personality. Both by its universal and by its
individualistic character it stood in contrast to classical Greek
culture. And while Greek culture proper succeeded, in a way and to a
degree that perhaps no culture has ever managed, in gaining
expression in great works, rich in content and perfect in form, later
antiquity was not able to gain such an expression. In particular the
two tendencies, the universal and the individual, often came into
sharp opposition to each other, insofar as they did not go their
separate ways. There came to be lacking the harmony, the union of
fullness and strength, that was the distinctive mark of classical
Greek culture. In the philosophical sphere it is the Stoic school
that takes the lead, and it is characteristic that the leading Stoic
thinkers came from wholly or half barbarian lands. Of the school's
founder
% --- p. 108 ---
\setcounter{footnote}{0}%
\opage{108}Zeno it was said in antiquity (by the Academic Polemo) that he
had clothed Plato's Ideas in Phoenician dress; Zeno was from Cyprus,
colonized by the Phoenicians. Later on, in the imperial age, the Stoic
school was recruited not only from the barbarian lands of the East
but also from those of the West. Epictetus, the most remarkable and
the profoundest of the Stoics of the imperial age\footnote[1]{Cf., on the Stoics of the imperial age, especially with regard to their relation to Christianity, my \emph{Mindre Arbejder}. Første Række p. 158---179 (``Hedenske Sandhedssøgere'').}, was from Phrygia; but Seneca
and Marcus Aurelius were of Spanish stock. Yet Marcus Aurelius wrote
his diaries, which after his death were published under the title
``Self-Contemplations'' (εἰς ἑαυτόν), in Greek, and he belongs to
Greek literature through language and thoughts, through his
intellectual content, even if the significance that this intellectual
content has in his personal life is very different from what we would
find in a Greek of the classical period. He feels himself a citizen
of the world, the word ``world'' taken both in the sense of the world
of men and of the great world of nature; but his notes concern
exclusively what stirs in his mind in the course of life, in his work
and his doubt, in his antagonism to his surroundings, or at the
thought of his connection with the whole of nature. The thoughts he
has appropriated get their essential significance for him as a help
to keep himself up, to be himself and to assert his personality in
spite of outer opposition and inner trials. His notes are one of the
first self-confessions that literature has to show. We are led into
the inner world of a personality that cannot be counted among the
great, but who with honesty, strength and inwardness reflected on the
conditions of life as they took shape for him. Not as a scientific
thinker but as a philosopher of life is he of great interest. And in
spite of all that the ages have brought, we can still see in him a
fellow combatant and learn from the thoughts that he wrote down in his
diary in the privy chamber of the imperial palace or in his tent in
the camp up by the Danube.

% --- p. 109 ---
\setcounter{footnote}{0}%
\opage{109}Marcus Aurelius\footnote[1]{His first name was Catilius Severus, but after being adopted by Antoninus Pius he received the name Marcus Aurelius (\emph{Capitolinus}: Antoninus Philosophus. I; V; VII.} was born in Rome in the year 121 A.D.
His family came, as already remarked, from Spain; under Vespasian it
had been admitted to the patrician order. As a boy he attracted the
attention of the Emperor Hadrian by his gentle, serious and pure
disposition, and when the emperor adopted Antoninus Pius, it was on
the condition that the latter in turn adopt the young Marcus. It was
with his adoptive father as his model that he prepared himself to
govern the Roman Empire. With great reverence and with a gratitude
that only significant personalities can feel, he has described not
only what he owes to his imperial model and his other relatives, but
also what he owes to his teachers, the grammarians, rhetoricians,
jurists and philosophers who directed his studies. He thanks
especially the Stoic Rusticus, because he led him to serious
absorption in philosophy and kept him away from sophistry, rhetoric
and belles-lettres, and instead laid the weight on what could serve
the development of character. And not least he thanks him because he
made him acquainted with the writings of Epictetus. For Marcus set
Epictetus highest of all philosophers; the philosophical emperor is
the disciple of the philosophical slave. The Peripatetic Severus,
with whom he entered into brotherhood, he thanks for the idea of a
free republic governed according to the principles of equality and
justice, and of a monarchy that honors the freedom of all. One sees
from this how the philosopher himself conceived his position as
emperor in Rome. The gods he thanks for his not having come to follow
the unfortunate inclinations that might lie in his nature, and for
the good relatives they have granted him, especially for his good wife
Faustina. This last thanks seems, according to what is known, to be
less well founded\footnote[2]{After Faustina's death he asked the Senate to pay her honors and to build her a temple, praising her, although she --- whether he did not know it or pretended not to know it --- was notorious for her debaucheries. (\emph{Capitolinus}. XXVI).}.

% --- p. 110 ---
\setcounter{footnote}{0}%
\opage{110}At forty years of age Marcus was drawn out of the quiet, inward-turned
life he had led until then. Antoninus Pius on his
deathbed named him his successor and had the image of the goddess of
victory carried from his own bedchamber into that of the new ruler.


He does not belong among Rome's great emperors. But his reign is
marked, in internal affairs, by a series of humane measures, such as
the establishment of a chief office of guardianship and the laying
down of the rule that every doubt should benefit the accused. During
the great calamities --- plague, earthquake, war, --- that afflicted
the empire, he brought relief to the best of his ability. And although
he is said not to have had any special talent as a general, he
nevertheless acted with strength, prudence and consistency against the
Germans storming forward from the north, and it was only a rebellion
in one of the provinces that prevented him from protecting the borders
of the empire in a lasting way. For almost fourteen years he, whom
Mommsen has called ``the warrior prince against his will'', lay at war
up in the border lands against Marcomanni and Sarmatians. In the camp
he wrote down some of his most interesting aphorisms, and it was in
the camp at Vindobona, the present Vienna, that he died in the year
180, 60 years old. ---

Thought, philosophy, has for Marcus Aurelius the significance of
being a stronghold, something that stands firm while everything else
is vanishing, fleetingly streaming away. The things in the world are
in constant motion; to attach one's heart to any of them would be
like falling in love with a bird flying past, which is at once out of
sight. And so it goes with the whole of life; the time between the
first breath drawn and the last breath given out is at bottom as
short and fleeting as the time each breath fills. But reflection helps
us to a firm standing amid all this change. In it our genius (daimon)
expresses itself, the god within us. Through it we can assert
ourselves against all toil and all allurement, and receive what meets
us with the consciousness that it stems from the same source from
which we ourselves have our origin, the great, all-embracing nature
whose laws determine everything, and which reigns not only in the
outer world but in our own inner selves. ``There are many who seek out
solitary places,
% --- p. 111 ---
\setcounter{footnote}{0}%
\opage{111}country regions, the shores of the sea or mountains; and you
yourself, too, long for such things. But it is at bottom foolish,
since you can, at whatever moment you will, withdraw into yourself!
No quieter and more undisturbed refuge can a human being find,
especially when he can sink himself, within, into thoughts that can
restore the harmony of his being, so that he can again meet the
wickedness of men (which is really folly), meet the changing fates and
the bodily discomforts, and master the urge for outward honor! Look
within! For there is the source of the good, a source that can ever
well forth anew, if only you dig! You shall hold yourself in awe and
not let your happiness depend on other souls than your own. Let the
god within you come to rule over a manly, venerable human being, a
citizen, a Roman! This god or daimon, whom Zeus has given to each one
as overseer and guide, a spark of Zeus' own being, is nothing other
than each one's own mind and thought\footnote[1]{The Stoics took over the idea of daimons from popular belief, where it denotes a kind of momentary gods, by whose intervention sudden turns and inspirations are explained. The transition to the Stoic idea is formed by the belief that each human being has his daimon. See \emph{Usener}: \textit{Götternamen} p. 290---297. --- In Marcus Aurelius the daimon is made one with man's mind or thought (see II, 6. 13, 16; III, 4---6; V, 27; XII, 1. 3. 26), as E. \emph{Rohde} remarks \textit{(Psyche.}\textsuperscript{2} II, p. 317). Moreover, Heraclitus had already said that a man's temper (ēthos) is his daimon. (\emph{Diels}: Vorsokratiker p. 82).}.''

Marcus Aurelius seldom mentions his work as a statesman and his
imperial dignity. He warns himself against chasing impossible ideals.
Do not expect Plato's state, but be glad of a little progress! --- He
is ironic about his own warlike career: A spider is proud when it has
caught a fly; there are others who rejoice at having shot hares or
bears, or at having caught sardines in a net; others again at having
conquered Sarmatians; --- but are they not all robbers, when one looks
closely? --- and he warns himself against taking on the color of his
position: See to it that you are not Caesarified; it happens so
easily, after all. But keep yourself such as accords
% --- p. 112 ---
\setcounter{footnote}{0}%
\opage{112}with your striving for wisdom, and live in all respects as
befits the disciple of Antoninus! ---

Court life has not pleased him. ``If you had at one and the same
time a stepmother and a mother, you would doubtless honor them both,
but you would always return to your mother. So it goes with you with
the court and philosophy. Return often to your spiritual mother, and
rest there; then court life too will become easier for you to bear,
and you yourself will become easier to bear there!'' --- He feels
himself constantly an object of slander, criticism and backbiting.
Then he consoles himself with the words of Antisthenes, that it is a
king's lot to do good and to be spoken ill of for it. He knows very
well that there are those who long for this schoolmaster, who has
indeed done them no harm but whose conduct was a silent accusation
against them, to depart by death. And when even those who stand close
to him, and for whose sake he has endured so much struggle, nourished
so many wishes and had so many cares, will feel relief at his passing,
then it must be easy for him himself to leave a life where things go
on in this way. Yet without bitterness. That experience shall not be
allowed to change his character! He knows, besides, that even if he
bursts with vexation over them, they will go on acting as they do. And
he neither esteems nor fears their judgment. Just try once, he says to
himself, to think your way into what stirs in these people. They
flatter one another and yet despise one another! They seek to raise
themselves above others and yet bow humbly to the very ones they wish
to surpass! --- There must, after all, be shameless and faithless
people in the world. It is your fault when you are surprised by it.
And when you are disappointed at meeting ingratitude, should that not
be because you expected another reward than the one that lies in
acting according to your own inmost nature?

Even those who wrong us we can love, when we consider that they are
akin to us in their being, --- that they act from ignorance, --- that
they, as well as we ourselves, will be dead in a short time, --- and
above all, that no harm has been done to us, for our mind has not
become worse than it was before. The jewel remains a jewel, whatever
one may say or
% --- p. 113 ---
\setcounter{footnote}{0}%
\opage{113}do. Let them only throw dirt into the pure and clear spring:
its vigorous welling will wash it away, and its water will not take
color from it! In spite of everything you can keep your thought pure,
understanding, sober-minded and just. You will have in you an
eternally flowing spring if you always assert your freedom in league
with a kindly, simple and modest disposition. ---

Marcus Aurelius's contemplations do not form a coherent book. His
thoughts and moods are expressed in short and often broken sentences,
occasioned by experiences that as a rule are not even hinted at, but
which one must infer. What kind of people, e.g., it was who called
forth the utterances last quoted, we are not told. This constant lack
of motivation is explained by the fact that the initial ideas or the
initial moods filled his mind in a living way, and that it was only the
thoughts they occasioned that he needed to write down. Therefore one
cannot read long at a time in the book, and in the first readings
there is much whose significance is not understood. But then it may
happen that, by opening the book at random and fixing on particular
sentences, one calls up the whole corresponding mood and with it the
understanding. It goes here in a similar way as when one reads a
collection of poems. Each single little piece demands its mood and
must be read by itself.

It is not always merely moods, however, that must find vent. The
imperial philosopher moralizes to himself, scolds himself, or urges
himself to have his thoughts ready for situations to which he would
otherwise easily succumb. One can see, e.g., that it has often been
hard for him to begin his day with zeal and confidence. At times he
has difficulty waking up fully and wishes to prolong his stay in the
arms of sleep on the soft and warm bed. Therefore he says to himself:
``In the morning, when you have difficulty getting up, you shall have
the thought ready at hand that you are awakened to do the work of a
human being; it is for this that you have come into the world, not to
lie warm between soft blankets. Your work is to act for the good of
society, and how can you hold this work in low esteem?'' At other
times he resolves,
% --- p. 114 ---
\setcounter{footnote}{0}%
\opage{114}first thing in the morning, to be prepared to meet ungrateful,
brutal, treacherous and envious people; what matters is to hold fast
to the knowledge that they act from ignorance, and that in spite of
everything they are of the same kin as he himself, destined like him
for a great, common work.

One sees here how he uses his philosophical ideas to strengthen his
zeal and to make it possible ever to take up the struggle of life
anew. These ideas he owes to Stoic philosophy, as it had taken shape
in the Stoics of the imperial age, especially in Epictetus. We gain an
insight into how this philosophy is used in the spiritual economy of
a serious human being living amid great affairs, as a defense against
everything that can dull, ensnare and disquiet.

The Stoics laid great weight on the great, lawful connection in all
things, on the thought of nature as a great whole, to whose law
everything individual is subject. This thought also forms the
background of Marcus Aurelius's view of life, although at times he
seems to move away from it. His thoughts oscillate, use different
points of view, but always return to that central point of view.

In order to free himself from what might ensnare and beguile, he will
see things in their nakedness and dissolve them into their elements,
separate them out from the whole connection. He analyzes at times in
a crude and cynical way that which is an object of enjoyment or of
pride. When we thus strip things bare, we rise more easily above
them. In contrast to the crude and the material there then appears
that loftiness of mind which is grounded in the consciousness of being
a citizen in the great world community, the highest state, in relation
to which all other communities are what the houses are in relation to
the city.

Above all there are two apparently opposed thoughts that come back
again and again. Everything changes --- and yet nothing new happens!
Existence is like a river in constant flow. Without interruption the
passage takes place from what stood as future, through present, to
past. Where now are the famous and mighty men of the past? A court
like the one you yourself have was had by Alexander, Caesar, Hadrian.
They are long since gone, after having filled a brief
% --- p. 115 ---
\setcounter{footnote}{0}%
\opage{115}moment in infinite time. But in spite of this change you need
not expect or fear anything new. If you have looked about you in the
world and gained an overview of what happens, you can be at ease
about the future. Everything will repeat itself. The same rhythm in
the course of the world will ever anew make itself felt. The events of
the future will in all essentials resemble those of the past. If you
have only reached your fortieth year, you can, at everything that
happens, say with good heart that it is something you have often
seen; it is the old that comes again in new forms; it is the same
drama that is played, only the actors are
different\footnote[1]{At times the repetition causes disgust and weariness of life. ``Just as you have grown tired of the games in the amphitheater and similar places, because one always sees the same thing there, and the monotony leads to disgust at the spectacle, so you will find it the whole of life through. From first to last things are the same and spring from the same causes. How long, then, is it to last!''}.

But here too the feeling of unity with nature again and again gains
the upper hand. The very repetition and the rhythm bear witness, after
all, to the great world law that reigns through all vicissitudes. All
things cooperate in maintaining the world order. Everything forms one
world; one god works through everything, one being, one law, one
thought common to all things and all men. His mood, which took on a
bitter character from seeing things in their nakedness, their change
and their repetition, alters the more he immerses himself in thoughts
of this kind. He no longer needs the cynical contemplations, but lays
his fate securely and confidently in the hands of the great order of
nature. ``Everything that is harmonious for you, O Nature (Kosmos), is
also in harmony for me. Nothing that is timely for you comes too early
or too late for me. Everything your seasons bring is fruit for me.
From you comes everything, in you everything subsists, to you
everything returns!''

This feeling of unity, which in the face of the course of nature
becomes mild, wistful resignation, and in the face of men heartfelt
sympathy, marks what is central in Marcus Aurelius's view of life. He
is not concerned to fix wholly definite ideas about the ultimate
grounds of things. He seeks only to
% --- p. 116 ---
\setcounter{footnote}{0}%
\opage{116}show that, however that may be, man is master of himself and
has the task of asserting himself according to his own peculiar
nature, a task that can always be accomplished. If there is a
providence that determines my fate as well, then I am well provided
for; if providence concerns only the general order of things, then it
is of no use, surely, for a single part to find fault with the whole;
if there is no world order at all, but everything is only a play of
atoms, and chance reigns, then see to it that there is at least
connection and meaning in yourself! If there is no one who cares for
me, then I can care for myself. And to my nature belongs thought, and
belongs participation in a community. As Antoninus, Rome is my
fatherland; as a human being, the world is my fatherland! It is
conduct and action in the single moment, in the particular, given
situation, that matter. Everything else is indifferent. --- Yet, as is
seen from the train of thought just cited, Antoninus always comes back
to the conviction that the concentration of will that makes us
independent of everything external is itself an element in a world
order; the genius (daimon) that expresses itself in our willing is an
emanation of the force that works in the great connection of the
world.

Death too is a link in this connection, as natural a link as birth
and growth. It too is something that nature wills. One should be
neither too preoccupied with it nor impatient or presumptuous toward
it, but wait for it as one waits for the birth of a child. Whether
death is merely a dissolution into elements, or whether it is a
transition to another life, or perhaps an absorption into the great
unity of the world, we do not know. You have gone aboard; you have
sailed; you have reached the end of the voyage; so leave the ship. If
you then pass over to another life, you will find gods there too, if
it is true that there are gods everywhere. If it is to an unconscious
state that you pass over, then you are at least set free from bondage
to the material, from the constant opposition between your mind, your
genius, on the one side, and earth and dirt on the other. The only
thing that could make us hold fast to life would be if we lived
together with people who looked on life with the same thoughts as we
% --- p. 117 ---
\setcounter{footnote}{0}%
\opage{117}ourselves. But now you see yourself what torment disharmony in
living together causes, so that you must cry out: Come quickly, O
death, lest it go so far that I forget myself!

And when, now, a given situation becomes unbearable, in that it
demands more than the strength we have to put in?
--- Then leave life by your own decision, not in bitterness and anger,
but in the same way as you go out of a room because there is smoke.
You should not imitate the beast fighters in the arena who, though
badly wounded, beg to be spared until the next day, but neither the
Christians who in blind zeal rush into death\footnote[1]{While Marcus Aurelius (XI, 3) derives the Christians' craving for a martyr's death from ``blind zeal'' (or ``love of contest'', parataxis), Epictetus (IV, 7, 6) derives their fearlessness from mere ``habit'', in contrast to the rational reflection that according to him is the condition of true courage. By ``habit'' Epictetus is presumably thinking of its having become a tradition in the Church that the martyr's death was to be coveted.}. It must be your own reflection that
prepares your departure, and this departure should take place, not
theatrically, but thoughtfully and solemnly.

Even if in one way or another you are called away from the stage
before you have played in all five acts, but only in three, your life
can nonetheless form a whole; it is not you, after all, who composed
the whole drama. Go away, therefore, with a kindly mind. The power that
leads you out of life is, after all, kindly minded too. --- These are
the last words of the book. And the gentle tone that expresses its
author's inmost mood sounds in them as well. ---

I have sought to render the thoughts that one of the noblest human
beings of the ancient world formed about life and the world. Free from
the Oriental mysticism that spread more and more over the lands, he
strove to look clearly and soberly at things and yet keep his mind
open to everything that makes life noble and high. And there reigns
in his contemplations an inwardness and a richness of mood that hardly
have their counterpart within ancient literature. It is not weariness
or inner division that brings him to think about life as he does, and
to find rest in the thoughts he forms. From youth on there was, in
% --- p. 118 ---
\setcounter{footnote}{0}%
\opage{118}the words of one of his first historians, ``gravity without
sadness''\footnote[1]{Capitolinus IV: sine tristitia gravis.}
characteristic of him. It is through this that he has felt the urge
to live himself into the world of thought that one of the last
independent Greek schools of philosophy had built up. He has taken
from it what he could use for spiritual self-preservation under the
circumstances in which his position and the conditions of the time
placed him. Such an individual urge to self-assertion, to
understanding of oneself and fidelity to oneself, had developed in
Greek antiquity from the turning point marked by Socrates. But it has
hardly anywhere found so unreserved an expression as in Marcus
Aurelius. It is hardly even a literary work proper that we have before
us. The emperor surely never imagined that what he wrote down would be
read by others. It is a sanctuary into which we have perhaps by chance
been initiated, --- one of that kind of sanctuaries that endure however
many temples fall. Precisely through the personal character of his
thoughts he belongs to all times.

And yet he stands clearly and definitely on ancient ground. This shows
itself especially in the limitation of his picture of the world. In
spite of all change there is a certain number of possibilities that
will ever be realized anew in rhythmic succession. Experience is for
him closed. And this limitation gains decisive importance for him,
because he builds his trust and security on his knowing what the
course of things could bring. However great the oppositions he has to
struggle with, none of them was so strong and deep as the one that
arises from the possibility that life should come to be led under
wholly new conditions, --- that the rhythm known until then should be
only a slight ripple in the great stream of the ages. With this
possibility all the problems of life are sharpened. What we may miss in
Marcus Aurelius is courage and confidence in the face of the new and
uncertain; they are lacking for the simple reason that he did not
believe in anything really new. Memory had consumed hope and fear, and
the world of memory was not to be enriched by new experiences. Greek
intellectual life
% --- p. 119 ---
\setcounter{footnote}{0}%
\opage{119}is essentially built on memory; on this rests its great
significance as well as its strange limitation.

Marcus Aurelius breathed his last in his camp during the struggle
against the barbarians storming forward, whose appearance in history
was precisely to mark the beginning of a wholly new period in the
history of the human race. Here he stood face to face with new forces
and new possibilities, although he himself probably regarded the
attacks as repetitions of the many border storms that the empire had
had to contend with ever since the time of Marius. This time, however,
the storm was, even if it took three centuries more, to lead to the
downfall of the Roman Empire and of ancient culture. But although the
swell in human intellectual life has since become stronger than the
Greeks imagined, there is such depth and truth in the Greek thoughts
that they have ever anew made themselves felt, and moreover, under the
new conditions, have proved to contain possibilities that did not come
to unfolding, and perhaps could not come to unfolding, in antiquity
itself. And where it is a matter of testing thoughts with regard to
their carrying power in the course of life, Marcus Aurelius's ideas
too, closely bound up as they are with his personality, will ever anew
awaken attention and win understanding.

% --- p. 120 ---
\setcounter{footnote}{0}%
\opage{120}\essay[Rousseau and the Nineteenth Century (1912)]{ROUSSEAU AND THE NINETEENTH \\ CENTURY.}
\streg

R\textsc{ousseau's} influence extends through the whole of the nineteenth
century. In an immediate way he acted upon the people of his own time,
and among them in turn especially upon a series of
eminent minds, through whose works his influence makes
itself felt in large circles that know nothing of him, and
perhaps do not want to know anything of him. We must therefore distinguish
between a mediate and an immediate influence.

Throughout the whole century, however, both his personality
and his thoughts are objects of strong sympathies and antipathies.
He stands continually as a psychological problem through the contradictions
of his personality, and the value of his thoughts is assessed differently
according to the different positions the assessors take toward
the ethical and social problems that Rousseau had posed.
The visitors' book at Les Charmettes bears witness to this. When George
Sand visited Les Charmettes in the year 1861, she found the pages of the visitors' book
``filled with coarse abuse or foolish accusations
against Rousseau'', ``with dirty taunts, with judgments and
hypocritical curses or pedantic reprimands'', --- and
in contrast to this ``tributes from democratic and
socialist workers''.\footnote[1]{\textit{A propos de Charmettes.} (Revue des deux mondes. 15 Nov. 1863.) --- When I visited Les Charmettes in the summer of 1909, I unfortunately neglected to study the visitors' book, so I cannot say whether such opposite testimonies are still given.} The battle over him still goes on. It
% --- p. 121 ---
\setcounter{footnote}{0}%
\opage{121}is now Conservatives and the Orthodox, now Liberals and Freethinkers,
who fall upon him.

That the battle remains joined over his personality and his
work has its reason in their typical character. His position
toward the ``Enlightenment'' of his time and toward the orthodoxy and conservatism
of his time presents analogies with standpoints that arise
in the course of the nineteenth century. One can point to a
series of men who are not all influenced by him, nor
all sympathize with him, but whose stance in religious
and social questions, partly because of their personalities,
partly because of the situations confronting them, comes to
show a kinship with his and therefore serves, in the history of ideas,
to illuminate and confirm its significance.

An inquiry into ``Rousseau and the nineteenth century''
must concern both his mediate and his immediate influence,
and must bring out not only this influence but
also the confirmation his ideas receive from the stance of more or less
kindred persons in analogous situations. To this
must be added a look at his anticipation of certain more special
points of view that have come to be developed in the course of the nineteenth
century.

Rousseau's main problem was conditioned by the division of labor
in the mental as well as in the material sphere, and by
the one-sided development of particular sides of personality that
was thereby brought about, to the detriment of the healthy and integrated unfolding
of life. His chief demand then became: Rendez l’homme un! In order
to satisfy it he sought down to the immediate and involuntary
forms of life, to the life of feeling in contrast to the life of the intellect,
to life in free nature, to the mood and the urge that the
changing fates of life call forth in men's minds. That was
what he meant by ``returning to nature''. It was,
if one keeps to the clarification his first paradoxes
received in his later writings, by no means his intention to deny
culture. He demanded only that culture should grow forth out of
nature, correspond to the stage of development of the individual or the people, and preferably
arise through their self-activity. These were thoughts that pointed
far beyond the eighteenth century.

% --- p. 122 ---
\setcounter{footnote}{0}%
\opage{122}\begin{center}I\end{center}


Most direct, naturally, was the influence on French
literature. \textsc{Chateaubriand} was in his youth intoxicated by Rousseau's
works; both his religious and his poetic tendencies
were born of this influence\footnote[1]{Cf. \emph{Giraud}: \textit{Les origines du Génie du Christianisme} (Compte rendu de l’Académie des sciences morales et politiques. 1911, p. 43).}, which does not become less significant
because other influences came later, or because
the young romantic was himself placed face to face with great nature,
learned in the school of solitude and poverty, and thus had
experiences like those by which Rousseau had been guided. \textsc{Madame
de Staël}, who in her youth wrote an enthusiastic and understanding
little book on Rousseau, was likewise led by his works onto
the romantic path. \textsc{Lamartine} too was a worshipper of Rousseau
in his youth; although his mother burned Rousseau's books when
she found them in his room,\footnote[2]{Cf. Annales de la société J. J. Rousseau. III, p. 271, 283.} their effect was strong
and lasting. Both in his nature poetry and in his religious poetry
after-effects of it are stirring. \textsc{Lamennais} was influenced by
Rousseau's democratic and social enthusiasm; it can be traced both
in his youth, when in the struggle against rationalism he pointed to
the immediate faith accessible to the simplest, and later,
when in the name of religion he demanded the social transformation
of which Rousseau had been the herald. ``Paroles
d’un croyant'' (A Believer's Words) has been called ``the Contrat social presented in
the form of a parable''.\footnote[3]{Charles Adam: \textit{La philosophie en France.} Paris 1894, p. 103.} There was, however, one important point on which Lamennais
(like so many of the nineteenth century's
social reformers after him) corrected Rousseau's position. It was the
weight he laid on the combination of the workers in small and
large associations. He says to the poor: ``Join together,
and be a mutual support and shelter for one another. As long as you
are divided, and each thinks only of himself, you can only
expect suffering, misfortune and oppression.... When you are asked
how many you are, you shall answer: We are only one, for our
brothers, they are ourselves, and we are one with our brothers''.\footnote[4]{\textit{Paroles d’un croyant.} (Paris 1833), Chap. VII.}
% --- p. 123 ---
\setcounter{footnote}{0}%
\opage{123}\textsc{George Sand} remained true all her life to the enthusiasm
that the reading of Rousseau's works had awakened in her. Let
us hear her own words: ``Jean Jacques's language and his manner
of thinking captivated me like a music heard while the sun
was shining brightly.'' She found again something of her own nature in him.
René Doumic calls the first chapter of his book on George
Sand: ``Psychologie d’une fille de Rousseau''. In her later
years she has (in the already mentioned article on Les Charmettes)
expressed in warm and beautiful words what she herself and others,
who perhaps would not admit it, owe him: ``For my
part I remained true to him, true as to the father to whom I
owe my life. For even if he has not left me his genius,
he has, in me as in all artists of my group,
awakened love of nature, enthusiasm for truth, contempt
for an artificial life and distaste for the vanity of the world''. ``The most powerful
geniuses of our time as well as its gentlest talents owe
him, however much they might wish to deny it, their decisive
initiation .... Rousseau's influence will continually spread
further, even to those who have not read him''.

Even a critic who has very little sympathy, and one
may add, very little understanding, for Rousseau must, half
a century after George Sand, confirm her words:
``Partly through himself, partly through the authors who have been influenced
by him, he still acts upon many of us, without our perhaps
knowing it ourselves''.\footnote[1]{Jules Lemaître: \textit{Jean Jacques Rousseau}, p. 352.}

On English literature Rousseau's influence has hardly been
considerable. Yet Southey and Wordsworth are strongly influenced
by him,\footnote[2]{Cf. Annales de la société J. J. Rousseau. III, p. 270.} and in splendid stanzas Byron has glorified the poet of the new
Héloïse.\footnote[3]{\textit{Childe Harold.} III, 76---104.}

On the other hand, his influence on German literature and intellectual life
has been exceedingly great, and it is especially here that the mediate
effect is clearly seen to proceed from the immediate. Through
Herder and Lessing, Kant and Fichte, Goethe and Schiller, the
great founders of German humanism, he has acquired decisive
% --- p. 124 ---
\setcounter{footnote}{0}%
\opage{124}significance for the whole intellectual life of the nineteenth century. The
men named\footnote[1]{Cf., as regards Kant, my work \textit{Den nyere Filosofis Historie.} 2nd ed. II p. 32 and p. 73 f., where I have tried to show that a Rousseauian influence made itself felt in Kant at two different points in time. --- Cf. furthermore \emph{Rich. Fester}: \textit{Rousseau und die deutsche Geschichtsphilosophie.} Stuttgart 1890. --- H. v. \emph{Stein}: \textit{Die Entstehung der neueren Aesthetik}, Stuttgart 1886, p. 266.} are the real creators of the intellectual content
with which so-called Romanticism appeared, and in their development
Rousseau intervened strongly. Here as everywhere it is
immediacy in contrast to the reflection of the intellect, life in
the whole instead of the divided life in the parts, the great understanding
of the simple and of the sublime, that conditions his
influence. It is above all through what he has been for the men named,
and through them for the whole European intellectual development,
that, as George Sand and Lemaître said, his influence
has been able to extend mediately to many who do not
suspect what they owe him.

\begin{center}II\end{center}


Toward the middle of the nineteenth century there arose,
in the sphere of spiritual culture, a situation that in several
respects recalls the conditions under which Rousseau appeared.
The romantic movement, which had had such great
effects in the spheres of art, religion and thought, was
slackening or stiffening. It was the age of the epigones. There was need,
as in Rousseau's time, of a new infusion of ``nature''.
At the same time the view was asserted that all spiritual
needs must and could find their satisfaction in
the forms of systematic thinking. Even the striving and yearning
that lead to the great problems in the sphere of the view of life
were now to be only an obscure and self-contradictory form of
what appears in full clarity in the form of pure thought.
This view reached its height in the Hegelian philosophy.
But from quite another side too it was maintained that personal
life could be wholly absorbed in intellectual forms. The great upswing
of natural science certainly led to a reaction
% --- p. 125 ---
\setcounter{footnote}{0}%
\opage{125}against abstract speculation, but at the same time led not infrequently
to the view that now a human being's personal relation
to life, every urge and every faith, could be determined through and through
by means of the results of natural science. The great question
Rousseau had raised was, seen from an essential side, precisely
whether man's immediate experience of life, won face to face
with the conditions of life themselves and deposited in a deep and original
life of feeling, should not be of some significance, indeed,
be the really decisive factor. What Rousseau meant by
``nature'' was precisely the original, spontaneous and immediate
in man's life of feeling.\footnote[1]{Cf. on Rousseau's concept of nature my book Rousseau og hans Filosofi. Second ed. (1912) p. 97---103.} The value of a culture, an intellectual
culture too, depends on whether it has its root in such a
foundation or has at least struck, or can strike, root in it. From
general scientific points of view man's purely
individual and immediate captivation by life and its course
cannot be derived. Pure deductions are not decisive here when
they do not meet with the results won through individual induction.

The analogous situation in the middle of the nineteenth and in the
middle of the eighteenth century called forth analogous problems
and an analogous attitude in those who seriously felt the significance
of the situation. In temperament and gifts too there was a certain
kinship between Rousseau and the men who now came forward.
We shall bring forward a few of these men as typical examples.

\textsc{Thomas Carlyle} is not essentially influenced by Rousseau,
but proceeds from Goethe and from the poetry
and speculation of German Romanticism. But on the basis of these presuppositions, and driven
on by his melancholy and by the after-swell of Scottish Calvinism
in his soul, he arrived at a position toward the spiritual
and material culture of his time that forms a striking analogy to Rousseau's
position toward his time. Against science, which in its speculative
form made everything into logic, and in its naturalistic form
made everything into mechanics, he demands that a distinction be made between
the outer garment of existence, which must continually be woven in new forms,
% --- p. 126 ---
\setcounter{footnote}{0}%
\opage{126}and its inner core, which reveals itself in man's inner being.
He demanded the great wonder and the great reverence before
life and its powers. ``He who cannot wonder, who does not continually
wonder (and revere), --- were he president of innumerable
scientific societies, and did he hold the whole of the Mécanique
céleste and Hegel's Logic together with an epitome of the results of all laboratories
and observatories in his head, --- he is but
a pair of spectacles behind which there are no eyes.''\footnote[1]{Sartor resartus. Book I, Chap. 10.} --- And as he
set the living inner being of personality against the forms of science,
so in the religious sphere he set the inner subjective
conviction, won in silence and perhaps inexpressible, in
contrast to all dogmas. --- Like Rousseau he points
beyond the culture of the time and its warring oppositions to a new
spiritual culture, to be reached through the unfolding of the innermost
powers of personal life itself. --- There is a great opposition
on many particular points between Rousseau and Carlyle. To
mention just a single point: Carlyle approaches a cult
of power as justified in and for itself; but Rousseau declared
precisely the proposition ``force ne fait pas droit'' (might is not
right) to be the main proposition of the ``Contrat social''. But they meet in
their conception of the relation between personality and culture.
--- Akin to Carlyle's stance is that of Emerson and of Ruskin,
and what has just been said holds for them too. ---

In the Scandinavian North \textsc{Kierkegaard's} stance forms
an analogy to Rousseau's and Carlyle's. In his paradoxical proposition:
``Personality is the truth'' he, like his two kindred spirits,
sets the struggling yearning of personal life in contrast to
objective science as well as to ecclesiastical institutions and
traditional faith. He too felt himself in sharp opposition to his
age, which he accused of having forgotten what it was
to exist in earnest. ``What good would it do me'', he wrote in his
youth in his journal, ``if I discovered a so-called objective
truth, when it had no deeper significance for myself and my
life''. He belonged to the kind of thinkers he himself
has called ``subjective thinkers''. A subjective thinking proceeds
% --- p. 127 ---
\setcounter{footnote}{0}%
\opage{127}from personal life and permeates it. It does not dissolve
life into abstractions or let it be absorbed in fantastic ideals,
but works to hold fast and carry through the knowledge won
in the midst of life's reality. The subjective thinker is more
artist than scientist; his thought and his will go
together as one, and since life, with its changes and with the many finitudes
and externalities, constantly resists inward concentration,
there arises a relation of tension between life and thought
that can bring suffering and pain. Hence the deep pathos in
the subjective thinker's life.\footnote[1]{Cf. my book Kierkegaard som Filosof, Chap. III, 6---7.}

Kierkegaard is so far from being influenced by Rousseau
that he would rather have felt antipathy toward him if
he had known him more closely. But their position toward
the time in which each of them lived is analogous, just as
is their attitude toward the problem of life they had come
to see. And the same can be said of a man who felt himself
to be Rousseau's antipode, and would also have felt himself
to be Kierkegaard's antipode had he known him,
--- \textsc{Friedrich Nietzsche}. He too did not find life exhausted by
rational knowledge, by the ``Enlightenment''.\footnote[2]{After a disappointment in his relation to two rationalist friends he said to the faithful Overbeck and his wife: ``er müsse doch immer wieder etwas anderes haben; die reine Aufklärung genüge ihm nicht; die beiden verstanden davon nichts''. C. A. \emph{Bernouilli}: \textit{Overbeck und Nietzsche.} 1908. I, p. 338.} His antipathy toward
Strauss forms a parallel to Rousseau's antipathy toward Voltaire.
In his cult of Dionysus Nietzsche found expression for the overwhelming
enthusiasm, the captivating stirring of feeling, that for
Rousseau too marked the high point of life.

As Nietzsche sets Dionysus above Apollo, and Apollo
above Socrates, so he sets life above art, and art
above science. This thought is entirely Rousseauian.
And yet Nietzsche despised poor Jean Jacques and saw in
him one of the worst representatives of ``slave morality''. He counts
him among those he calls ``meine Unmöglichen'' (whereby Rousseau,
by the way, comes into quite good company). ``Rousseau, oder die
% --- p. 128 ---
\setcounter{footnote}{0}%
\opage{128}Rückkehr zur Natur in impuris naturalibus'' was an abomination to him.
Everything that stems from Rousseau was in Nietzsche's eyes
``false, contrived, bellows, exaggerated''.\footnote[1]{\textit{Götzendämmerung.} Streifzüge eines Unzeitgemässen.}

What is interesting about the figures analogous to Rousseau
in the nineteenth century is precisely that in many respects
they at the same time stand as his antipodes. It is against their
will that they come to bear witness to the value
of Rousseau's stance and to the justification of his problem.\footnote[2]{In a work by Seillière noticed in the Annales de la Société J. J. Rousseau (II, p. 284) there is even talk of Nietzsche's ``Rousseauisme''.}

In authors such as Guyau, William James and Henri
Bergson too one could find motives and courses of thought
that present analogies with Rousseau's, without there seeming, however,
to be any influence. ---

In \textsc{Leo Tolstoi}, on the other hand, both influence
and analogy can be demonstrated.

In the letter by which he enrolled in the ``Société J. J.
Rousseau'', Tolstoi stated (1905): ``Rousseau has been my
teacher since I was fifteen years old. Rousseau and the Gospel have been
the two great and good things that have determined my life. Rousseau
does not grow old. Just recently I happened to read some of
his works again, and I felt then the same elevation and admiration
as when I read him in my youth''.\footnote[3]{Annales de la Société J. J. Rousseau, I, p. 7.}

What Tolstoi has stated here is confirmed by his biography,
as Paul Birukof has presented it and Tolstoi himself
has looked it over. In a list of the books that have been
of significance for his development, he adds at ``Confessions''
and ``Émile'' that their significance has been ``enormous'', a predicate
which among the other books only the Sermon on the Mount and ``David
Copperfield'' receive. ``I have read the whole of Rousseau,'' he says,
``and there was a time when I more than worshipped him.... ...... Many
pages in his books are so close to my being that it seems to me that
I have written them myself.''

Tolstoi's relation to the literary circles in Petersburg and
Moscow presents a resemblance to Rousseau's relation to the
% --- p. 129 ---
\setcounter{footnote}{0}%
\opage{129}Encyclopedists. And like Rousseau he gave up urban civilization and
went out to the country to save his soul. Tolstoi stood, like
Rousseau, facing a circle of authors and publicists who regarded
themselves as the apostles of true culture, the leaders and teachers
of mankind, who were in possession of the highest enlightenment.
He found that they lacked self-criticism, that they were wanting in nature
and originality. Among the people with whom his life as a soldier had
brought him into contact he had found greater humanity.
And toward the great question of the aim and
meaning of life, the question that posed itself for him in ever sharper
form, the literary circles stood indifferent. But it was
precisely brooding over this question that led Tolstoi onto
new paths.

The opposition in which Rousseau and Tolstoi stood to
the literary circles of their times will repeat itself again and again. It
is an opposition between the two directions of intellectual life, the centripetal
and the centrifugal, unity and multiplicity,
inwardness and outwardness. With passion
Rousseau had taken his place in one of these extremes, and from
it he set forth his paradoxes. He and the minds akin
to him stand as prophets who address admonishing words to
the generation not to sacrifice the kernel for the shell and let life be
absorbed in externality. ---

A significant opposition between Rousseau and Tolstoi
appears in the fact that the fundamental thought of Rousseau's ethics is self-unfolding,
expansion, which springs from self-assertion (amour
de soi, distinct from amour propre), whereas Tolstoi preaches an
absolute altruism and tries to take up again the ethics of the Sermon on the Mount.
The distinction between amour de soi and amour propre
Tolstoi does not know. ``If you work for yourself alone, for
your personal future, then you must know that what awaits
you in the future is death .... Everyone must remember that when
he strives to preserve his life, he loses \textit{life.}
Jesus repeats this again and again. The true life is that which
increases the good that earlier generations have accumulated, makes this
inheritance greater in the present and leaves it to future generations.
In order for man to have part in this life, he must give up
% --- p. 130 ---
\setcounter{footnote}{0}%
\opage{130}his personal will''.\footnote[1]{\textit{Ma Religion.} 2\textsuperscript{e} éd. p. 141, 144.} --- Rousseau is far more an individualist
than Tolstoi. But he shows in an interesting way that it is
only through comparison with others, their circumstances and their
happiness, that amour de soi can pass over into amour propre. When,
on the other hand, self-unfolding has its free course, compassion
and love of humanity arise as expansive extensions of self-assertion.
Rousseau gains by this doctrine\footnote[2]{Cf. my book Rousseau og hans Filosofi. 2nd ed. p. 101---104.} firmer ground under his
feet than Tolstoi, since he demonstrates psychologically the possibility
of harmony between self-assertion and altruism, whereas Tolstoi,
through his rationalistic interpretation of the Sermon on the Mount, introduces
altruism as a mystical power. With this is connected the difference
which Benrubi has aptly brought out in his essay
``Tolstoi continuateur de Rousseau'',\footnote[3]{Annales de la Société J. J. Rousseau, III, p. 104 f. --- The Norwegian literary historian Chr. \emph{Collin} has also presented Tolstoi as one who takes up again Rousseau's problem of culture. See \textit{Leo Tolstoi og Nutidens Kulturkrise.} Christiania 1911.} that Rousseau lays the main weight
on the education of the child, whereas Tolstoi, who had after all
tasted more of worldly life than he, sets his trust especially in the
conversion of adults. Rousseau's men are ``once-born'',
Tolstoi's ``twice-born''.

\begin{center}III\end{center}


Besides through his influence on modern poetry and through
his position on the problem of culture, Rousseau also has significance
for several more special philosophical questions and their treatment.
A philosopher in the strictly scientific sense he was not, but
the conception of human intellectual life that he expressed with such
great warmth in his writings came to influence the position of later philosophers.

a. Here belongs first and foremost his influence on
psychology, especially his assertion of the originality
and independence of the life of feeling in relation to the other sides of mental life. As
predecessors here he has partly English psychology, which especially
influenced him through Diderot, partly the mystical tendency
% --- p. 131 ---
\setcounter{footnote}{0}%
\opage{131}within Cartesianism (Malebranche, Fénélon). But his glowing
enthusiasm, his sense for the sublime as well as for
the idyllic, for spiritual beauty (belle-âme in contrast to
bel-esprit) as well as for natural beauty, awakened the sense of and interest
in the life of feeling, so that one may rightly date a new
epoch for this part of psychology from his appearance. In the
nineteenth century, it is true, attempts have been made from various quarters
to reduce the life of feeling to a secondary side of mental life,
--- now (as in Hegel) to an obscure utterance, appearing in the form of ``immediacy'',
of what in clear and reflected
form comes forth in the life of thought, now (as in Herbart) to an
effect of the mutual friction of ideas. But these attempts
have not acquired any lasting significance. In modern
pragmatism, on the other hand, the life of feeling is asserted as
determining the life of thought, in sharp contrast to the view of those intellectualist
philosophers.

In Rousseau not only is great weight laid on the life of feeling,
but it is often isolated from the other sides of mental life, and
it becomes an object not only of contemplation but also of
enjoyment. Hence the sentimentality that Rousseau not only shared
with his age but himself helped to increase. Yet it must be
held fast that this isolation conflicts with his own psychology.
The life of feeling that conditioned the happiness he found in life was
not the blind state, untouched by all thoughts, that he glorified
in his first writings; it was a life of mood that gained
in inwardness and in scope by --- now as cause, now as
effect --- being bound up with thoughts. Nor, on
his own view, could feeling be isolated from will, for he laid
(in his later writings) great weight on the instinctive
foundation of mental life, the involuntary striving he called amour de soi,
and from which the most diverse mental processes of development
can proceed. When Schopenhauer in his concept ``Wille zum Leben''
comprised feeling and will under one, this was in accordance
with Rousseau's opinion (in his non-sentimental
moments).

The influence Rousseau has had on modern psychology
has for the most part been mediate, with \textsc{Kant} as
% --- p. 132 ---
\setcounter{footnote}{0}%
\opage{132}intermediary. But even where he merely stands as a predecessor,
later development has shown that he had an eye for an
essential side of mental life and gave the decisive impetus for its
being recognized and understood.

b. What is interesting in Rousseau's ethical conception is
the way in which in his later writings he shows the possibility
that the urge of self-preservation (amour de soi) can develop
into compassion, love of humanity and justice.
Earlier (in Discours sur l’inégalité) he had, on the contrary, set up
the urge of self-preservation and compassion as two original, mutually
opposed tendencies. In this first phase of Rousseau's
ethical conception \textsc{Schopenhauer} found a confirmation of his
ethics built on the originality of compassion. He certainly did
not arrive at this conception through Rousseau. It
is due to his own experience, and it stands moreover in close connection
with his general philosophical view, which asserts
the unity of all beings in spite of individual differences,
whereby it becomes possible for the pain of others to be felt as our
own. But Schopenhauer was very glad to find his
thought again in Rousseau. And although in another connection,
namely when he thinks of Rousseau's optimism and of his
faith (his ``Pastorentheologie'')\footnote[1]{\textit{Welt als Wille und Vorstellung.} II, Chap. 46.}, he speaks very disparagingly
of him, he praises him in high tones as an ethicist. ``My grounding'',
he says\footnote[2]{\textit{Die beiden Grundprobleme der Ethik}\textsuperscript{2}, p. 246.}, ``rests on the authority of the greatest moralist
of all modern times. For this is without doubt J. J. Rousseau,
the profound knower of the human heart, who drew his
wisdom not from books but from life, and destined his doctrine not for
the lecture chair but for humanity, --- Rousseau, this enemy
of all prejudices, this foster-son of nature, on whom she had
bestowed the gift of being able to moralize without being tedious,
because he hit upon the truth and moved the heart.''

In relation to Rousseau's later doctrine Schopenhauer (as
an ethicist, not as a psychologist) stands in a relation of opposition similar to
that of Tolstoi (who, after all, is also strongly influenced by Schopenhauer).
% --- p. 133 ---
\setcounter{footnote}{0}%
\opage{133}This later doctrine, which brings the life of feeling and volitional life into close
inner connection, nevertheless certainly marks an advance that cannot
be taken back.

c. Rousseau's strong emphasis on immediate and
involuntary mental life was bound to acquire quite special significance in religious
respects. Against the outward customs of the church and the dogmatic
doctrines, against books and speeches, he set the heart's
captivation, the feeling for life in its greatness and its smallness,
its happiness and its suffering. It was an open-air religion he
taught. But since all religion must have begun as open-air
religion, springing from the involuntary experiences of life, he
worked toward a renewal of religious life. He pointed
back to the source from which every serious view of life
must stem. He maintained that religious ideas can only
be interpretations of experiences, and he thereby initiated
a revision of the traditional forms of religion. To be sure,
he himself made an easy and rapid transition to dogmatics,
--- not to the church's dogmas, but to the dogmas of so-called ``natural'' religion,
and he was not free of fanaticism in his zeal to
maintain this ``natural'' theology. But here too,
as on other points, through his appeal to ``nature'' he points beyond
the warring oppositions of the time, the ecclesiastical repose in authority on
the one side, the scorn or blasé indifference of the Encyclopedists on
the other side.

In the nineteenth century these oppositions appear
anew, even if it happens in new forms. But again and again
the appeal to ``nature'' is heard. Schleiermacher's assertion of
the life of feeling as the real foundation of religion, while all
dogmas are only interpretations, is certainly not due to the influence
of Rousseau, but to the after-effect of a Moravian upbringing.
But Rousseau's problem is thereby posed anew. And the
same must be said of Ludwig Feuerbach's philosophy of religion, which
is marked by psychological sympathy with immediate religious
life and its involuntary manifestations. The whole of modern psychology of religion
is due to a resumption of Rousseau's point of view.
By the help of the history of religion and by the study of
the religious life of the present one seeks to get sight of the mental
% --- p. 134 ---
\setcounter{footnote}{0}%
\opage{134}motives that sustain religion and continually make it a spiritual
power.

It is a certain special kind of religious experience and faith
that Rousseau describes. It cannot be otherwise, since the description
is drawn from his own experience. Here again it is
self-assertion (amour de soi) that lies at the basis. It can
give birth to an urge to live wholly and fully, to strive beyond all
limits, to surrender oneself in admiration and enthusiasm. Then it becomes
religious feeling in Rousseau's sense. In spite of the disharmonies of his
life and the narrow circumstances within which it moved,
the source of expansion and enthusiasm did not dry up
in him. That this type of feeling is not the only possible one,
he did not see; only a comparative consideration could
show that. And it is especially in this respect that the
nineteenth century, with its greater spiritual horizon and
more comprehensive historical knowledge, has corrected and supplemented
Rousseau's thought.

The type in the psychology of religion that Rousseau describes, the expansive one,
recalls, as has been said, the one that English and American authors
in our day have called ``born once'' (once-born), in contrast
to the type that develops through crises and catastrophes
and can be called ``born twice'' (twice-born)\footnote[1]{William James has given a brilliant description of this difference. See my \textit{Mindre Arbejder.} Anden Række, p. 134 f.}. It was
characteristic of the age of the Enlightenment and of Romanticism that the former
type stood in the foreground. It was through the sharp
oppositions that the nineteenth century partly brought with it,
partly revealed, when the age of great hopes was over, that the
disharmonious type was properly noticed. This latter type
had, however, been the one the church had especially inculcated and favored,
and Rousseau's psychology of religion was a justified
reaction against the church's one-sidedness.

It is, however, not only Rousseau's psychology of religion, it is
also his theology, that is of interest in this connection.
While, as regards the psychological side of religion, he stood in
the sharpest opposition to Voltaire, the two opponents were
% --- p. 135 ---
\setcounter{footnote}{0}%
\opage{135}really agreed in their theology. Both maintained that the
evil in the world made it impossible to assume a God who
was both almighty and all-good. Significant thinkers
in the nineteenth century have arrived at the same result, without there
always being any question here of influence. Since the problem has been treated
under the same presuppositions and from the same point of view, the result
has quite naturally become the same. \textsc{Stuart Mill} was
indignant at the sophisms by which one sought to justify
the deity in the face of the suffering and injustice found
in the world. On the other hand he finds it possible to believe in a divine
being, to whom all good in nature is due, while
the material forces hinder its working. We may then in our
ethical striving stand as the deity's fellow workers\footnote[1]{\textit{Den nyere Filosofis Historie.} 2nd ed. II, p. 433.}. \textsc{Charles Renouvier}
too assumed that the power of the deity was limited,
and that therefore the development of disharmonious conditions in the universe
cannot be ascribed to it\footnote[2]{\textit{Moderne Filosofer,} p. 75. --- Cf. \textit{Les derniers entretiens de Ch. Renouvier.} Recueillis par Piat, p. 62 f.}. Renouvier is altogether a great
admirer of Rousseau and on several points strongly influenced by
him, especially as regards the basis for the treatment of the religious
problem\footnote[3]{See \textit{Esquisse d’une classification des doctrines,} I, p. 276 ff.; 459, II, p. 306 f.}.

Much suggests that not only philosophizing theists but
many people from practical life arrive at such results on account of their
experiences of life. One often meets in
life with a faith of this kind, especially in people in
whom clarity of thought is united with a living and warm
love of humanity. Love of humanity is angered at the attempts
to veil the misfortunes, and thought sees through the sophisms
in such attempts. It is a sign of the times that a
faith such as the one mentioned so often occurs, and one of the signs of the
times that give occasion to remember Jean Jacques Rousseau.

d. What is interesting in following Rousseau's modifications
and reshapings of the paradoxes he put forward in his first
% --- p. 136 ---
\setcounter{footnote}{0}%
\opage{136}writings on the relation between nature and culture lies in the fact
that he gradually worked his way forward to a modern concept of development
in the spirit of the doctrine of epigenesis. His protest against culture
in the name of nature he later modified in such a way that he
distinguished between a culture that is a natural fruit of individual
or social development, and one that does not stand as the
result of natural growth and self-activity. He denied that
culture could without further ado be transferred from one people to another.
In his pedagogy he emphasized the qualitative difference between
the different ages of life, especially between childhood and
adult age. The child is not an adult en miniature; but
childhood is a peculiar form of existence that makes its
demands and has its rights. At about the same time C. F. Wolff,
the founder of the doctrine of epigenesis in anatomy, maintained that the transition from
the germ to the fully grown organism takes place through successive
changes, without the earlier stages needing to resemble the
later. Growth is something other and more than enlargement.

Rousseau's conception of development led him to lay
weight on what he called negative education, which essentially
consists in keeping obstacles away so that natural growth
can take its course, and in seeing to it that influences come
only, and tasks are set only, when this natural growth is
sufficiently advanced.

It is through the clear and sure application of this concept of development
that Rousseau stands as the founder of modern pedagogy.
All significant thoughts in the pedagogy of the nineteenth century
point back to him.

\begin{center}IV.\end{center}


Through his criticism of the one-sidednesses of modern culture
Rousseau also stands as the one who posed the social problem
that since his time has followed the generation on its way.
Against the isolating and mechanizing effects of the division of labor
he demands that every human being shall have the opportunity to
develop wholly and fully, according to his individual peculiarity.
His worth does not depend merely on his being a link in a
machine. It was this line of thought that led Kant to
% --- p. 137 ---
\setcounter{footnote}{0}%
\opage{137}set up what has since been called the principle of personality, the
demand that every human being shall be regarded as an end, never
as a mere means. This principle is tacitly presupposed in
modern socialism, as in Karl Marx, as well as
in the modern attempts to ground morality on the principle of solidarity\footnote[1]{See my \textit{Etik.} 3rd ed. p. 390 (on Marx). --- \emph{Leon Bourgeois} in \textit{Essai d’une philosophie de la solidarité,} Paris 1902, p. 38.}.


Rousseau could of course not foresee the ways by
which this thought was to be carried over into reality. He
still stood polemically opposed to social formations, in his fear
that natural freedom and immediacy would suffer.
Social formations such as craft guilds and fraternities his time knew only
as obstructive factors, degenerate forms of the past. Turgot,
that great admirer of Rousseau, in his suspicion
of all associations\footnote[2]{\textit{Vie de Turgot} (par Condorcet), Londres 1786, p. 81; 228 f. Of interest here also is Turgot's judgment of the clique that the Encyclopedists formed (p. 28).}, forbade as minister all combination of
workers, and this prohibition was repeated during the Revolution.
Here faith in ``natural'' development had an obstructive effect.
For it has turned out that only by way of association, through
the great combinations of workers, has it been possible to achieve
progress for the fourth estate and a share for it in the goods of
culture. As we have seen above, Lamennais expressed
this in clear words, ahead of the English trade-union movement
and ahead of the program of Marxism. --- Development has, however,
not entirely shown Rousseau's and Turgot's misgivings to be unjustified.
The great question of the rights of minorities
shows this.

Rousseau is therefore not a socialist in the modern sense of the word.
A certain state socialism he does, however, advocate in his
article ``Économie politique'' in Diderot's Encyclopedia. But
he is a predecessor of the socialism of the nineteenth century
through his whole conception of the problem of culture and through his demonstration
of the disadvantages of the division of labor. Far beyond the limits of his own
% --- p. 138 ---
\setcounter{footnote}{0}%
\opage{138}century the words, animated by indignation and enthusiasm,
with which he has pointed to the Achilles' heel of modern culture,
have their effect. One of the modern historians of socialism\footnote[1]{\emph{William Graham}: \textit{Socialism new and old,} London 1890, p. 55.}
says of this: ``Not since Luther's voice was heard, indeed, hardly
since the words of the Gospel were proclaimed, had words sounded so charged
with far-reaching consequences, words that awakened thinkers and
poets, the middle class and the common people, words that have been the source
of all later revolutionary, communist and socialist literature,
and whose influence will be felt as long as the earth circles
in space.''

Whether this mighty word has acted most directly or indirectly
in the course of the nineteenth century is difficult to
decide. For a large part, here too it has been the analogy
in the situation that has called forth kindred thoughts. But
the main thing is that he was the first to see, with clarity and with feeling,
the great problem whose sting continually sets thoughts and
passions in motion. ---

In relation to the political thoughts and movements of the nineteenth century,
too, one is led back to Rousseau. He is the
first consistent champion of popular sovereignty. The natural-law
movement, which had already begun in the Middle Ages, found
its conclusion in \textit{Contrat social.} Whereas all earlier teachers of natural law\footnote[2]{With the exception of Hobbes, who, unlike Rousseau, let the ``pact'' of the formation of society be absorbed in the ``pact'' of submission to an authority.}
distinguished between two ``pacts'', one by which
society was founded, another by which power was transferred
to a ruling authority, Rousseau excluded
this latter ``pact'' and thereby completed the theory of popular sovereignty.


This Rousseauian doctrine has in the nineteenth century
been attacked from two very different sides.

The so-called historical school, which stood in close connection
with Romanticism, conceived the state as a divine order
grown up through the ages, against which particular
% --- p. 139 ---
\setcounter{footnote}{0}%
\opage{139}individuals could not set their accidental experiences, ideas
or wishes. The state is the organized conscience of society,
in relation to which the conscience of the individual is of vanishing
significance. The present cannot be isolated from the past, nor the
individual living for the day from the eternally persisting state. \textsc{Hegel}
stood in close relation to the historical school, insofar as he
too could not see in particular individuals living points of departure
for the sense of right and the formation of law. He acknowledges
it as a merit of Rousseau that he made will,
and thus an ideal principle, the principle of the state, but reproaches
him that by ``will'' he thinks partly of the will of individuals,
partly of what is common to these wills, not of ``das an und für
sich Vernünftige des Willens'', which holds whether the
individual acknowledges it or not\footnote[1]{\textit{Die Philosophie des Rechts}\textsuperscript{2}, Berlin 1833, § 258 p. 314.}.

It is not only from the conservative side that Rousseau is attacked
for his individualism. From the socialist, especially the syndicalist,
side he is attacked as an anarchist and put on a line with Stirner,
Reclus, Kropotkin and Tolstoi. These ``bourgeois'' regard,
in their fantasy, the individual as something absolute and do not see that the individual
means something only through his productive labor!
While socialism would have the state take over all production and
thereby take all individuals into its service, syndicalism hopes
to be able to make the state superfluous and impossible through free
organization of labor and through the joy in work that
thereby becomes possible. This latter tendency therefore attacks
Rousseau both as an anarchist and as one who believes in the state's
continued persistence\footnote[2]{Cf. \emph{Challaie}: \textit{Le syndicalisme revolutionaire} (Revue de Métaphysique et de Morale, 1907). --- \emph{Guy Grand}: \textit{Le procès de la démocratie} (ibid. 1910).}.

It cannot now be denied that the Contrat social invites
criticism. Rousseau himself was not satisfied with it. ``As for my
Contrat social'', he says\footnote[3]{Reported by \emph{Dusaulx}: \textit{De mes rapports avec J. J. Rousseau,} Paris 1798, p. 102.}, ``those who boast of
understanding it are cleverer than I. It is a book that must be
% --- p. 140 ---
\setcounter{footnote}{0}%
\opage{140}rewritten.'' --- Three different principles clash in the book: the individual
wills in their multiplicity, --- their interaction and
crossing, --- and finally the fundamental will that represents the people's
self-assertion through the ages, and that Rousseau calls
la volonté générale (the people's fundamental will) and definitely distinguishes from
la volonté de tous (all the wills of the individuals taken together). The
first principle is anarchistic, the second mechanical, the third
mystical. He did not succeed in getting an inner connection
among these three principles. And in a certain sense no theory of the state
has yet succeeded in getting such a connection.
But it is nevertheless Rousseau's clear endeavor to
distinguish between the interest of the people as that which is to persist through
the ages, and the particular individuals or generations of individuals,
none of whom may arrogate to themselves the right once and for all
to speak in its name. When Rousseau says (II, 3. 6) that the fundamental will
does not err, it is precisely because the people or society
cannot will its own dissolution, however many illusions
may prevail among particular individuals and in particular
generations. But on the other hand the fundamental will expresses itself
in particular cases only through particular individuals and in particular
generations, --- and here the problem comes in.

Whether this problem can be solved or not,
Rousseau is right that a faith in the people's ability and will to
persist lies as a hidden presupposition behind all politics and all
political arguments\footnote[1]{Cf. the brilliant discussion of democracy in \textit{Bulletin de la société française de philosophie,} VII (21 Dec. 1906). It was in large part a discussion about Rousseau.}. And Rousseau saw clearly that the life of the state
cannot at any point in time begin absolutely afresh.
He expressly insists (III, 8) that the people's past and its
conditions of life must be determining for its constitution, and it
was for this reason that he answered dismissively the naive requests
he received from several countries (Poland, Corsica) to write a constitution
for them. If he, he answered, had lived for ten years
in their country, he could perhaps write its history; but
even then he would not be able to write its constitution\footnote[2]{See my book Rousseau og hans Filosofi. Second ed. p. 87.}.
% --- p. 141 ---
\setcounter{footnote}{0}%
\opage{141}Rousseau thus does not lack, as has often been thought, historical
sense. He does not place the individual or the individual generation
outside the historical connection; in this respect
he can very well be regarded as a forerunner of the historical
consciousness of which the nineteenth century has been so
proud. The zeal with which, in the various countries of Europe, in
the course of the nineteenth century and in remarkable contrast
to ``the historical consciousness'', new state constitutions have been fabricated,
Rousseau criticized in advance. Just as little as
he agreed with the Encyclopedists that a spiritual culture
could without further ado be transferred from one people to another, just as
little did he consider it possible that a people could fetch its constitution
from another people. His doctrine of the connection of culture
with nature has found its confirmation through a
series of fateful experiences in the course of the nineteenth century.

% --- p. 142 ---
\setcounter{footnote}{0}%
\opage{142}\essay[Rousseau and Religion (1912)]{ROUSSEAU AND RELIGION.}
\streg

\begin{center}1.\end{center}

R\textsc{ousseau's} psychological conception has sprung from his
own personal experience. When he wrote his ``Confessions'',
one of his motives was that he wished to present his own
mental life as a psychological example. And --- so far as
we can rely on his account\footnote[1]{Cf. my book: Rousseau og hans Filosofi. p. 13---22.} --- we see how the fundamental
psychological traits on which he builds in his theory, especially in
his conception of religion, are generalizations of what he
has experienced and found in himself. To the Confessions are joined
his letters, La Nouvelle Héloïse, Émile, Dialogues (Rousseau
juge de Jean Jacques) and Rêveries d’un promeneur solitaire.

There appears in Rousseau, not only in his self-characterization,
but in his whole conduct, a great opposition between
feeling on the one side, cognition and will on the
other side. His urge is to live out along one straight
line. Circling and compound movements are unnatural to him.
Every mood that emerges spreads and rules over the
mind until its energy is exhausted, --- and its energy is exhausted
the more quickly, the more vivid it is. When, then, a new
mood replaces an earlier one, perhaps a wholly opposite one, the
memory of what was thought and resolved in the earlier one is not preserved.
Therefore he can go from extreme to extreme without
noticing the inconsistencies that thereby arise. The
% --- p. 143 ---
\setcounter{footnote}{0}%
\opage{143}involuntary moods and the self-surrender they make possible are the
most valuable; but what he resolves in them is often not
carried out, because the whole force of the will is exhausted in the resolution, and
no force remains for the execution, since he sinks back
into the dreaming state which is now a kind of idleness,
now a poetic rapture, but which in both cases is
his favorite state. Through this passive character of his emotional life
he became highly dependent on his surroundings. ``Je
suis ce qui plaît aux hommes'', he says himself. Others could
win him over, or tease him, or put him in embarrassment.
He lacked presence of mind for reply and counterattack. Only
in one period, --- when the problem of culture had dawned on him, ---
when he had undertaken the ``moral reform'' of his way of life
and thereafter felt himself the prophet of his age, --- did indignation make
an energetic bearing possible for him, which astonished his
whole circle. But when he had left Paris, and when in
the solitude of country life there was no occasion for wrath, the
energy of prophecy was again replaced by rapture. The undivided states
held sway again, and the resistance he had met, and the enemies
he had had (or believed he had had), were then at once
forgotten. In his solitude it was the immediate
feeling of life, the delightful sensation of existing, that filled
him, and that, especially when he roamed about in the woods or let
his boat drift with the current, found vent in smiling memories
and fantasies.

Is it possible to be consistent? he asks. He does not believe
so, and we understand from the foregoing why he does not
believe so. Each single state gains complete power over him.
To think, to compare, was always for him a painful and little
attractive activity. Thinking presupposes, after all, that the ideas
in which the changing states find expression can be
preserved and held together. But that was precisely what Rousseau
feared and shunned, because such holding together, and the relations
it leads one to establish, conflicted with the ecstatic absorption in
the single state.

Even if it is not possible to be wholly consistent, it
is nevertheless, Rousseau thought, possible to be honest. And he
% --- p. 144 ---
\setcounter{footnote}{0}%
\opage{144}prides himself on his honesty and on his repentance. He writes in a letter:
``Voici ce qui me distingue de tous les hommes que je connais:
c’est qu’au milieu de mes fautes, je me les suis toujours
reprochées''. --- But honesty and repentance surely presuppose thinking, comparison! He must, then, have been able to gain
a survey of his changing states, however extremely opposed
they might be; in him there was an unceasing oscillation between
the highest and the lowest. Or else repentance in him
had a passive, dreaming character.

And yet Jean Jacques was a brave man. He could not
in the long run appear as a prophet quick to strike back, but when his ideal of life
--- the whole, undivided living in the immediate state of the mind
--- had in earnest dawned on him, he sought out solitude
in spite of his friends' mockery and unmoved by all their attempts
to play Providence and draw him back to the one saving
salon life. He now guarded his independence as the most
precious good of his life, and the fear of losing it overwhelmed
his fantasy to such a degree that it degenerated into persecution mania.

And from his solitude he then sent out the great works
that were directed against the radical as well as against the reactionary
tendencies of the time and therefore aroused bitterness on both sides.
In contrast to Voltaire he put his name on his books and
took the danger and persecution that followed. Here too he was
both honest and brave. However hard he found it to
think, he nevertheless managed, through a series of fruitful
impulses, to deepen his first paradoxes about the opposition
between nature and culture and to hold up a great ideal of freedom and originality
before the age.

Finally he was also, in spite of the constant alternation of the
high and the low in his life and in his moods, clear about
where the high points of his life lay. They were the moments of great enthusiasm,
when the movement of the mind was raised to dizzying ecstasy,
and when no thought and no word could express what he
felt. In such moments he felt himself one with the innermost
of existence. And such moments he called his religious
moments. Then ``nature'' had overcome the fragmentation
% --- p. 145 ---
\setcounter{footnote}{0}%
\opage{145}and externality of ``culture'' and proved its unity in spite of all resistance. This
highest feeling of unity was Rousseau's personal religion. ---

It will now appear that the merits and imperfections
of his religious standpoint are closely connected with his own personality
and the experiences it made possible.

\begin{center}2.\end{center}

In Lettres de la Montagne Rousseau says that the confessions of faith
of the Savoyard Vicar and of Julie agree so closely that
they can mutually illuminate each other, and that their author
sympathizes with both. On closer consideration, however,
a certain difference appears between the standpoint of Julie, of the Vicar, and Rousseau's
own. These are three types whose difference has its
psychological interest.

Julie is first and foremost the loving woman, like
the medieval Héloïse. But just as love at its
height can take on the character of enthusiasm and devotion, so
conversely the religious feeling borrows the language of love
in order properly to find expression. And the kinship that shows itself here
between erotic love and religion Rousseau sought to portray in Julie's
character and course of development, without, however, succeeding
in making this portrayal natural and probable. Julie's character
was to be an example of the thought that Plato expressed in so brilliant
a way in the ``Symposium''. But --- it is not the
aesthetic criticism of ``Julie'' that concerns us here. As
she is portrayed, it is her urge to love that leads her
to religion, because nothing earthly and finite could take full
possession of her ``affections surabondantes''. Here she recalls
Ste Thérèse; but her religious love does not, like
Ste Thérèse's, stand in contrast to the love of human beings, so
that it must exclude the latter; the religious and the human
love mutually strengthen each other.

The Vicar's confession of faith is the fruit of a seeking, a criticism,
called forth by the opposition to the positive dogmatism of the Church
on the one side, the negative dogmatism of the freethinkers
on the other side. He appeals to the inner light, to
% --- p. 146 ---
\setcounter{footnote}{0}%
\opage{146}the testimony of the heart. His only concern is to satisfy the heart's
urge; all further-reaching theological questions he puts
aside. And yet he makes use of philosophical reflections, gives a whole
system of ``natural theology'' and applies with full consciousness
the analogy with human will to explain the nature of the Deity.
His religion is the expression of his gratitude
for existing, which leads him to admiration and worship of
the highest power. ``De mon premier retour sur moi naît dans
mon cœur un sentiment de reconnaissance et de bénédiction
pour l’Auteur de mon espèce, et de ce sentiment mon premier
hommage à la divinité bien-faisante. J’adore la puissance suprême,
et je m’attendris sur ses bienfaits. Je n’ai pas besoin qu’on
m’enseigne ce culte, il m’est dicté par la nature elle-même. N’est-ce pas une conséquence naturelle de l’amour de soi,
d’honorer ce qui nous protège, et d’aimer ce qui nous veut du
bien?'' (Émile. Livre IV).

For Rousseau's own personal part, religiosity stands
in the closest connection with his urge for independence
and solitude, --- to live out of himself. In free
nature, far from human beings, his inner life could properly unfold.
From nature and its glory he rose to the ideal world of fantasy
and raised himself --- not least when the contrast
with reality made itself felt --- to the idea of an infinite
being which no thought, no word could express. In
the third letter to Malesherbes we have the most characteristic
expression of his cosmic feeling of life. ---

Common to all three types is the great opposition that
underlies Rousseau's psychology and is at work in his whole
personality, --- the opposition between ``nature'' and ``culture'',
as he has presented it most clearly in his Dialogues
(Rousseau juge de Jean Jacques). To this opposition corresponds
the opposition between absolute and relative feeling, --- between states
in which one lives wholly and undividedly, and states that
presuppose opposition and comparison. His psychological
theory is built on the experiences of his own life, on the great
opposition between times in which he could give himself over
wholly to the unfolding of his rich nature, and other times when he
% --- p. 147 ---
\setcounter{footnote}{0}%
\opage{147}was hampered and embittered by his relation to people with
whom he could not feel himself one. In the distinction between absolute
and relative feeling, amour de soi and amour propre,
he has a predecessor in Vauvenargues and older Cartesians. The
absolute feeling is one with the instinct of self-preservation, is an involuntary
urge to assert and expand one's being. So long as it
can easily overcome the resistance it meets, its
nature is not altered. But when the obstacles become stronger, so
that they draw attention to themselves at the expense of the original
end, which was mere self-assertion, relative feelings arise
such as anger, hatred and fear. And when imagination and
thought make it possible for the individual to compare himself with
others, his lot with theirs, the relative feelings are favored
at the expense of the absolute. Amour de soi becomes amour propre.
Needs, and with them occasions for discontent, fear
and hatred, grow.\footnote[1]{It is perhaps characteristic that in ``Rousseau juge de Jean Jacques'' Rousseau particularly stresses the provoking effects of the obstacles, whereas in his earlier accounts he lets fantasy and comparison unfold and work more without external occasion.}

When Rousseau speaks of comparison as the cause of
relative feeling, he is thinking especially of the perception of difference. Likeness
to other beings, on the other hand, will not hinder self-unfolding.
It can lead us to take others up into the involuntary stream
of our self-unfolding, involuntarily to put ourselves in their place,
to rejoice and suffer with them and help them. Pity and love
therefore belong (in Rousseau's later accounts) to the
domain of absolute feeling: La force d’une âme expansive
m’identifie avec mon semblable! And therefore love of mankind
can spring from the urge to self-preservation.

This theory of absolute and relative feeling has sprung
from deep psychological understanding. It reappears, remarkably
enough, in Nietzsche, who was nevertheless hardly influenced by Rousseau,
whom he deeply despised. But whereas Nietzsche counts precisely
pity, love and the feeling of equality among the relative
feelings, in Rousseau it is fear and hatred that stand as the
most important relative feelings. An inconsistency such as the one Nietzsche
% --- p. 148 ---
\setcounter{footnote}{0}%
\opage{148}commits in extolling the feeling of distance (das Pathos der Distanz),
directed from above downward, especially contempt, as the mark of the great,
is not found in Rousseau. On the other hand it is an imperfection
and an inconsistency in Rousseau when he regards all relations of difference,
especially all obstacles, as unfavorable to self-assertion.
That means impotence and isolation in the face of reality.
When self-preservation and self-unfolding have sufficient energy,
they will be able to bridge the distance and become richer by
entering the world of differences. Rousseau was only himself,
was only ``what nature [in his opinion] had willed'',
when on his solitary walks he was free of all obstacles
(sans diversion, sans obstacle, Rêveries II). This is connected
with the fact that his emotional life so easily came to stand in an
inverse relation to his life of thought and his volitional life, and thereby took on a
dreaming character. Obstacles can, after all, release new forces,
hitherto hidden qualities, and thereby indirectly further
emotional life too. And reverie gains content and value precisely only
when it stands in reciprocal action with work. Precisely the opposition,
the relation of contrast, has its significance here, which Rousseau
himself has experienced, when he says that it is to his adversaries
that he owes the ecstasies on his solitary walks.
His first, paradoxical opposition between nature and culture
is still at work here. And it had in turn sprung from his romantic
and vagabond life, in which every relation came
to stand as an intolerable bond.

Religious feeling is now, according to Rousseau, like love,
an involuntary expression of self-assertion. It leads us
to believe in life, to feel ourselves one with nature, ``se fondre
pour ainsi dire dans le système des êtres'', and at its height
to religious ecstasy (Rêveries VII; 3e lettre à Malesherbes).
Only in this state do all relations, all possibility of
comparison, fall away. No thought, no image would be able to
express what is experienced in such moments. Both Julie
and the Savoyard Vicar speak of this in the same way
as Rousseau when he speaks in his own name. There are three
reasons why no thoughts or images suffice: the fullness
of what is experienced; --- there is no relation between what is
% --- p. 149 ---
\setcounter{footnote}{0}%
\opage{149}experienced and what is otherwise known; --- the total concentration
of all faculties leaves no force over for comparing.

In Rousseau's last years these ecstatic states took on a
somewhat different character than before. There was lacking the violent
force with which the inner longing could formerly gain
free course. In return his solitary experiences were marked by
a resignation that had long been hindered by amour-propre.
Now, when the external relations and the comparisons they had
given occasion for had receded, amour propre could
again become amour de soi. This made possible the beautiful,
mildly resigned mood that prevails in Rêveries d’un promeneur
solitaire. ---

Through the religious feeling, as absolute feeling, being set
in opposition to all relations, since these constantly threaten to
annul it, there arises the great one-sidedness of Rousseau's psychology of religion.
The religious feeling does not become a feeling
that includes in itself and expresses the result of life's struggle and
suffering, of its tragedy and comedy; its stamp is determined
precisely by life and its oppositions being kept away. There
is lacking an element of courage and an element of fellow-feeling, which cannot
be lacking in the feeling for which the many-sided experience of life
is to form the foundation. ---

Where thought is not a hindrance to feeling, it is,
according to Rousseau, an effect of feeling, at any rate in
the religious domain. The Savoyard Vicar limits his
search for truth to what is of interest to him in a practical
respect; the rest he leaves aside. His ideas
of God he grounds on the feeling of the relation in which God
stands to him. Rousseau is clearly aware that reason in
the long run takes on ``le pli que le cœur lui donne''.

From all this it follows that the religious problem could not
present itself to Rousseau in its sharpest form. The purely objective
view of the world never came, in him, to be set in relation
to the heart's urge. In the strongest spiritual strife he
has never been placed. The cosmological problem has not
presented itself to him in its difference from the religious problem. He
% --- p. 150 ---
\setcounter{footnote}{0}%
\opage{150}has never doubted. In the remarkable letter to M*** of 15 janvier
1769 we read:

``Vous me marquez, Monsieur, que le résultat de vos recherches
sur l’Auteur des choses est un état de doute. Je ne
puis juger de cet état, parce qu’il n’a jamais été le mien. J’ai
cru dans mon enfance par autorité, dans ma jeunesse par sentiment,
dans mon âge mûr par raison; maintenant je crois, parce
que j’ai toujours cru''.

Perhaps it is so with most people. Even those who set
themselves in the sharpest opposition to the authorities of their childhood
and to the feelings of their youth are determined, directly or indirectly,
through continuity or through contrast (which is also
a kind of continuity), by the subterranean currents of their
earlier life. Rousseau belonged to ``the once-born''. He
went through his various conversions --- from Protestantism
to Catholicism and back again, and from both to his emotional
rationalism, --- and through his various circles
(pious Cartesians and impious Encyclopedists), --- untouched,
with the voice of ``nature'' in his heart. Although, as we shall
see, he criticizes the doctrines of the Church on essential points, and
although the problem of evil in particular has had an influence on
his religious ideas, he has nevertheless not carried out any real
work of thought in the domain of the religious problem. His
merit lies in his having decisively referred the religious
question to feeling and to personal urge. That is a
great merit, even if not all questions are thereby settled.
He finds in the ideas of so-called ``natural'' religion
the adequate and only valid expression of the highest,
without examining the justification of the analogies on which
these ideas rest. The Cartesian philosophy and the
deist tradition have exerted influence on him here. There is
in his religion something other and more than what immediate feeling
and a faith grown out of experience of life must necessarily assume.
% --- p. 151 ---
\setcounter{footnote}{0}%
\opage{151}\begin{center}3.\end{center}

Rousseau has an entire ``natural'' theology, which he energetically
sets in opposition partly to ecclesiastical orthodoxy, partly to materialism.

God is, the Vicar teaches, an infinite being which none of
our concepts or images can express. Even the concept of spirit
is not adequate; for God and our spirit cannot be of
the same kind. For God all truths are contained in a single
thought, all places in a single point, all times in a single
moment. Whether he has created the world is not known; but the indestructibility
of matter and the reality of evil speak decidedly against it;
the very word creation in an absolute sense is unintelligible, and it
is in any case too abstract to be able to be a biblical
concept. But the order that prevails in the world, and the evident
way in which everything leads in the direction of the same ends,
shows that there is a first moving force, a God, who
directs the course of things in the world.

Against the materialists he emphasizes that we must think
every cause of motion in analogy with our own will.
When we begin the analysis of the universe with our own being,
instead of with the dead masses, we shall find the key to the
arrangement of the universe in ourselves. That the soul must be a being
distinct from matter follows from the very faculty of thought, the faculty of
comparing, a faculty that cannot be ascribed to matter. It
is a faculty of activity that stands in opposition to passive
sensation, which is due to the influence of the body on the soul.

The distinction between an active and a passive principle which
Rousseau is thus led to make in accordance with his ``analysis of the universe''
now finds its application to his position on the
problem of evil. Here too, moreover, his psychological conception
makes itself felt.

This problem has, by Rousseau's own declaration, never
made any strong impression on him. It arises only when
one regards God as almighty creator, and when one overlooks
to how great a degree the misfortunes in the world are due to man himself.
For his own part Rousseau solves the problem by conceiving
% --- p. 152 ---
\setcounter{footnote}{0}%
\opage{152}God as an active principle that can intervene, ordering and altering,
in the material processes, but yet often meets with resistance
from them, --- and moreover by pointing to the source of happiness
that already lies in the immediate feeling of life, when
all comparisons and refinements that are due to human beings
themselves are kept away. Our unhappiness consists in the disproportion between
urge and faculty, and it is fantasy and comparison that have
inflamed the urge, so that the faculty no longer suffices.

In the famous letter to Voltaire (1756), written on the occasion
of the latter's poem Sur le désastre de Lisbonne, Rousseau's
position on the religious problem, and especially on the problem
of evil, appears in a characteristic way. The comparison
which this letter occasions between Rousseau's and Voltaire's
religious standpoints is moreover of considerable interest
from the standpoint of the psychology of religion.

In his Poëme sur le désastre de Lisbonne Voltaire had
given vent to his indignation at the belief that the evil in
the world could be reconciled with the world's origin in an almighty
and all-good being. With great force the usual
insipid grounds of consolation and edifying phrases with which
people meet great misfortunes are rejected, and there is also polemic
against the aesthetic-contemplative way of regarding the disharmonies
of existence that makes itself felt in many optimists,
not least in Leibniz. The problem is posed sharply and
severely, and the whole poem has a dark and bitter character. As
usual Voltaire declares, which in his mouth is of course only a
phrase, that only a revelation can solve the riddle.
--- In later editions of the poem Voltaire, ``instruit par la
vieillesse'', changed its ending and inserted some lines in which
the hope of a better future finds voice:
\begin{quote}
``Un jour tout sera bien'', voilà notre espèrance,\\
``Tout est bien aujourd’hui'', voilà l’illusion.
\end{quote}

But it was not in this form that the poem lay before Rousseau, whose statements
become unintelligible if one does not keep to the poem's
original form.

In his reply to Voltaire's poem Rousseau finds use for
% --- p. 153 ---
\setcounter{footnote}{0}%
\opage{153}his favorite theme, the opposition between nature and culture. How
many evils are there not that human beings bring upon themselves
--- not least by moving together into great cities? A disaster
like that in Lisbon would not strike a scattered rural population,
which moreover does not live in houses of 6---7 stories, so hard. And as regards
the question of the value of life, experience shows that human beings
prefer existence to non-existence, which can only be
explained by the immediate feeling of life having the character
of a feeling of pleasure. But this doux sentiment de l’existence,
which can be found, and is most readily found, where life is led under
limited circumstances, is as a rule overlooked when the account
is drawn up. And in addition, one poses the question wrongly
when one infers from the value of life and of existence to God's
existence. It is the opposite way one must go: he who
believes in God, and thereby has the foundation for a great hope,
will always see existence as valuable. Rousseau
later has Julie in her ``swan song'' give the ground for this thus:
Toutes les misères s’évanouissent devant un plus grand objet. ---

In the letter to Voltaire the subjective, personal turn
that Rousseau has given the discussion of the religious problem
has found its first energetic expression. Henceforth the
discussion could really turn only on whether an urge such as
the one Rousseau refers to is universally human, --- how
it arises, --- whether it always expresses itself in the same way and
leads to the same postulates, --- and what significance it has
for human intellectual life. Rousseau does not go further
into this, any more than Kant did later. ---

For the first time, too, this letter brings out the theological
conception that Rousseau later presented in detail in
Émile, in the letter to the Archbishop of Paris and in the letter to M***.
He reproaches Voltaire for not sacrificing omnipotence
rather than goodness. ``Si l’embarras de l’origine du mal vous forçait
d’altérer quelcune des perfectures de Dieu, pourquoi vouloir
justifier sa puissance aux dépens de sa bonté? S’il faut choisir
entre deux erreurs, j’aime encore mieux la première''. --- His
later reflection came, starting from the general assumption of an
active and a passive principle in existence, to the same result.
% --- p. 154 ---
\setcounter{footnote}{0}%
\opage{154}The divine spirit constantly has a resistance to overcome
in matter and therefore cannot always, and not immediately, realize
its ends. When Rousseau said in the letter to Voltaire
that instead of the almighty God who in Voltaire's poems was depicted
as the author of all evil, although he could have prevented
it, he would prefer Manichaeism, he later shows
(in the letter to the Archbishop) that his doctrine is not Manichaeism,
since he assumes not two gods, two active principles,
but one active and one passive principle.

When Rousseau defends Leibniz against Voltaire, it is
because he has seen that Leibniz in fact also assumed two principles,
only that he (like Böhme) conceived them as elements
within the Deity itself: the one the selecting, realizing
will, the other the intellect, which presents the many possibilities
among which a choice must be made\footnote[1]{See \textit{Den nyere Filosofis Historie}\textsuperscript{2} I p. 368.}. For Rousseau himself
only the first element belongs to God's nature, while the
many possibilities among which a choice must be made are contained
in matter.

The theological turn that Rousseau intimated in the letter to
Voltaire and later carried through in \textit{Émile} (1762) and in the \textit{Letter to
the Archbishop} (1763), Voltaire for his part made
later, in \textit{Dictionaire philosophique} (1764)\footnote[2]{The view at which Rousseau and Voltaire thus arrived had already been expressed earlier by \textsc{Ludvig Holberg}, as the result of his brooding over the problem of evil raised by Bayle. Cf. \textit{Danske Filosofer} p. 14 f. Among later thinkers who have arrived at a similar result must be mentioned \textsc{Stuart Mill}, \textsc{Charles Renouvier} and \textsc{William James}. (Cf., on the latter's religious view, the excellent characterization by \textsc{Th. Flournoy}: \textit{La philosophie de William James.} 1911. Chap. IX).}. And thus
the two men, who in so many respects were antipodes, came
to agree in their theology. They stood together in the struggle
against orthodoxy. But this agreement does not mean
much, since such a negative attitude can unite the
most opposed standpoints. On the other hand it is of the greatest psychological
significance to see the great personal opposition that
can be bound up with agreement in religious
% --- p. 155 ---
\setcounter{footnote}{0}%
\opage{155}ideas. The difference lies chiefly in this, that Voltaire is filled with
joy at the ``culture'' and refinement of the age, when his indignation
is not aroused by remnants of medieval brutality among the
rulers, with whom he otherwise tries to keep on good terms,
whereas Rousseau deeply felt the externality and superficiality of
the culture of the age, its lack of originality and naturalness.
Moreover, Rousseau had in his joy in nature and his enthusiasm
sources of life from which he could constantly draw new spiritual nourishment,
and by whose help he was kept up through illness, want and
persecution. This gave his religiosity inwardness and originality
and paved the way for a new way of treating the religious
problem. It led to religious ideas
being regarded as derived in relation to the lived experiences and
practical experiences that are the real psychological foundation.

\begin{center}4.\end{center}

In Rousseau's religious conception the Christian tradition
makes itself strongly felt. In the period of his life when his thought
was most searching, there were especially two currents by which he was influenced:
the Cartesian mystics of the 17th century (Pascal,
Fénélon) and the pietist movement that from Germany penetrated
Romance Switzerland, and by which his foster-mother
Mme de Warrens had been so strongly influenced. When,
after the wanderings of his youth from confession to confession and
through the purgatory of the Encyclopedists, he drew up his account,
the result was that he claimed the right to call himself a Christian,
although he set aside essential aspects of the Church's Christianity.
In the proud letter to the Archbishop of Paris he says:

``Monseigneur, je suis Chrétien, et sincérement Chrétien,
selon la doctrine de l’Évangile. Je suis Chrétien, non comme
un disciple des prêtres, mais comme un disciple de Jésus-Christ.
Mon Maître a peu subtilisé sur le dogme, et beaucoup insisté
sur les devoirs; il préscrivoit moins d’articles de foi que de
bonnes œuvres; . . . et il m’a dit... que celui qui aime son
frère a accompli la loi''.

In Émile and in Lettres de la Montagne he expresses himself more fully.
% --- p. 156 ---
\setcounter{footnote}{0}%
\opage{156}He believes the Gospel --- but not on external authority. And as soon as
there is talk of historical facts, we build on
human testimony, and then books and learning come to
play a part. If the essential in life could only be
known from books, what means would one then have had to come to
know it before there were books? It is impossible that God should
punish human beings for lack of learning. Therefore the Vicar has
closed all books. But the great book of nature and the book of our own heart
lie always open to us. The Gospel is the most
sublime of all books, but it is still a book, and God's law
stands written not on leaves but in the heart of man. That we
understand the Gospel rightly, and that it is its spirit we nourish our
hearts with, we notice by the peace of soul we feel.

Rousseau here anticipates Lessing's famous proposition (in the controversy
with Götze) that historical truths cannot be proof
of necessary truths. It is connected with Rousseau's
democratic spirit that he wishes to free human beings from dependence
on learning. Locke had already asserted the point of view
that hard-working people did not have strength
and time to spare for investigating the sources. The truth must
therefore be findable by a simpler way. Incipient rationalism
has a democratic character, wishes to emancipate the layman
from learned orthodoxy. Later the orthodoxy reawakening in the
period of Romanticism was guided by a similar
motive; learned rationalism met the same fate as
learned orthodoxy.

In putting away the special dogmas, Rousseau is
clear that his type of Christianity is St. James's, not St. Paul's.
And if one would deny the name of Christians to those who do not have
the developed faith of the Church, then Christ's first disciples ---
he says --- could not be called Christians either. He wishes to keep to
the character of the doctrine and of the teacher, to let it work on him. ``Nous
reconnaissons l’autorité du Jésus-Christ, parce que notre intelligence
acquiesce à ses précepts et nous en découvre la sublimité.
Elle nous dit qu’il convient aux hommes de suivre ces précepts,
mais qu’il était au dessus d’eux de les trouver''. (Lettres de la
Montagne. I).

% --- p. 157 ---
\setcounter{footnote}{0}%
\opage{157}As for the miracles, they certainly, Rousseau admits,
made a strong impression on those for whom they were intended.
That is the main thing. What does it matter, then, whether they were
real miracles or not? But now they are rather obstacles
than proofs for Christianity. Take the miracles away from the Gospel,
and the whole earth will lie at the feet of Jesus Christ!
--- One cannot assert that something is a miracle without ascribing
omniscience to oneself. For one must know all the laws of nature in order to
be able to maintain that some event or other is an exception
to them. Whether miracles happen or not, it is
in any case impossible to convince oneself that in the single
case a miracle is present.

One may very well accept a doctrine, after all, even though one does not
accept all the proofs with which people have sought to support it. The proofs
are only means, faith is the end. And faith I have, if only
one single reason convinces me! ---

It was, then, a Christianity without dogmas and miracles
that Rousseau embraced. As for worship, it was for him only
ceremony, but it had significance for him through the spiritual
community among human beings, and during his stay
in Neufchâtel he took part with great enthusiasm in Holy Communion, feeling
with heartfelt emotion: Je suis avec mes frères! ---

The books in which Rousseau expressed his conception of religion
and of Christianity were banned and burned both
in Protestant and in Catholic countries. But he has expressed thoughts
on whose validity the future both of religion in general and of Christianity
in particular will in large part depend. It
was indeed in part his thoughts, too, that in turn paved the way for
a religious renaissance. This excuses him with his ecclesiastical
opponents, but is the great count of the indictment against him
with his anticlerical opponents. But for a historical consideration
it is clear that both his appearance and the consequences
it has had are testimony that he has pointed to a
side of human mental life that is not satisfied by
intellectual and aesthetic culture but demands more thorough nourishment.
Whether the reformation of Christianity for which he spoke
could have led to the goal is another matter. And the whole
% --- p. 158 ---
\setcounter{footnote}{0}%
\opage{158}problem lies deeper than he, the son of the eighteenth century,
could see. Nevertheless it is great to have seen what he
has seen.

There is one side of the question whose significance he does not see,
--- the purely historical. He is too afraid of learning.
But however much one sublimates and rationalizes Christianity,
it is and remains a historical phenomenon,
whose origin and development it is the task of historical research
to trace. One cannot here be content to take what one
can use and leave the rest aside. And it is, as it happens,
not in St. James's type that Christianity has become a great
power in the world. It is through the apotheosis of the figure of Christ,
through the grand mythology that formed around this
figure, and through the exalted moods and deep emotions
that found nourishment through its worship, that Christianity
has become a world-historical phenomenon. It is one
of history's greatest tasks to portray the origin
and growth of this phenomenon. And the results of this research will gradually
come to be of significance for the future of the religious problem.

The individual naturally always has the right to take from a
historical formation what he needs in the conduct of his personal
life. It is characteristic of Rousseau
that what he selects from Christianity is something that he
believes all human beings will be able to rally around. He imagines
(in the letter to the Archbishop) that human beings (i.e. Jews,
Mohammedans and Christians), weary of religious quarrels,
assemble to agree on a common religion. At first
it does not succeed, because every speaker maintains his own
religion as the true one and rejects all others. But then
theologians of all religions are excluded, and this furthers
the negotiations considerably. Agreement is now easily reached on --- the
Savoyard Vicar's confession, without each giving up his
national worship for that. --- This thought recalls Jean Bodin's
Colloqvium heptaplomeres\footnote[1]{--- See \textit{Den nyere Filosofis Historie}\textsuperscript{2}. I. p. 57---61.}, which, however, had not been printed in Rousseau's
time.

% --- p. 159 ---
\setcounter{footnote}{0}%
\opage{159}It is not only the positive religions that have their
history. ``Natural'' religion too has its history, especially
its history of dogma. Rousseau's religious ideas point back
to older ``deists'' as well as to the Cartesian philosophy.
Here too he has taken what he could use. Within
Rousseau's private religion one must distinguish between what is
due to personal experience of life and what he has, naturally under
the influence of the latter, taken over from speculation and tradition.
However great or small his own existential contribution is,
it is nevertheless that which has the greatest psychological interest, and we
shall therefore return once more to his personality
in order to characterize this contribution.

\begin{center}5.\end{center}

Some years before Émile and Lettres de Montagne appeared,
David Hume (1757) had published his famous Natural History of
Religion and had thereby laid the foundation of the modern psychology of religion.
There is no trace of Rousseau's having known this work
when he wrote his books on the philosophy of religion. The contrast
between the two men's relation to religion is characteristic.
Hume stands as an observer, as a spectator at the play which
is performed in religion on the theater of mental life. But by his faculty
for psychological observation and analysis he manages to let the
various sides or tendencies in religion come more into
their own than Rousseau manages, absorbed as he is in what
stirs within himself, and without a sense for other
types than his own. Rousseau's significance rests in large
part precisely on this one-sidedness.

Hume has demonstrated a double tendency in religious consciousness.
According to the one, the distance between the religious
object and the religious subject, between gods and men, is stressed;
the divine world is raised high above the human,
so high sometimes that Hume cannot believe that a sound and
natural feeling can be nourished toward so remote an object.
According to the other tendency, religious consciousness seeks to
bring its object as close to itself as possible, to imagine
% --- p. 160 ---
\setcounter{footnote}{0}%
\opage{160}it in a perceptible way, through concrete images and as
fellow-combatant and fellow-sufferer in the course of fate.

Rousseau's conception of religion lets only the first
side come into its own. The religious feeling has for him
the character of expansion; it is an enthusiasm that rises
and wells up and leads beyond all limitation, so that at last
no thought, no image can express what is experienced. Where
Rousseau's Cartesian reminiscences do not make themselves felt,
the idea of the infinite is not formed purely intellectually, so as
thereafter to determine feeling; it is, conversely, the latter that
by its rising drives thought beyond every limitation.
This view is surely more correct than Hume's, who holds
that men do not dare to set limits to the nature and power of their deity,
but consider it safest to apply the
highest terms of praise, for which reason they try to screw their feelings up to
the height of the great words, which produces affectation (affected
ravishment and devotion).

On the other hand, the second tendency in religion comes out
very little in Rousseau. Religion leads beyond everything
finite, the soul finding vent for its longing and being freed from the
tension and the confinement that life in the finite,
especially the human, world brings with it. He hastens out
beyond experience to a world of ideals, but the great struggle
between the world of ideals and the world of experience, and the experiences
that are had in this struggle, gain no positive influence
on Rousseau's religion. It has no expression for the living through
of life's oppositions which all higher positive religions
express, especially in their worship, --- the opposition between
exaltation and abasement, between light and darkness, between life
and death. His religion has not sprung from a synthesis of
life's opposite experiences, which at once holds fast to their
reality and overcomes them in victorious resignation. He goes
off along the tangent to the circle of life, and not always as victor. ---

On the other hand, Rousseau has an eye for an element in the religious
problem that does not come into its own in Hume's psychology,
namely the concentrated energy with which religion in
% --- p. 161 ---
\setcounter{footnote}{0}%
\opage{161}its classical periods can work. It is from this point of view
that he defends fanaticism, although he grants Bayle
that it is more harmful than atheism. In spite of everything, fanaticism is,
he thinks, ``une passion grande et forte qui éleve le
cœur de l’homme, qui lui fait mépriser la mort, qui lui donne
un essor prodigieux, et qu’il ne faut que mieux diriger pour
en tirer des plus sublimes vertus'' (Émile. Profession du Vicaire
Savoyard). In connection with the proposition that is so important
in Rousseau's pedagogy: ``on n’a de prise sur les passions que
par les passions'' (a proposition in which he has Bacon and
Spinoza as predecessors), this statement shows that Rousseau
has an eye for the significance of the principle of equivalence for all spiritual
development. In the domain of the religious problem this means
that so concentrated a form of spiritual life as religion
has been in its classical times cannot disappear without
an equivalent. A more distributed or parceled-out application of
spiritual energy, the consequence of entropy in the spiritual domain,
will bring with it a decline in value. This was one of the sides of
the religious problem which the century of the Enlightenment did not
have an eye for, and which many freethinkers still overlook in their shortsightedness
and their psychological ignorance. Until
spiritual development has led to a form of life as concentrated
as the one the race has possessed in religion, a form of life
that must be developed through experience and work, without the
transcendent ways out that religion believed it could open,
--- until then the problem has not gone out of the world. It is Rousseau's
honor that he has seen this, although this point of view has not
come into its own in his personal attitude. ---

In all his frailty, with all his contradictions, all his
oscillations between the highest and the lowest, Jean Jacques has
nevertheless seen more clearly than his contemporaries in one of the most important domains of intellectual life.
By bringing the religious problem into close
connection with the problem of culture as a whole, with the great problem of the division of labor,
he has illuminated it from a central point of view.
Whatever the fate of this problem may otherwise be,
his name will not be able to be erased from its history.

% --- p. 162 ---
\setcounter{footnote}{0}%
\opage{162}\essay[Goethe and Philosophy (1907)]{GOETHE AND PHILOSOPHY.}
\streg

W\textsc{hen} Napoleon was staying at Erfurt in the year 1808, he had
Goethe summoned to him. When the poet entered, the
Emperor looked at him attentively and then said: Vous êtes un
homme! --- It must have been Goethe's manly and handsome figure,
in connection with what Napoleon knew of his works
--- he had himself, on the expedition to Egypt, immersed himself in
``Werthers Leiden,'' --- and of the course of his life, that called forth
this remark.

And Goethe was a man to whom nothing human
was foreign. He lived life fully and warmly; only for that reason is
there life in his poetry. And each time he had transformed his
experience of life into image and thought, he returned with freshness
to life. Heart and sense, imagination and thought stayed
alive in him into his high old age. By his fullness of life and his
concentration he stands as a great model, the greatest figure of modern
humanism.

In his spiritual economy as a man and as a poet
philosophical thinking too has played its part. Goethe was
not a philosopher, and he did not wish to be one. But philosophical study
and thinking were an element in his world, a tone
that came to sound along in the harmony of his life and helped
to determine its timbre.

\begin{center}I.\end{center}


In the ferment of his youth he often took refuge in
Spinoza. It was not only the clarity and the calm, in contrast
% --- p. 163 ---
\setcounter{footnote}{0}%
\opage{163}to his own restless yearning, that worked on him. He also found
in the philosopher, at that time still so misjudged, thoughts
that brought coherence into his striving to understand nature
and to live life.

He found in Spinoza a great and free view out over the world of nature
and of mind, borne by a deep assurance that
we have before us here not two different spheres, but in the end a
single great and unified connection. Beyond the contending
spiritualistic and materialistic conceptions a doctrine of unity
was here pointed to. The strictly logical-mathematical
form in which Spinoza believed he could ground and present
this doctrine, Goethe left aside. What especially appealed to him
was the point at which Spinoza himself, without being clearly
conscious of it, indicates that every completion of our view of the world
and of life takes on an artistic stamp, namely his doctrine
of the highest knowledge as a vision, a scientia intuitiva,
by which we --- after having traversed the regions of observation and
of rational thinking --- at last see the particular and
individual in its peculiarity and yet as a member of the
great whole, by whose inner forces it is borne. It was ``a
vision under the aspect of eternity'' like the one Faust so fervently
wished to attain, in which ``the quiet world lay at his feet in an
eternal evening sun''. Such a great vision was for Goethe
the highest; he knew it as an immediate divine gift
in his poetic imagination, and as an inquirer and thinker he strove
constantly up toward it through his scattered studies, even though it
was not his way to go through the strict school that Spinoza
demanded beforehand. But it was not only intellectual satisfaction
that he sought and found in Spinoza's works. In his autobiography
he emphasizes especially Spinoza's great unselfishness and
lofty resignation. Whereas most men muddle
through life by knocking off a little here and a little there, and let
their limitation and resignation be forced upon them from outside, by circumstances
and by the course of fate, he found in Spinoza's Ethics a
lofty resignation that sees once and for all the necessity of
limitation. The partial resignations are forced upon us, but
the great resignation is the personality's own work and rests
% --- p. 164 ---
\setcounter{footnote}{0}%
\opage{164}on one's having ``convinced oneself of the eternal, the necessary, the
lawful, and formed for oneself thoughts that are unshakable, indeed
are actually confirmed by the contemplation of the transitory instead
of being annulled by this contemplation''.

The reading of Spinoza appears at various points in Goethe's
life. The earliest acquaintance was made while he was still in Frankfurt,
and a visit to Jacobi played a strong part here. Later, in Weimar,
he took up Spinoza again, and several notes (in particular a
fragment brought to light by the Goethe Archive only 15 years ago,
which Goethe dictated to Frau v. Stein) testify to how absorbed
he was in Spinoza's line of thought. His association with Herder nourished
this interest. Later too, for instance during the theological
dispute between Schelling and Jacobi (1811), he took out his old copy
of Spinoza's Ethics, furnished with many underlinings,
and found confirmation of his old thought of the unity of nature
with the Godhead. A Spinozist Goethe never
was; but it speaks to the honor both of him and of Spinoza that the
stern philosopher's thoughts could be of such great value to
him as a poet, as a thinker and as a man. ---

Partly of a kindred, partly of an opposite character was the
influence Goethe received from the great philosopher who was his
older contemporary, Immanuel Kant. Here too it is not only in
one period, but repeatedly, that Goethe's attention
is turned to philosophy. It was not only his interest in everything
human, but also his own peculiar striving,
that guided him here. After his return from Italy
(1788) he became aware of the influence Kant was beginning to exert,
and which soon found its hearth in nearby Jena, where
Goethe so often went for the sake of his scientific interests.
He tried to orient himself, made excerpts from
Kant's writings, which have recently been brought to light from the Goethe Archive\footnote[1]{Cf. \textit{K. Vorländer's} essay: \textit{Goethes Verhältnis zu Kant in seiner historischen Entwickelung.} (The journal ``Kant-Studien'' I---II).},
and he had a younger scholar prepare a short survey of
the new philosophy. It was not at first the philosophical
master's chief work (Kritik der reinen Vernunft) that worked
% --- p. 165 ---
\setcounter{footnote}{0}%
\opage{165}on him; the intricate course of the inquiry and the many distinctions
confused him. On the other hand he greeted with joy
Kant's concluding work (Kritik der Urtheilskraft), in which ``the
Old Man from the King's Mountain'' (der Alte vom Königsberge), as
Goethe calls him, at last declares his many critical distinctions
between the subjective and the objective, phenomena
and the thing in itself, the natural and the moral, to be
of merely provisional significance, and, supported by the study of natural beauty,
art and genius and of organic life, points toward a
great unity in all things beyond all differences and oppositions. Here
he won Goethe's heart. ``I owe,'' he said much
later, ``to Kant's Kritik der Urtheilskraft a most happy period of my life.''
And he has aptly characterized Kant's standpoint
and procedure when he says: ``When I sought to
make use of, if not to penetrate, Kant's doctrine, it often seemed
to me as if the excellent man proceeded in a roguish
and ironic way, in that he now seemed to strive
to draw narrow limits for the faculty of knowledge, now with
a sidelong hint pointed beyond the limits he himself had drawn.''
Goethe was especially glad of Kant's suggestion of a
union of intellect and intuition; here he found his Spinozistic
favorite thought again, and the passage in Kant where this suggestion
is expressed is underlined in Goethe's copy.

When the intimate collaboration with Schiller began,
his friend's great interest in Kant's philosophy had a great effect on
Goethe. Here it is the critical fundamental thought itself, in its application
to the theory of knowledge and to aesthetics, that exerts its influence.
He later declared that through this occupation he
gained a higher concept of art and science than he had previously
commanded. Yet in relation to Schiller he always felt himself
to be ``a philosophical child of nature'', although he admitted that the
critical purgatory was beneficial to him. Goethe's natural urge aimed
at total intuition, at great, comprehensive images, but he saw
clearly that combination has value only when the oppositions have
been given freedom to assert themselves sharply. And he believed that
he could for a time calmly give himself over to critical discrimination
and to the many distinctions of analysis. In an
% --- p. 166 ---
\setcounter{footnote}{0}%
\opage{166}extraordinarily apt way he expressed himself on this in a letter to
Schiller (10 February 1798): ``Philosophy becomes ever more valuable to me,
because it teaches me daily to separate myself from myself
(mich von mir selbst zu scheiden), which I can do all the more,
as my nature --- like separated globules of quicksilver
--- easily and quickly unites again''.

When Goethe later wished to write the history of his thoughts and studies,
he names Shakespeare and Spinoza and, next
after them, Linné\footnote[1]{``Nach Shakespeare und Spinoza ist die grösste Wirkung auf mich von Linné ausgegangen, und zwar gerade durch den Widerstreit, zu welchem er mich aufforderte. Denn indem ich sein scharfes, geistreiches Absondern, seine treffenden, zweckmässigen, oft aber willkürlichen Gesetze in mich aufzunehmen vermochte, ging in meinem Innern ein Zwiespalt vor: das was er mit Gewalt aus einander zu halten vermochte, musste, \textit{nach dem innersten Bedürfniss meines Wesens,} zur Vereinigung anstreben.'' (Geschichte meines botanischen Studiums.) Linné's influence was of a kind similar to that of Kant's criticism; but then Kant in turn gave Goethe a happy time by his suggestion of a doctrine of unity.} as those to whom he owes most. But it is
clear that on his whole conception of thinking, nature and
art in his later years Kant exerted a deeply penetrating
influence\footnote[2]{See the section ``Zur Naturwissenschaft im Algemeinen'' in Goethe's works, and under it especially the fragments ``Einwirkung der neueren Philosophie'' and ``Anschauende Urtheilskraft''.}. Two of the greatest figures of modern philosophy have
thus played a strong part in the great poet's intellectual life, and
this although he, as already said, neither was nor wished to be
a philosopher.

\begin{center}II.\end{center}


Goethe had his own philosophy of life, which had developed
in him on the basis of his experience of life and as an expression
of his personality. We know it from his poems and letters
and from his conversations with Eckermann. It has drawn nourishment
and found confirmation in his philosophical reading, but it
stands as a distinct, independent contribution to philosophy in the
wider sense of the word.

% --- p. 167 ---
\setcounter{footnote}{0}%
\opage{167}He did not believe in a real opposition between the claim of the
single moment and the demand that life should form
a coherent whole. Every single state, every single
moment has infinite value, for they represent a whole
eternity! Take up the nearest task --- then you will get
through! And he believed that in this respect there was much
to learn from women. It is this that especially captivated him
when, at 75 years of age, he fell in love with the young girl he
came to know in Marienbad:

\begin{quote}
Es ist als wenn sie sagte: Stund um Stunde\\
Wird uns das Leben freundlich dargeboten.\\
Das Gestrige liess uns geringe Kunde;\\
Das Morgende, zu wissen ist’s verboten;\\
Und wenn ich je mich vor dem Abend scheute,\\
Die Sonne sank und sah noch was mich freute.\\
Drum thu’ wie ich und schaue, froh verständig,\\
Dem Augenblick ins Auge! Kein Verschieben!\\
Begegn’ ihm schnell, wohlwollend wie lebendig,\\
Im Handeln sey’s, zur Freude, sey’s dem Lieben;\\
Nur wo bist, sey alles, immer kindlich,\\
So bist du alles, bist unüberwindlich!
\end{quote}

But the very elegy in which these lines are found shows
that it was not always easy for him to fit the moment, as an eternal
element, into the whole. The moment of farewell and separation
came:
\begin{quote}
Was hilft es mir so hohe Weisheit lernen!\\
Nun bin ich fern! Die jetzige Minute\\
Was ziemt denn der!
\end{quote}

Thus even in his high old age he had to learn how
difficult it can be to carry through the great doctrine of resignation
that in his youth he had admired in Spinoza. But that
he was able to do it, --- that the sorrow could fade away, --- was due
not only to ``the divine value of tones and tears'',
which he (see the poem ``Aussöhnung'') experienced shortly after that separation,
but also to the fact that his whole view of life had
% --- p. 168 ---
\setcounter{footnote}{0}%
\opage{168}impressed upon him the great law of resignation and self-conquest.
It was his conviction that only he who has stood in the struggle
of self-conquest has a real personality. 40 years before
that episode he had written (Die Geheimnisse):

\begin{quote}
Wenn einen Menschen die Natur erhoben,\\
Ist es kein Wunder, wenn ihm viel gelingt.\\
Doch wenn ein Mann von allen Lebensproben\\
Die sauerste besteht, sich selbst bezwingt,\\
Dann kann man ihn mit Freuden andern zeigen\\
Und sagen: das ist er, das ist sein eigen!\\
Von der Gewalt, die alle Wesen bindet,\\
Befreit der Mensch sich, der sich überwindet.
\end{quote}

When Goethe here and elsewhere preaches the gospel of
self-conquest, it is nevertheless not his meaning to praise self-compulsion
and asceticism. He had learned from his Spinoza and first and
foremost from his own experience that an impulse can only be checked
by another impulse. The struggle against hindrance and resistance must be
waged in a positive way. Evil must be displaced by
putting something good in its place. Great spiritual interests must
fill the place from which the momentary claim had to give way. And
so Goethe also took his stand toward the bad and the wrong
that he saw around him: ``He who would truly do good work must not
scold; he should not trouble himself at all about the wrong, but
only always do the good. For what matters is not that something
is torn down, but that something is built up in which
mankind can feel a pure joy''. And when he had now and then
spoken in bitter words of the attacks to which he was subjected in his later
days, he broke himself off: ``Not another word
about the bad business, lest I myself become unreasonable
in combating the unreasonable!'' ---

One of the positive ways Goethe struck out on, when he
had to work his way beyond a sorrow or a disharmony, was
that of transforming his experience into poetry. It was not his
intention to forget, still less to enjoy, in this
% --- p. 169 ---
\setcounter{footnote}{0}%
\opage{169}way\footnote[1]{Cf. my essay on Goethe in ``Mindre Arbejder'' II p. 231---234.}. But in the poetic form the experience was seen ``under
the aspect of eternity'' --- in its peculiarity and yet
as a single member in the great connection of life. In Werther, in
Tasso, in Faust and Mephistopheles, in Wahlverwandschaften, in his
lyric poems he has carried out such a transformation, into the clear
form of art, of what during the experience filled him with
passionate infatuation, with sorrow, with bitterness or with
remorse. Of ``Werther'' he could say in his old age
that it was written with his heart's blood, and that he was afraid
to read it again. His poetry was not a road by
which he withdrew from life, but a part of his own
personal conduct of life, a way in which he worked
what he had lived into his personality, so that it could thereby
also gain significance for others.

It was far from being the case that he carried out this transformation from life
to poetry with full consciousness, let alone with calculation.
Where it happened perfectly, it took place unconsciously and in silence.
The great images formed themselves in him without his suspecting
it, and then suddenly thrust themselves forward before him. He therefore could
not, like Schiller, talk with others beforehand about his
poetic plans. ``Ich trug Alles still mit mir herum.'' The details
he could work out afterwards with clear consciousness.
But the main thing was ``the surprising gift'', the outcome of
the ``daemonic'' power that he found not only in the imagination,
but in all involuntary modes of action, in the whole of nature. He
therefore could not abide the word ``composition'', because it suggested
an external and mechanical putting together. And since what matters
above all is the quiet growth, he advises not to force
artistic work, not to ``whip oneself up in the bad time'',
but to wait for the good times. In his art too he found
use for his philosophy of life.

\begin{center}III.\end{center}


When the work has come into being through quiet growth and conscious
after-work, it stands as a total image that can neither
% --- p. 170 ---
\setcounter{footnote}{0}%
\opage{170}be expressed in a general idea nor be explained from definite,
actual experiences and circumstances. Life, imagination and thought have
worked together in so intimate a way that neither a speculative-idealistic
nor a psychological-realistic aesthetics can
unravel the living web that is the fruit of this collaboration.
What ``the daemonic'' has joined together, criticism cannot and shall not
put asunder.

Again and again Goethe declares that every work of art is
incommensurable and irrational. ``Man kann über Shakespeare
gar nicht reden; es ist Alles unzulänglich.'' ``Je inkommensurabler
und für den Verstand unfasslicher eine poetische Produktion
ist, desto besser.'' He firmly rebuffed those who asked
what the ``idea'' in Tasso or Faust was. He could not say
it. It was life in its fullness and in its oppositions that he
wished to present, on the basis of what he himself had lived, and
it could not be strung on a single thread. ``Der Faust ist doch etwas
Inkommensurables!'' And of ``Wilhelm Meister'' he says:
``Es gehört dieses Werk zu den inkalkulablesten Produkten,
wozu mir fast selbst der Schlüssel fehlt!''

Nor do the definite facts that were the occasion of the
work or were used as motifs in it have decisive significance
for the understanding of it. He was amused that people
wanted to know which town on the Rhine was meant in ``Hermann
und Dorothea'', and that King Ludwig of Bavaria wished for information
about the fact that lay at the basis of the ``Roman Elegies''.
``People seldom consider that the poet is for the most part able
to create something good out of slight occasions.''

The conception he thus asserted of the work of art
as a whole that could be grasped neither from an
idea nor from its elements agrees entirely with the one he
asserted in his conception of nature and in his inquiry.

One should, he maintains, keep as long as possible to
immediate intuition and keep analysis and judgments away.
``Anschauen'' and ``Urtheilen'' form sharp oppositions in
Goethe. He speaks of ``the quiet fruitfulness of such impressions
as one takes in and enjoys without dissolving judgment'' (ohne
zersplitterndes Urtheil). During his Italian journey he writes:
% --- p. 171 ---
\setcounter{footnote}{0}%
\opage{171}``I merely keep my eyes always open and impress the objects
firmly upon myself. I would rather not judge at all, if it
were only possible''.

This principle he followed in his inquiry. Whether
he is studying colors or plant and animal forms, his concern is
not to find laws and mechanical connection, but
to find typical intuitions, which he could find again everywhere,
only in different nuances. It was a procedure that
suited the comparative study of plant and
animal forms better than a physical theory of colors. When a typical
phenomenon, what Goethe called a primal phenomenon, had been found,
the task was to show its nuances. Nature was for
him a play of forms, and the task was to see the fundamental form
that expresses itself in this interplay. But he is indignant with those
who want to go further back and seek an explanation of the primal phenomenon
itself. At this one must stop; it is the limit of
our vision; from it one must proceed by way of analogy
to the multiplicity of the special phenomena, but behind it one
cannot get. For behind it stands the Godhead itself, the active force
of nature. Goethe's theory of colors (to keep to that)
has a more aesthetic than scientific character, but he
trusted it to such a degree that he engaged in a vehement
polemic against modern physics, as it had developed
from Newton's time on. The conception of nature should, according to
Goethe, come into being in us just like the work of art, through the involuntary
emergence of typical images, only that the natural scientist finds these
images, the primal phenomena in the various domains, by observation
and comparison. Thus in his theory of colors Goethe wished
to explain all colors as different mixtures of light and
darkness; his primal phenomenon here was the yellow or blue shadows
that are seen when the sun shines on faintly darkened media (with
a lighter or darker background). A further explanation of this
primal phenomenon itself he held to be impossible. Here, he believed, we must
resign ourselves, because we stand at the limit of human capacity.
It was his fear of letting go of intuition that called forth
this hasty resignation. It was the poet in him that
held the inquirer back.

% --- p. 172 ---
\setcounter{footnote}{0}%
\opage{172}Goethe knew very well that the work of thought consists both
in distinguishing and in combining, both in analysis and synthesis.
He himself compares this double procedure with
breathing out and breathing in, which are both equally necessary. And
he maintains that one understands nature only when one distinguishes
sufficiently to be able to combine, and combines sufficiently
to be able to distinguish. But he could not, as exact
natural science and critical philosophy demanded, regard every
image, the ``primal phenomenon'' too, as provisional and insufficient,
or as a mere example. He clung to intuition
and harbored distrust of concepts and judgments. He
wished to explain everything by analogy with his primal phenomena and did
not see that these themselves are phenomena like other phenomena and
can have no privilege over them.

The reason was not that he held a completed knowledge of nature,
an absolute system, to be possible. Nature is for him
incommensurable like the work of art. Nature has no system;
it moves from an unknown center to an unknowable
limit. And man does not exist in order to solve the problems of the world,
but to seek where the problematic begins,
and then to keep within the comprehensible! --- When a student
laid before Goethe questions that lay beyond this limit,
Goethe wrote in his album: ``Der liebe Gott hat die Nüsse
wohl geschaffen; er hat sie aber nicht auch geknackt''.

When Goethe has thus come to the limit of knowledge,
he turns aside and returns to his proper domain,
the art of poetry. Science can never make thought and
experience wholly cover each other, but ends in the eternal
strife between what is intuited and what is thought; only art speaks
a language that expresses nature, that great weaver, whose
work is not mechanically joined together from scattered threads, but
where whole and parts are in eternal unity.

\begin{center}IV.\end{center}


In Goethe's religious conception we find fundamental features that
correspond to his conception of art, of science and of nature.


% --- p. 173 ---
\setcounter{footnote}{0}%
\opage{173}He believed in a Godhead that worked in nature from within,
through nature's own forms and forces. A God who
pushes against nature from outside to set it going could not
be a true God. If one asks how he conceives
the Godhead, he answers that the highest truth, which is
one with the divine, cannot be known directly, --- ``only
in reflection, example, symbol''. But he is convinced that
``the spirit's highest fiery yearning has enough in image and likeness''.
And he himself could not be content with a single symbol, but
had to have changing symbols, according to the urge and capacity
in him that especially asserted itself in the particular state.
He writes to his theistic friend Jacobi (6 January
1813): ``I for my part cannot, on account of the manifold
directions of my nature, be content with a single way of thinking:
as a poet and artist I am a polytheist, as a naturalist
I am a pantheist, and the one as decidedly as the other. And if I
then need a God for my personality as an ethical being,
that too has no doubt been provided for''. A similar statement
is found on a slip among his papers. A dogmatic formulation
and working together of his religious ideas he
did not feel any urge for. For this reason he was attacked as irreligious.
``I believed in God and nature and in the victory of the good over
the bad; but that was not enough for the pious souls.''

He believed in an immortality for all those in whom
a spiritual energy, a firm type, had been developed; such a one,
he believed, nature cannot do without. But there are degrees in
this respect, and we are not immortal in the same way.
Only for him who constantly acts is nature obliged to provide
a new form of existence when the earlier one is dissolved. But then
one also has enough to do with one's striving and acting in the
present life and can leave the future life to itself,
all the more since no definite ideas of its nature
are possible. Goethe shunned those who wallowed in thoughts of immortality
and were busy examining the faith of others in
this respect, and when some importunate ladies examined
him on the matter, he declared that he would in any case
hope not to meet those who had already
% --- p. 174 ---
\setcounter{footnote}{0}%
\opage{174}been busy with immortality down here, for it would
be unbearable when they came rushing up and cried:
``Haben wir nicht recht gehabt? Haben wir es nicht vorhergesagt
?'' ---

More markedly than in the form of ideas, Goethe's
religiousness expresses itself in feeling. This already comes out in Faust's
confession of faith, when he is catechized by Margaret: ``Gefühl ist
Alles! Name ist Schall und Rauch!'' And as the worshiper of nature
Goethe was, --- as one who had a sense for the mysterious
(daemonic) in nature and in man, as well as for
the great and life-furthering forces in the world, for man's
connection with nature, as well as for the limits of mankind, --- his feeling
toward existence had to take on the character
of reverence. In the work of his old age ``Wilhelm Meisters
Wanderjahre'' this feeling is described in such a way that we are
reminded in an interesting manner of what he admired in his youth
in Spinoza.

There is one thing, it says, that no one brings with him into the world,
and on which nevertheless everything depends if man is in all respects
to be man: it is reverence. The natural man
is tossed back and forth between fear and freedom, and only
with difficulty does he work his way out beyond this swell.
Or rather: he does not work his way out of it, but it is a
higher sense, lying in his nature, that leads him to
reverence, which at once lets him see his dependence
and his independence. Herein consists the worth and significance of all genuine religions.
--- Reverence stands here as akin to the
lofty resignation Goethe had come to know in Spinoza.
But it is a stronger and deeper feeling that is described here;
in the meantime Goethe has lived life and seen nature. Reverence
has three different directions: upward toward the great
powers of nature; --- downward toward the beings that stand below us,
or that suffer and are crushed; --- forward toward the other men,
our fellows in the battle of life. There are three religions that spring
from these directions: the first is the heathen, the
second the Christian, the third the philosophical. The mature
% --- p. 175 ---
\setcounter{footnote}{0}%
\opage{175}man unites all three in reverence for himself, which
leads to finding his right place in existence.

Goethe has here characterized the fundamental mood that, throughout his
life, animated his poetry and his inquiry and his whole
conduct. When he ended his greatest poem with the thought that he who constantly strives
can be redeemed, reverence lies at the basis here too;
for a world order that makes high and great striving possible itself
takes on the character of sublimity, and reverence must be the fundamental mood
that the thought of it awakens.

% --- p. 176 ---
\setcounter{footnote}{0}%
\opage{176}\essay[At the Schiller Festival (1909)]{AT THE SCHILLER COMMEMORATION}
\undertitel{25 NOVEMBER 1909.}
\streg

T\textsc{he} man\footnote[1]{Georg Brandes.} who was the one rightly called to open this
festival is held captive by illness. If, despite much misgiving,
I undertook to take his place, it was because
I trusted that I did have one qualification for carrying out
the task: the great admiration and the deep gratitude
toward the great figure whose memory we are to celebrate this evening.
I wish especially to speak in the name of the many for whom Schiller's
lyric poems, weighty in thought and shaped in powerful rhythms,
have been, from our young years on through the years of manhood, a
constant source of courage for life, a constant spur to new and higher
striving up above everything that can dull or drag down.
The manly force and the great pathos will at all times make
his most significant poems some of the best examples of
the poetry of life, a poetry that does not lead into a dream world
or in romantic fashion veil the disharmonies, but gives
expression to the feeling of life of one who stands in the midst of life's work
and toil, striving despite everything to hold fast to a great goal.
And then his tragedies, in which we hear the heavy tread of mighty
fate, --- where clearly drawn, if not picturesquely
nuanced, figures come to meet us, --- and where from
their lips sound striking and energetic words about the conditions of life,
about the demands of life and about the battle of life, words that have become winged
words even for many who do not know their
% --- p. 177 ---
\setcounter{footnote}{0}%
\opage{177}origin! For my part Schiller as a thinker too has
been of great interest. In the great Renaissance in German intellectual life
toward the end of the eighteenth century --- in many
respects the most significant spiritual movement since
the days of the Italian Renaissance --- Schiller's figure stands
between the two giants Goethe and Kant, akin to both
and a characteristic intermediate link between them. Not that he
could be, or himself believed he could be, set as their equal.
Thinking was, especially in a certain period, an essential element
in his spiritual work, but only an element. And as for
poetry, he saw clearly Goethe's unique worth.
``Gegen Goethen bin und bleib’ ich ein poetischer Lump'' (letter
to Körner 27 June 1796). But he was not unworthy to be
named together with the two spiritual heroes. As his friendship
procured Goethe a new spiritual youth, so
his conviction of the significant place of the artistic
in human life has helped to harmonize and humanize what was
sharp and hard in Kant's ethical ideas. --- And now finally the
man Schiller himself, as we know him through the course of his life,
from the fermenting youth in the period of Storm and Stress, through
the period of work with its groping attempts, to ever increasing
clarity in understanding and firmness in artistic activity, despite
want and illness always striving forward toward the high ideals,
which indeed stood before him in different shapes at different times,
but always signified for him the great goal for which alone it
was worth living. A peculiar luster lies over the
beautiful life together with Goethe, in which each of them stands
at once as helper and as helped forward. Outside the theater
in Weimar their proud figures stand, shoulder to shoulder,
looking upward, up toward the goal toward which each by his own ways strove,
and which ennobled life for them both. --- ---

It would not be my province to give you a thorough
characterization of Schiller as an artist and of his place in the history
of literature. Such a characterization you have missed
because the man did not come to speak who alone could speak
with authority on it.

What I shall try, in all brevity, is to show Schiller's own
% --- p. 178 ---
\setcounter{footnote}{0}%
\opage{178}conception of the relation between art and life. When Schiller
philosophized, it was especially the significance of art for human life
that occupied him.

He had a great idea of what art had
been and is for the race and in the life of the individual. The raw and
unruly forces of nature and impulses of nature, which the stern law of duty
was unable to bend without crushing them, indeed which it
perhaps even provoked with its severity, --- shape and limit
themselves involuntarily when they stand before the works of the plastic arts,
or when the harmonies of tones sound for them. And art itself
springs involuntarily from man's nature. Other living
beings surpass man in strength and in instincts; but
art is man's distinguishing mark:
Die Kunst, o Mensch, hast du allein!

Art, then, is the great educator; it is the first
form in which man comes to know spiritual enjoyment;
it leads up from rawness or asceticism to complete and harmonious
human nature, where no single element rules at the expense of the other
elements, and where the high demand of the ideal can
be met with the enthusiasm that art's images of the great and
sublime or of the gentle and graceful are able to awaken.
But it is not merely an educator. Only by becoming art
does a human work reach its perfection, whether it be
the work of thought or of the hand. Art is both
beginning and end, both means and goal. The dignity of mankind
lies in the hands of the artists; they are appointed to
assert and preserve it!

These thoughts Schiller expressed in the great poem ``Die Künstler'',
which was written in his thirtieth year. Some years later
he took up this line of thought again in a longer treatise: ``Briefe
über die ästhetische Erziehung des Menschen''. It is here
especially the significance of art for modern man that he
discusses. The great dark sides of our culture are connected
with something that, seen from another side, is a great advance,
and that in any case has come as a necessity,
--- the division of labor in the material and in the spiritual
% --- p. 179 ---
\setcounter{footnote}{0}%
\opage{179}domain. The individuals become more and more members of a great machine.
Each gets his own, often insignificant and monotonous,
task to perform, a task that lays claim to only a
single side of his nature, perhaps requires only a single, simple
manipulation. Whole and full living, in which all sides of man's
being come to work together, becomes more and more
difficult to attain. And even those who can work in the world
of the spirit suffer under the division. They lead one life in their
intellect, another in their heart, a third in their imagination.
But in art --- both in those who are influenced by it and
in those who create it --- all the powers of the soul work together.
What shows itself in childlike form in play, a use of the powers
merely for the sake of using them, without any external end, without fear
and without desire, --- appears on a grand scale in art, when
the mind is taken up by the images in which the fates of life are presented
in their changing forms, showing us now the sublime,
now the graceful. It is no external need that leads to
such presentation; it springs from the inner urge of the mind.
Give man moments in which he can rise from the scattered and isolated
life to dwelling on a great and comprehensive
image of life, --- and from such moments there will radiate
after-effects whose ennobling influence is felt when he again
returns to his place of work! The very fact that he feels
joy and enthusiasm over something that is not a sensuous reality
is a testimony of freedom and self-activity, of independence
of narrow interests and impulses. --- ---

Such, roughly, were Schiller's thoughts on the significance of art
for life. The artist was for him a prophet. And this
prophethood he celebrated in poems and in treatises. It was
in the period from his thirtieth to his thirty-sixth year that he
especially brooded on this question. It was characteristic
of him that he felt this urge to assure himself that what he
sacrificed his life to was also a great and indispensable human
good. It was his heavy nature, ever struggling with inner and outer difficulties,
that thus for a time let him
brood over art instead of working in art. He himself felt
the self-contradiction in it: he, the poet, who prized his
% --- p. 180 ---
\setcounter{footnote}{0}%
\opage{180}art as the true, full life, occupied himself with speculation,
this most one-sided occupation! He even writes
to Goethe (7 January 1795): ``So much is certain, the poet is
the only true \textit{man,} and the best philosopher is only a
caricature beside him,'' --- a remark every philosopher will approve,
\textit{when} it is a Schiller who says it to a Goethe.

But in Schiller's artistic work itself too we find
this urge to brooding and reflection in contrast to
full self-surrender. It was not, says Goethe, in Schiller's nature to
work involuntarily, as if by instinct; a certain deliberateness was
always present. The particular traits and images were often brought into
the world with effort and as a result of deliberation. They
did not come, as Goethe tells of his good inspirations, ``like
free children of God, who suddenly stood there and said: Here we are!''
His work was severe, though it was his highest life. He
could be far down in doubt, darkness and spiritual trial as long as
the preparatory work went on. He expresses himself very strongly
on this in a letter to Körner (25 May 1792): ``It often happens
to me that I am ashamed of the way in which my
works, even the most successful, come into being. It is often said that
the poet must be fully absorbed in his subject while he
writes. But in my case it is often a single and
not always an important side of the subject that makes me work on
it, and only during the work does the one idea then separate itself
from the other... But how is it possible that from
so unpoetic a procedure something excellent can nevertheless arise?''
Schiller finds the solution of this question in the fact
that it is often an indeterminate urge to give vent to feelings striving
forth that leads to the poetic works. As long as
this urge had not yet found means and ways, there could
be heavy times. But according to what his friend \textsc{Wilhelm v.}
\textsc{Humboldt}, who in some of his most important years of breakthrough was
with him daily, has said, everything that
made the work hard for him served to increase the perfection of
the matured work. --- Unfortunately it is not all hard work
of which this can be said.

I cannot in this connection refrain from thinking
% --- p. 181 ---
\setcounter{footnote}{0}%
\opage{181}of the Song of the Bell, which we are later to hear interpreted here.
As so often in Schiller, there are released from the images he
gives other, more or less symbolic images, which consciously
or unconsciously may also have stirred in the poet
during his work. In the progressing work of
the casting there is not only occasion --- often through beautiful and inspired
hints in the situation --- to depict the course of human life
from cradle to grave, through happiness and misfortune,
through war and peace, but we also hear the master's moods
during the work. Will what is now being prepared here in the darkness
of the casting pit really become a shining bell, whose
tones can be heard far and wide from up in the tower?

\begin{quote}
Wird’s auch schön zu Tage kommen,\\
Dass es Fleiss und Kunst vergilt?\\
Wenn der Guss misslang?\\
Wenn die Form zersprang?\\
Ach! vielleicht, indem wir hoffen,\\
Hat uns Unheil schon getroffen.
\end{quote}

In this anxiety speaks Schiller's own experience of his artistic
work. From his correspondence with Goethe we know that
he had trouble getting this poem finished. --- Both
as a man and as an artist his life was full of struggling
oppositions. But then he also knew the master's
joy when the work succeeded and could by its tones interpret
and guide the course of life. It was his conviction of
the high worth of art that kept him up in the dark times.
He sacrificed, as Goethe has said of him, his very life
to bringing forth the great images of life. ---

When we here this evening celebrate Schiller's memory, it is of interest
to recall that from Denmark a helping
hand was extended to him at a time when economic pressure threatened
to halt his poetic activity, precisely in the years when he
was, through brooding, to work his way forward to the last great
period of his poetic life, the one in which Wallenstein, Maria
Stuart and Wilhelm Tell came into being. The treatise I mentioned before,
``Briefe über die ästhetische Erziehung des Menschen'',
% --- p. 182 ---
\setcounter{footnote}{0}%
\opage{182}is dedicated to Duke Frederik Kristian of Augustenborg, brother-in-law
of Crown Prince Frederik and Danish minister of education. The dedication
was a thanks for a high-hearted support that the duke together
with another Danish minister, Schimmelmann, had given the
sick and needy poet. It was the greatest Danish
poet of the time, Baggesen, whose enthusiasm for Schiller had communicated
itself to the two magnates, who were also his patrons. The summer
before their gift to Schiller, he had been
dangerously ill, and the rumor spread that he was dead. His
Danish admirers, not least Baggesen, were seized by the deepest
grief, and at Schimmelmann's estate in Hellebæk a memorial celebration
was held, at which Baggesen recited the Hymn to
Joy, accompanied by music. This token of sympathy too
moved Schiller deeply, and it is kind thoughts that he in his last
years sent up here. His influence on Danish literature
has surely been great; I need only mention names like
Oehlenschläger and Grundtvig; the latter names not only
his poetry but also his view of life as something that
has had significance for his development. --- And this assembly
here this evening is a testimony that he still exerts his
influence among us and finds understanding and admiration.

We have, then, reason to be glad and proud that
from here an indispensable help was given him.

But proud above all may the people be to whose best sons
he belongs. No poet has, as he has, become the German people's
property. If it can be disputed that he should stand in the first
rank of the great spirits of the poetic world, the rank in which
only four or five places can be filled, there is no
dispute that he has brought high poetry into human minds and
human homes where works of poetry otherwise do not reach. And
this is connected with the relation of his art to life. His
art wished to carry the dark image of reality over into its own
bright realm, the seriousness of life over into its cheerful world, so that
reality and life could in turn be strengthened and elevated by the
spiritual forces that art gives occasion to exercise.

When the ``Lied von der Glocke'' was performed at Weimar's theater
a few months after Schiller's death, Goethe had written an
% --- p. 183 ---
\setcounter{footnote}{0}%
\opage{183}epilogue to it, in which it sounds again and again with pride:
``Er war unser!'' This applies here especially to Weimar itself, the
little town on the Ilm, --- ``wie Bethlehem in Juda klein und
gross'', --- where Schiller had found shelter in his last years.
Goethe recalls how people here had seen him in his
daily life, his severe work, his constant struggle with
illness, but above all had been witness to how his spirit
had steadily advanced:
\begin{quote}
Indessen schritt sein Geist gewaltig fort\\
Ins Ewige des Wahren, Guten, Schönen,\\
Und hinter ihm, im wesenlosen Scheine,\\
Lag, was uns alle bändigt, das Gemeine.
\end{quote}

A more beautiful memorial, a more beautiful characterization than Goethe
has given in these lines cannot be imagined. --- ---

``Er war unser!'' --- Beyond the little town's, and beyond
the German land's borders this word can be repeated. He was
and is ours, as truly and as long as the spiritual treasures he
has raised, and the great conception of life and of art that
he has expressed in poem and in thought, win understanding and appropriation
among us.

High above life's confused paths, above the noise of egoisms,
above the much that can ensnare and crush, his
clear, manly spirit points constantly, through an art that from the very
midst of life leads upward toward its heights. Never may the
time come, not in our fatherland either, when there is not sense and
enthusiasm for a hero of art, of thought, of life like
Friedrich Schiller!

% --- p. 184 ---
\setcounter{footnote}{0}%
\opage{184}\essay[Napoleon and Philosophy (1906)]{NAPOLEON AND PHILOSOPHY.}
\streg

I\textsc{f} one wonders what Napoleon and philosophy could have to do with each other, one must consider that
power, whether it appears as the power of the state or the power of the church, as absolutism, as aristocracy or as
democracy, cannot do without thought, without ideas, to guide the use of power and, first and foremost, to justify
the right of power. And on the other hand thought cannot in the long run do without power. What is fostered in the
secret chamber and is tended and developed in the smaller circles will in the end come to set those in power tasks
they cannot get around. But besides, those in power themselves, and not least the man in power we have before us
here, are a rich subject for philosophy, through their character, their genius, their conduct and their fate.

Napoleon's relation to the philosophy of his time is in any case a not uninteresting chapter of the history of
philosophy.

Napoleon, the son of the eighteenth century, came to power by the help of the ideas of the Revolution. In their
service he fought his first battles and won his first victories. And from the first he saw the power he could win as
a servant of ideas and of scientific inquiry. When, after the Peace of Campo Formio, he was received into the
National Institute, France's academy of sciences, he said: ``The true conquests, the only ones that give rise to no
sorrow, are those that are made over ignorance. The true power of the French Republic must henceforth consist in not
permitting a single thought to arise that does not belong to it.''

% --- p. 185 ---
\setcounter{footnote}{0}%
\opage{185}There was accordingly, for the time being, mutual sympathy between him and the French philosophers of
the time.

The philosophers had been persecuted under Robespierre's reign of terror. Power had here shown itself suspicious of
thought. Destutt de Tracy, for instance, had been imprisoned. He continued his work in captivity, from which he was
released only at Robespierre's fall on the ninth of Thermidor (1794). Cabanis, next to Destutt de Tracy the most
eminent of the French philosophers of the time, had been spared only because of the respect his love of humanity
inspired; out in Auteuil, where he lived, he practiced medicine without payment among the poor.

Napoleon visited Cabanis in Auteuil and sought in vain to get Destutt de Tracy to come with him to Egypt. When he
came home, the philosophers sympathized with his action against the Directory. People hoped that he would become
another Washington, and greeted the eighteenth Brumaire with joy. Under the new constitution Destutt de Tracy
became a senator.

The group of philosophers in question here has received the common name of Ideologists. The name is owed to Destutt
de Tracy's chief work \textit{Éléments d’Idéologie,} which appeared in 1801, an exposition of what we would now call
psychology and theory of knowledge, built on a series of more special treatises that had been published in the
Transactions of the National Institute in the preceding years. Ideology properly means merely the doctrine of
representations or ideas. It aims to give an analysis of the life of thought and of mental life. A careful
distinction is drawn between facts, hypotheses and speculations. Psychology is brought into close connection with
physiology. Characteristic of the fundamental psychological conception is --- in critical contrast to Condillac's
attempt to reduce mental life to a mere product of passively received sensations --- the weight laid on activity in
mental life, on the will, in connection with the original tendencies to movement, as they already appear in the
movements of the fetus in the mother's womb. Great weight is laid on resistance as the mark of reality, and on our
capacity to offer resistance to what is not ourselves. ``To will and to offer
% --- p. 186 ---
\setcounter{footnote}{0}%
\opage{186}resistance,'' says Destutt de Tracy, ``that is in truth to be!''

Ideology is a forerunner of the more recent empirical psychology. It has exercised no slight influence on later
inquiry; philosophers and psychologists such as Schopenhauer, Bain and Taine were
% print: forsatte (misprint)
influenced by it and continued its work. Ste Beuve thought of writing the history of this school, but did not carry
out this plan. A learned and thorough account of it has been given by F. Picavet (Les Idéologues Paris 1891), and
this, besides the study of the works of Destutt de Tracy and Cabanis, underlies this sketch.

One sees that the word ideology designates a sober inquiry, far from speculation and utopia. But through a great
collision between thought and power an entirely different meaning was forced on the word, which --- so great was the
influence of power --- it has kept in common usage! ---

After having saved France, Napoleon did not let go of the grip his iron hand had got on power. From now on power
gradually became everything to him, although in the first period of his rule he used it to carry through several of
the thoughts of the Revolution. The philosophers' criticism began to become unpleasant to him as he proceeded to
reactionary measures. The Concordat in particular met with opposition from their side. It should be well noted that
it was in the first instance the political opposition that embittered him; but he traced it back to philosophical
criticism as its ultimate cause. From now on, then, the word ideologist became in his mouth a term of abuse,
synonymous with a nebulous and pedantic metaphysician; ideology became one with speculative cobweb-spinning. But the
real meaning of the word ideology in his mouth Lanfrey has probably rightly stated thus: the sum of all the ideas
that were hateful to Napoleon.

The reason for the opposition he met from the side of the ``Ideologists'' lay in his despotic measures. The freedom
of the press was abolished. The Concordat was established. Special courts were instituted for political crimes. From
dislike of free inquiry, instruction in history and jurisprudence was abolished in the higher schools; philosophical
instruction was restricted, and the
% --- p. 187 ---
\setcounter{footnote}{0}%
\opage{187}philosophical section of the National Institute was abolished. Books too were persecuted; Mme de Staël's
``Allemagne'' went to the paper mill, and Destutt de Tracy's commentary on Montesquieu had to be printed in America.

Thus he defied the ideas on which the right of his power rested. He forgot, as so many in power constantly forget
it, what Rousseau, the object of his youthful admiration, had declared: fait ne fait pas droit! Power, which had
been a means, now became his only end.

The opposition in France he could subdue. Yet here too it became his nemesis that he could as little tolerate proud
souls around him as he could tolerate free ideas. And it became his downfall that he lacked respect for every kind
of spiritual movement. The religious movement that lay behind the opposition of the Catholics, and the national
movement that lay behind the opposition of the conquered peoples, were for him a kind of ideology that he believed
he could defy. Yet it was the German ``Ideologists'' who in high degree helped to awaken the enthusiasm that led to
his fall. Fichte and Schleiermacher,
% print: Humholdt (misprint)
Wilhelm von Humboldt and Henrik Steffens showed by their national activity that ideas are, after all, a power.

When the catastrophe then came (1814), it was precisely Destutt de Tracy who in the French Senate made the motion
for his deposition. --- ---


Let us hear about the fate of the two men later on.

After his death Napoleon came to honor again. In the time of reaction that followed he stood as the hero of the new
age, a symbol of the revolutionary thoughts, and his figure towered in the popular consciousness high above the
small men who now governed Europe.

Goethe has expressed this mood in a poem. The Devil appears before God's throne with a long register of the sins of
the fallen greatness. But his long speech bores the Deity, and the following words cut short the accusation:
\begin{quote}
Wiederhol’s nicht vor göttlichen Ohren!\\
Du sprichst wie die deutschen Professoren.\\
Wir wissen Alles, mach es kurz! --- ---\\
Getraust du dich ihn anzugreifen,\\
So magst du ihn nach der Hölle schleifen! ---
\end{quote}

% --- p. 188 ---
\setcounter{footnote}{0}%
\opage{188}In the July Days of the year 1830, when the reaction was overthrown, an old gentleman with silk
stockings, a long cane and a green shade over his eyes was seen moving among the barricades. It was Destutt de Tracy,
who despite the commands of Robespierre, Napoleon and the Bourbons had not lost his faith in the principles of 1789,
which he now heard proclaimed again, although he was not satisfied with Louis Philippe's conduct. He lived, however,
to see the philosophical section of the Institut de France restored before his death, which occurred in the year
1836.

% --- p. 189 ---
\setcounter{footnote}{0}%
\opage{189}\essay[Kierkegaard and Nietzsche (1911)]{KIERKEGAARD AND NIETZSCHE.}
\streg

I\textsc{t} might seem a strange thought to place side by side these two men, who on decisive points stand in the
sharpest opposition to each other.

The one asserted Christianity in its strictest form, sharpened orthodoxy into paradox and presented it in
expressions that set its doctrine in absolute and hostile opposition to natural, human ways of thinking. In
particular he set Christian ethics in the sharpest conflict with human demands and endeavors. The other was a radical
freethinker with so strong an antipathy to Christianity and to all views of life related to it that he could not
find words strong enough for his aversion. On the other hand he asserted with ecstatic enthusiasm the health and
strength of human nature and believed in its unlimited possibilities of development.

For the one the ideal lay in the past. Only a single time in history had the divine light shone into the darkness
of existence, and the relation to this point of light was decisive for the life of the human race and for the
personal value of each individual human being. The other held just as energetically to an ideal of the future,
proclaimed the coming of the superman and, in connection with this, a revaluation of all values.

In their outward circumstances too the two men form a remarkable contrast. Kierkegaard wandered about in the narrow
streets of Copenhagen, and what occupied the small, gossiping town penetrated all too much into the cloistered
solitude in which he lived. Apart from a few stays for study in Berlin and his excursions
% --- p. 190 ---
\setcounter{footnote}{0}%
\opage{190}to the forests of North Zealand, he walked back and forth in his rooms, now and then stopping at his desk
to write down his thoughts. For the rest he immersed himself in strictly ascetic books of devotion. Nietzsche had
spent his youth in the study of classical antiquity and had for a time been a teacher in this subject; later he
spent his time on the Elysian heights of the Engadine or on the Riviera. His thoughts were born out in open nature.

But precisely these extreme contrasts could be evidence of a spiritual kinship. Rivers that flow in opposite
directions may for all that have their source on the same mountain. And the philosopher is more interested in the
source than in the course, --- more in the deepest-lying demands than in the different ways in which their
satisfaction is sought, --- more in the leading thoughts than in the dogmas or symbols by which they are expressed,
--- more in the experiences of life than in their interpretation.


What Kierkegaard and Nietzsche have in common is that they both look far out beyond the usual, bounded, rational
tracks of life. They assert what is innermost and original in man, something that (provisionally at least) cannot be
analyzed and formulated in exact fashion, but which can nevertheless possess the highest power and the highest
value. They could not find that modern culture, with all its science, art and industry, satisfied the urge of
personality, and they were led besides, though from entirely opposite points of view, into sharp opposition to the
prevailing religious and ecclesiastical system. They are symptoms of something that is stirring in the age, and that
has here found a volcanic expression.

There can be no question of any historical connection between them. Kierkegaard (b. 1813, d. 1855) was in
Nietzsche's time (b. 1844, insane 1889, d. 1900) very little known in Germany. --- On the other hand I find very
significant analogies between them on three points.

The first point concerns the relation to truth.

Kierkegaard asserted the proposition that subjectivity is truth. What matters is not so much what the truth you
embrace contains as that you have won or appropriated it through your own personal work and live in it
% --- p. 191 ---
\setcounter{footnote}{0}%
\opage{191}with all the inwardness of your soul! Two men go up to pray. The one lives in a heathen land and prays to
an idol; but he puts his whole soul into his prayer and gives himself up in it: then in reality he prays to the true
God! The other lives in a Christian land and prays to the true God; but he prays dully and out of habit; his
personality is not in the prayer: then in reality he prays to an idol!

Kierkegaard's emphasis on personality led him out beyond all scientific thinking. What interested him as a thinker
were the boundary problems. His thought had an irresistible urge to follow the line to its end, to seek out what
cannot be thought but lies out beyond the limits of thought. Hereby he came into conflict with the doctrine of the
church, which it was nevertheless precisely his intention to hold fast to, but which presupposes an objective
concept of truth.


Nietzsche's watchword was ``will to power''. It is this will that determines what is to be truth. Art stands above
science, but life stands above art; Apollo stands above Socrates, but Dionysus stands above Apollo! The wonderful and
overflowing fullness of life can be contained neither in the forms that thought has put in order nor in the forms
that art can create. In particular Nietzsche strove for a liberation from ``the theoretical man''. What mattered to
him was ``die grösste Erhöhung des Kraftbewusstseins''. To his closest friend (Overbeck) he once said: ``Die reine
Aufklärung genügt mir nicht; es muss doch etwas Anderes geben''. This makes his attack on Strauss understandable; his
position here recalls Rousseau's relation to Voltaire and the Encyclopedists.

The two men have, each in his own way, known the constant tension between thought and life. Life, both the inner and
the outer life, brings us such a fullness of content that thought cannot at once master it. Therefore life
constantly has a head start, and only gradually and in part can thought produce the rational forms in which the
content of life can be ordered. This relation is best expressed in a journal entry of Kierkegaard's: ``Life must be
understood backwards; but we live forwards; life in temporality therefore never becomes properly understandable,
% --- p. 192 ---
\setcounter{footnote}{0}%
\opage{192}because at no moment can I get complete quiet to take up the position backwards.'' To this correspond
Nietzsche's words: ``Erkenntnis und Werden schliesst sich aus''. ---

Another analogy between the two thinkers I find in the standard they both apply in order to determine the height
of a personality's life and development. Their standard testifies to a deep understanding of the nature of
personality.

Kierkegaard has presented, partly in poetic, partly in philosophical form, a series of types (or, as he calls them,
``stages'') of different views of life and ways of taking life, --- types whose difference from one another rests on
the degree and range of the oppositions that life comprises, and thereby of the resistance that must be overcome.
Personal life stands the higher, the greater and sharper the oppositions are that can persist without the unity and
concentration of the personality being annulled. Lowest stands he who is absorbed in the enjoyment of the
\textit{moment}, in the arbitrary playing with the circumstances of life that characterizes the ``aesthetic'' view of
life. After this come, by degrees, the ironist, the ethicist, the humorist and the religious man. Increased content of
life and increased energy, but also increased pain, is the law of life on the higher stages. The pain is caused by
the conflict of the oppositions, finally --- for the religious man --- between time and eternity. By abating the
oppositions one lowers the level of life. The highest is the great courage that is required when spirits are tried
in strife.

This courage was for Nietzsche too the highest virtue. Throughout his whole activity as an author he fought to strike
down all pessimism. All dwelling on the dark sides of life, on the weak and the pitiable, all pity, all mistrust of
happiness and of the happy were to be rooted out. But his passionate optimism did not make him blind to suffering.
He saw, like Kierkegaard, that when great demands are made on life, great possibilities of suffering arise at the
same time. Suffering he knew from his own experience, ill as he was long before he succumbed to insanity; but it
could not bring him to bear false witness about life. He was a brave man. Longest he had to struggle
% --- p. 193 ---
\setcounter{footnote}{0}%
\opage{193}with his contempt for ``the great mob'', a feeling that underlay the opposition between master morality
and slave morality which was perhaps what made the most impression on his contemporaries. It is clear, of course,
that both the contempt and the social dualism between masters and slaves mark a great inconsistency in one who wants
to strike down pessimism. One can indeed perceive, especially through fragments and drafts, an unceasing striving to
remove this inconsistency. It was Nietzsche's tragedy that he did not come to complete the development of his
thoughts on this subject.

Kierkegaard saw the danger especially in aesthetic play with life; Nietzsche saw it in weakness and complaint; but
both maintained that life was to be conceived as a struggle to preserve both the changing richness of life and its
unity and coherence.

A danger, finally, of an entirely opposite kind is one to which neither of them is blind either, --- the one that
repetition brings.

Kierkegaard carries out the finest analyses here, and the problem of repetition appears at every point in his
reflection on life. The main point is the dulling power of repetition. It robs the life of feeling of its freshness
and strength. One can, with the ``aesthetic'' type, avoid repetition by keeping oneself hovering, so that one
withdraws from every lasting relation in life. But one can also struggle with repetition, by directing the will
toward preserving the originality and ideality of life amid the daily work and within the given social conditions.
Here too Kierkegaard found a standard for the degree of development of a personality: ``He who wills repetition, he
is a man''. --- Repetition plays a special role in his religious polemic. It was his last great accusation against
the church and Christendom that they had succumbed to the dulling power of repetition, which began to work from the
moment when Christianity became a member of existing human institutions and forgot its divine origin and its first
enthusiasm.


Nietzsche came upon the problem of repetition more from outside. His Greek studies led him to the thought that
everything must
% --- p. 194 ---
\setcounter{footnote}{0}%
\opage{194}repeat itself. The Greeks did not have the idea of a constantly progressing development, but in connection
with the idea of the limitation of the universe in space there awoke among them (especially among the Stoics) the
thought of a rhythmic recurrence of the same situations. And Nietzsche believed he found this thought confirmed by his
studies of modern natural science. For since one --- so he believed --- must assume that the number of atoms was
limited, the number of their combinations must also be limited. --- This thought (which has wrongly been regarded as
a symptom of insanity) filled Nietzsche with the greatest horror. For a long time he went about carrying it as a great
secret. It must, it seemed to him, lead back to the pessimism he had set himself the task of rooting out of the
world. The hourglass of existence is turned ever anew! Time is a ring, and the wheel of the world rolls eternally
without moving from the spot! --- Like Kierkegaard he saw the only way out in a mighty, creative will that constantly
brings forth new values. In the planned conclusion of the Zarathustra book the proclamation of such an enthusiastic
and energetic will to life in defiance of the fact of repetition was to be a main feature\footnote[1]{On the
influence of repetition on the life of feeling see my Psychology VI E, 4---5. --- While the law of repetition filled
Nietzsche with the greatest horror and sets Kierkegaard one of the most decisive problems in his doctrine of the
``stages'', it had for Marcus Aurelius the great significance of a calming thought, which contained the assurance
that he would meet nothing on his way that his will should not be prepared to receive and take up. Only on a rare
occasion does the thought of the dulling working of repetition come forward. (See the essay on Marcus Aurelius
above). --- This opposite way of taking a stand toward repetition is characteristic of the difference between the
ancient and the modern outlook on life. (See on this also the essay on Plato's Symposium above).}. ---

Kierkegaard and Nietzsche do not stand alone in the weight they lay on the relation of tension between life and
thought. In the course of the last century we see a whole series of thinkers who, like them, are symptoms of
something that is stirring in the age, and that will surely prove to be stronger both than
% --- p. 195 ---
\setcounter{footnote}{0}%
\opage{195}the old dogmas and than the new doubts. Carlyle, Guyau and William James bear witness in the same
direction, though not with the passion that drove Kierkegaard and Nietzsche out into their paradoxes.

Reflection --- what a poet calls ``the deep calm of the still, the faithful reflection'' --- is the philosopher's
spiritual home. But reflection presupposes experience. Thought gets its content and its tasks from life, and without
faith in the life in us and around us we cannot believe in thought. But the time of thought comes ever anew; it need
not hamper or subdue life for fear of not being able to master it. Nor need it take refuge in theological dogmas or
metaphysical speculations. Life itself will best be able to show in what forms of thought it can be grasped. As long
as life moves forward, its whole content cannot be wholly and completely put in order by thought; but it will always
be a matter of keeping the mind open to the new that life and experience can bring.

% --- p. 196 ---
\setcounter{footnote}{0}%
\opage{196}\essay[At the Krøyer Commemoration (1909)]{AT THE COMMEMORATION OF PETER SEVERIN \\ KRØYER}
\undertitel{4 DECEMBER 1909.}
\streg

O\textsc{ne} beautiful summer afternoon about 27 years ago a band of younger men walked across the Eremitage Plain.
They were members of a society that had been founded a couple of years before, and that consisted for the most part
of painters and writers. Krøyer, who had recently come home from a long stay abroad, walked at the head, and I see
him clearly before me, gazing out over the glorious countryside and pointing --- as he used to --- with his whole
hand (as in Viggo Johansen's evening picture) out toward the horizon. He had won his first victories and now moved
in our midst, bright and clear, fresh and warm-hearted, and all were glad to have him back.

Together with many others in the little band of painters he marked a new age in Danish painting. There had been a
decline and a slackening when the masters who in the '40s and '50s had produced such beautiful works had for the most
part passed away, and the judgment that had been passed abroad on our art (at the exhibitions in Vienna in 1873 and
in Paris in 1878) had been severe. Now there was a new generation with fresh courage and a new vision, and in it
Krøyer stood in the front rank. There was a joy throughout the whole band over every work of value that appeared. It
was art that mattered, and the comrades rejoiced just as much as the single fortunate master. There was in this
respect, as I remember it, a greater fellow-feeling among the painters than among the writers; and our society was
in fact dissolved because of frictions among
% --- p. 197 ---
\setcounter{footnote}{0}%
\opage{197}the writers. I believe that this fine relation among the artists has held. Oppositions may have come
forward between different schools and points of view. But I am still of the opinion that the visual artists, in
respect of fellow-feeling on behalf of art, stand above the writers. It was a great joy to move among them.

There was a question that was much discussed at that time in this circle too. It was the question of the national
element in art. Our two great art historians of the last century, Høyen and Julius Lange, viewed this question
somewhat differently. \textsc{Høyen} laid great weight on visual art's taking its subjects from the people's own
nature and history. Art had come to us from outside; but it was to become the people's property by taking its
subjects from its nature and its history or mythology. \textsc{Julius Lange} maintained, on the contrary, that what
mattered first and foremost was that the work of art should spring from the artist's own innermost personality, no
matter what subjects he then chose. Art is life, and all life must be lived at first hand. It is the nation that is
to come to see the demands of art, not art that is to accommodate itself to the nation's. But from this point of view
he came to a result similar to Høyen's; for he believed that native subjects would most easily be able to set the
artist's own personality in motion. In foreign subjects he will not so easily be able to become at home. And to
travel abroad to learn technique Lange finds unreasonable, since technique has no worth except as the form of a
personal work --- and every such work brings its own technique with it.

Yet Lange too surely viewed this matter too anxiously. If there is a real and original gift, it can, without harm to
its distinctiveness, learn a great many things from others, even from complete strangers. True geniuses are
distinguished from ordinary people not least by being able to learn more than others, because they have a more open
sense and a deeper urge. And as for the national element, it will --- without any stress being laid on it ---
involuntarily and unconsciously assert itself when the artist, as Lange indeed demands,
% --- p. 198 ---
\setcounter{footnote}{0}%
\opage{198}works with his whole innermost personality. When all the strings of the soul sound, the national
distinctiveness will assert itself; for it is not something external that is added, but lies in the way in which the
individual involuntarily does his work. Therefore it will also be quite able to bear the learning from abroad, where
this really has value.

Krøyer is a great and fine example of this. His personality and his art were developed by what he learned in
Bonnat's atelier and by the study of the older Spanish painters. The foreign became his property, and his finest
works have in rare degree become the property of the Danish people. In rare degree our people has found in them what
can gladden its eye and its mind. Through his development and his work a proof has been given of the independence of
Danish intellectual life. Even in the small country the human mind, the human eye and the human hand can work in a
new and distinctive way, so that its art even gains significance abroad and becomes more than an offshoot of the art
of other countries. --- Krøyer is far from the only artist of whom this holds. And it holds besides also for our
literature, our science and for the development in the practical spheres of life. But he is one of the best and
greatest examples.

What, then, is the distinctiveness of Krøyer's art? --- He is in the first place the painter of light and of open
nature, shows us the glory of life in sunshine and clarity. It has been objected from several sides against his
pictures of nature that the human beings vanish before the great effects of light and air. But however absorbed he
is in the glory of light, man is never vanishing or mere staffage for him. In the large picture of a summer night at
Skagen the fishermen are rather the main element; their powerful figures recall the hard work; they wait only for
the breeze to begin it; the idyll of the summer night thereby gets a serious background. In the glorious evening
picture where the two young women walk along the beach, and in the coastal picture with the bathing boys, it is
precisely something essential that the human beings are there; the festiveness and clarity of nature pass over into
their minds and stamp themselves on their movements. --- It has
% --- p. 199 ---
\setcounter{footnote}{0}%
\opage{199}also been held against him that only once, in the picture of the Italian hatmakers, has he followed the
light into the dark, wretched rooms where so many human lives are spent. When the picture first appeared, the
opposite complaint was raised against it: it was marked by a piercing pessimism and was meant to serve socialist
agitation! Both objections sin against the artist's individuality. He simply wanted to follow his favorite subject,
light, on its way, and thereby for once came in with it into the wretched hut. When otherwise he seeks out its
working in free, open nature, he follows his urge, seizes what lies near to his mind, his eye and his hand. He has
seized aspects of our country's nature that we constantly need to have brought out; gray weather and sleet we truly
have enough of. His works will continue to be a source of joy in life, both in themselves and by teaching us how we
should look at our country and its nature. And besides --- the glory that is represented is the glory of free nature
itself, of light and of air, which can become accessible to all. The boys who leap about in the water in the clear
summer sun are boys of the common people, and the picture shows us that there are goods and joys from which no one
\textit{needs} to be excluded. Quite as much as a realistic depiction of the dark dwellings of the unfortunately
placed, a picture such as this will testify to human sympathy. There lies a wish behind it --- a wish that no child
in our country might be excluded from the nature and the joy of life it depicts! --- There is here an opposition that
asserts itself in analogous fashion in all spheres of life: is the law or is the gospel to be put in the front rank?
is it Noureddin or Aladdin who is to provide the type of life?

A single artist cannot exhaust everything. There is much that Krøyer did not take in, because it was not his
concern. Henrik Ibsen once in a conversation compared the Danish landscape to a soulful face that constantly changed
its expression with the changing moods. This side of Danish nature --- the role of the clouds and the wind --- it did
not lie in Krøyer to depict. But what he has seen has an imperishable value.

% --- p. 200 ---
\setcounter{footnote}{0}%
\opage{200}There is another remembered image of Krøyer that has often come forward for me in these days. About twelve
years ago he regularly attended the meetings of the Academy of Sciences, in order to make preliminary studies for his
well-known painting. Sitting behind the blackboard he made sketches for the arrangement of the picture, or for single
figures. Attention often left the scientific lecture at a glance at the clear, sharply observing eyes that from their
hiding place were directed at the assembly or at a single member of it. In the picture that thus came into being, his
great gift for grouping as well as for the characterization of the single individuals proved itself. He took joy in
the diversity of the types and sought to bring it out by a sharp stress on one or another particular peculiarity in
the individuals. Through this bringing out there came a greater fullness of content into the picture. It has often
been thought that he came close to caricature through this strong emphasis on single traits. Far rather, however,
this sprang from an interest in the special peculiarities and from an eagerness to avoid repetition and uniformity.
He even said to me in jest that, as a contrast to the many attentive and listening faces at the lecture, he would
have liked to have one who sat asleep; but of course no model could be found for that. Sometimes he could be
mistaken in his inference from physiognomy to character. There was one member whom he, in a conversation I had with
him, described as a Viking nature, which that decent man was very far from being, as Krøyer too discovered on closer
acquaintance.

Behind his striving to bring out the characteristic peculiarities lay a great human sympathy, a joy in the fullness
and multiplicity of life within humanity too. And when the conversation turned to common friends, an amiable warmth
and cordiality came out in him. Even one who has met him only a few times on his way will keep the memory of his
open, cheerful, art-loving and friendly mind.

% --- p. 201 ---
\setcounter{footnote}{0}%
\opage{201}I cannot here go through Krøyer's whole world of pictures, one of the most precious treasures our
generation leaves behind. And such a survey would also only have significance if it were undertaken by an art
theorist or art historian. I have understood my task here tonight to be that I should speak on behalf of the many who
have rejoiced in what he has given his people, without their presuming to have the ability and the authority to
determine his place within the world of art and in relation to his fellows, who have often developed along quite
other paths than he. --- I merely recall that he has given, in groups and in single pictures, striking and
characteristic pictures of our fishermen, our artists, our businessmen, besides our scholars. He leaves to the
following generations a glorious record of how our nature and how the people of our time looked when seen with a
great artist's eye, an eye that had a predilection for light in all its nuances, and for human characters in all
their multiplicity.

And the master's own bright image, which he himself has indeed repeatedly given us, will accompany his works and live
in memory, the image of one of the most glorious artistic personalities our people has fostered.

Like every artist, like every intellectual worker, he has had his struggles to endure: the struggle with the dark
times when nothing would succeed and when even the inner possession seemed to fail, --- the struggle with refractory
subjects that would not take shape, --- the struggle with the prejudices and habits of the beholders. And now in
later years the struggle with sorrow and illness. Despite all, he stands as a figure of light and will continue to
stand so in his people's memory and in its gratitude.

He was a son of light. Light and honor he has brought our people.

Honor be to his memory!

% --- p. 202 ---
\setcounter{footnote}{0}%
\opage{202}\essay[Darwin and Philosophy (1909)]{CHARLES DARWIN AND PHILOSOPHY\footnote[1]{Printed in English in the Festschrift \textit{Darwin and Modern Science} (1909), published in Cambridge.}.}
\undertitel{A consideration in the centenary of Darwin's birth and the fiftieth year\\ since the publication of the ``Origin of Species''.}
\streg

I\textsc{t} is not so easy as many believe to draw a sharp line between philosophy and natural
science. The naturalist who introduces a new principle, or demonstrates
a fact which throws a new light on existence, not only renders an
important service to philosophy but is himself a philosopher in the
broader sense of the word. The aim of philosophy in the stricter sense
is to attain points of view from which the fundamental phenomena and
the principles of the special sciences can be seen in their relative
importance and connection. But philosophy in this stricter sense
has always been influenced by philosophy in the broader sense. Greek
philosophy came especially under the influence of logic and mathematics, modern
philosophy under the influence of natural science. And here the name of Charles
Darwin stands with those of Copernicus, Galileo, Newton, and Robert Mayer.
These names denote new problems and great alterations in our conception of the
universe.

First of all we must in this connection lay stress on Darwin's own personality. His deep
love of truth, his indefatigable inquiry, his wide horizon, and his
steady self-criticism make him a scientific model, even if his results
and theories should eventually come to possess only an historical
interest.
% --- p. 203 ---
\setcounter{footnote}{0}%
\opage{203}In all domains of thought the primary object is to reach
high summits from which wide surveys are possible, to reach them toiling
honestly upwards by way of experience, and then not to turn dizzy when
a summit is gained. Darwinians have sometimes turned dizzy, but Darwin
never. He saw from the first the great importance of his hypothesis---not
only because of its solution of the old problem as to the value of the
concept of species, not only because of the grand picture of successive
evolution which it unrolls before thought and imagination, but also, and not least, because of the life and inspiration
it would impart as a principle of method, by calling forth investigations in the study of comparative anatomy, of instinct
and of heredity, and finally because of the influence it would exert on
the whole conception of existence. He wrote as early as the
year 1837, shortly after his return from the great voyage which furnished him with the most important material for his hypothesis: ``My theory would give zest to recent and fossil comparative
anatomy; it would lead to the study of instinct, heredity, and
mind-heredity, whole (of) metaphysics.''

We can distinguish four main points in which Darwin's investigations
possess philosophical importance.

The evolution hypothesis is older than Darwin; it is, indeed, one
of the oldest guessings of human thought. In the eighteenth century it
was put forward in an interesting manner by Diderot and Lamettrie. Kant embraced it (as is seen from his essay
``Muthmaaslicher Anfang des Menschengeschlechts'' 1786), though in the ``Kritik der Urteilskraft'' (1790) he makes critical reservations.
As we shall see later, it was held also by several philosophers in the
first half of the nineteenth century. In his preface to ``The Origin of
Species'', Darwin himself mentions the naturalists who were his forerunners.
--- But he has set forth the old hypothesis in so energetic
and thorough a manner that it perforce attracts the attention of
all thoughtful men in a much higher degree than before.

And this leads us to the second point. He,
much more than his predecessors, even than Lamarck, sought a foundation
for his hypothesis in definite facts.
% --- p. 204 ---
\setcounter{footnote}{0}%
\opage{204}Modern science began with the demand for what
Kepler and Newton called vera causa, i.e. with the principle that the cause to which a phenomenon is referred must be demonstrable in experience, just as much as the phenomenon itself which is to be explained. This great principle Darwin set himself to satisfy by means of the wealth of
material which he collected by his observations, his experiments and his reading.
He not only revived an old hypothesis, but he saw the necessity of
verifying it by facts. Whether the special cause on which he founded the
explanation of the origin of species---``Natural Selection''---is the true cause,
is now a subject of doubt and investigation. He himself had some doubt in regard
to this question, and the criticisms which are directed against his
hypothesis hit ``Darwinism'' rather than Darwin. In his indefatigable search
for empirical evidence he has been a model even for his antagonists.
He has compelled them to approach the problems of life along other lines and by other methods than
those which were formerly followed.

Whether the special cause to which Darwin appealed is sufficient or not,
at least to it is due the greater part of the influence which
he has exerted on the general trend of thought. ``Struggle for existence''
or ``natural selection'' are principles which have been applied in every department of thought. Recent research, it is true, has
discovered greater leaps (``mutations''), more discontinuity in evolution as it appears to observation and experiment, than
Darwin supposed. He believed in the importance of small variations slowly
accumulated. It has also been shown---by the experimental method, which
in recent biological work has succeeded Darwin's more historical method---that types once constituted possess great permanence, the changes (``fluctuations'')
being restricted within clearly defined boundaries. The problem has
become more precise, both as to variation and as to heredity. The inner
conditions of life have in both respects shown a greater independence
than Darwin had supposed in his theory,---though he always admitted that
the cause of variation was
% --- p. 205 ---
\setcounter{footnote}{0}%
\opage{205}to him a great enigma, ``a most perplexing
problem,'' and that the struggle for life could only occur where
variation already existed. But, at any rate, it was of the greatest importance
that Darwin gave his age so living an impression of the struggle for life which is
everywhere going on, and to which even the highest forms and ideals
must be amenable. The philosophical importance of this idea does not
stand or fall with the answer to the question, whether it can furnish a sufficient explanation of the origin of species. It has
an independent, positive value for everyone who will observe life and
reality with an unbiased mind.

Finally, it is of interest to see that in accentuating the struggle for life Darwin stands as a
characteristically English thinker. He continues a train of ideas which
Hobbes and Malthus had already begun. Moreover in his critical views as
to the conception of species he had English forerunners---in the middle
ages Occam and Duns Scotus, in the eighteenth century Berkeley and Hume.
In his moral philosophy he is an adherent of the
school which is represented by Hutcheson, Hume and Adam Smith.
--- Just because
he is no philosopher in the stricter sense of the term, it is of great
interest to see that his attitude of mind is that of the great thinkers
of his nation. ---

In considering Darwin's influence on philosophy more specially I will begin with an
examination of the attitude of philosophy to the conception of evolution
at the time when ``The Origin of Species'' appeared. I will then examine
the effects which the theory of evolution, and especially the idea
of the struggle for life, has had, and naturally must have, on the
discussion of philosophical problems.

\begin{center}II.\end{center}


When ``The Origin of Species'' appeared fifty years ago Romantic
speculation, Schelling's and Hegel's philosophy, still reigned on the
continent, while in England Positivism, the philosophy of Comte and
Stuart Mill, represented the most important trend of thought.

a. German speculative philosophy had much to say on evolution,
% --- p. 206 ---
\setcounter{footnote}{0}%
\opage{206}and pretended to be a ``philosophy of evolution''. But then the word evolution was to be taken
in an ideal, not in a real, sense. To speculative thought the forms and
types of nature formed a system of ideas, within which any form could
lead us by continuous transitions to any other. It was a classificatory
system which was regarded as a divine world of thought or images, within
which metamorphoses could go on---a condition comparable with that in
the mind of the poet when one image follows another with imperceptible
changes. \textsc{Goethe's} ideas of evolution, as expressed in his ``Metamorphosen
der Pflanzen und der Thiere'', belong to this category; it is, therefore,
incorrect to call him a forerunner of Darwin. \textsc{Schelling} and \textsc{Hegel}
held the same idea; Hegel expressly rejected the conception of a real
evolution in time as coarse and materialistic. ``Nature,'' he says, ``is
to be considered as a system of stages, the one necessarily arising from
the other, and being the nearest truth of that from which it proceeds;
but not in such a way that the one is naturally generated by the other;
on the contrary (their connection lies) in the inner idea which is the
ground of nature. The metamorphosis can be ascribed only to the notion
as such, because it alone is evolution ... It has been a clumsy idea in
the older as well as in the newer philosophy of nature, to regard the
transformation and the transition from one natural form and sphere to
another as an outward and actual production''. (Encyclopädie der philosophischen
Wissenschaften § 249).

The only one of the philosophers of Romanticism who believed in a
real, historical evolution was \textsc{Oken} (Lehrbuch der Naturphilosophie.
1809). In Denmark---in express opposition to German speculation---the historical reality of evolution was maintained by \textsc{Treschow} (Elementer til Historiens Filosofi 1812)
and by \textsc{Sibbern} (Spekulativ Kosmologi 1846), as also, from a more scientific starting-point, by \textsc{Olufsen} (Om Menneskets
Rolle i den fysiske Verden. 1822). ---

\textsc{Schopenhauer's} philosophy has a more realistic character than
that of Schelling's and Hegel's, his diametrical opposites, though he
also belongs to the romantic school of thought. His physiological
% --- p. 207 ---
\setcounter{footnote}{0}%
\opage{207}and psychological views were greatly influenced by French naturalists and
philosophers, especially by Cabanis and Lamarck. He praises the ``ever
memorable Lamarck,'' because he laid so much stress on the ``will to
live.'' But he repudiates as a ``wonderful error'' the idea that the
organs of animals should have reached their present perfection through
a development in time, during the course of innumerable generations. It
was, he said, a consequence of the low standard of contemporary French
philosophy, that Lamarck came to the idea of the construction of the species of living
beings in time through succession (Über den Willen in der Natur. 1836). ---

b. The positivistic stream of thought was in its way not more in favor of a real
evolution than was the Romantic school. Its aim was to adhere to
positive facts or groups of facts: it looked with suspicion on far-reaching hypotheses.
\textsc{Comte} laid great stress on the difference, the discontinuity found between the different
domains of experience, as well as within each single domain. As he
regarded as unscientific every attempt to reduce the number of physical
forces, so he rejected entirely the hypothesis of Lamarck concerning the
evolution of species; the idea of species would in his eyes absolutely
lose its importance if a transition from species to species under the
influence of conditions of life were admitted. His disciples (\textsc{Littré},
\textsc{Robin}) continued to direct against Darwin the polemics which their
master had employed against Lamarck.

\textsc{Stuart Mill}, who, in the theory
of knowledge, represented the empirical or positivistic movement in
philosophy---like his English forerunners from Locke to Hume---founded
his theory of knowledge and morals on the experience of the single
individual. He sympathized with the theory of the original likeness of
all individuals and derived their differences (on which he practically
laid much stress) from the influence both of
experience and education, and, generally, of physical and social causes.
He admitted an individual evolution, and, in the human species, an
evolution based on social progress; but no biological-physiological
% --- p. 208 ---
\setcounter{footnote}{0}%
\opage{208}evolution of
species. He was afraid that the hypothesis of heredity would carry us
back to the old theory of innate ideas. ---

Darwin was more empirical than Comte and Mill; experience disclosed to
him a deeper continuity than they could find. Closer than before the
nature and fate of the single individual were shown to be interwoven in
the great web binding the life of the species with nature as a whole.
And the continuity which so many idealistic philosophers could find
only in the world of thought, he showed to be present in the world of
reality.

\begin{center}III.\end{center}


Darwin's energetic renewal of the old idea of evolution had its chief
importance in strengthening the conviction of a real continuity in
the world, of continuity in the series of forms and events. It was a
great support for all those who were prepared to base their conception
of life on scientific grounds. Together with the recently discovered law
of the conservation of energy, and in general with the new era which began for natural science about 1840, it helped to produce the great realistic
movement which characterized the last decades of the nineteenth century.
After the decline of the Romantic movement people wished to have firmer
ground under their feet and reality now asserted itself in a more
emphatic manner than in the period of Romanticism. It was easy for Hegel
to proclaim that ``the real'' was ``the rational,'' and that ``the
rational'' was ``the real'': reality itself existed for him only in the
interpretation of ideal reason, and if there was anything which could
not be merged in the higher unity of thought, then it was only an
example of the ``impotence of nature to hold to the idea.'' But now
concepts are to be founded on nature and not on any system of categories
too confidently deduced a priori. The new devotion to nature had its
recompense in itself, because the new points of view made us see
% --- p. 209 ---
\setcounter{footnote}{0}%
\opage{209}that nature could indeed ``hold to ideas,'' though not to those which
we had cogitated beforehand. ---

A most important question for philosophers to answer was whether the
new views were compatible with an idealistic conception of life and
existence. Some proclaimed that we have now no need of any philosophy
beyond the principles of the conservation of matter and energy and of the
natural evolution of organic life. Existence should be completely explained by the laws of material nature. But abler
thinkers saw that the thing was not so simple. They saw it as their task to
give the new views their just place in the whole conception of life and of the world and to examine what alterations the
old views must undergo in order to be brought into harmony with the new
data.

The realistic character of Darwin's theory was shown not only in the
idea of natural continuity, but also, and not least, in the idea of the
cause whereby organic life was supposed to advance step by step. The idea of
the struggle for life implied that nothing could persist, if it had no
power to maintain itself under the given conditions. Inner value
alone does not decide. Idealism was here put to its hardest trial. In
continuous evolution it could perhaps still find an analogy to the inner
evolution of ideas in the mind (and it was precisely on such an analogy that the philosophy of Romanticism was founded, without being conscious of it); but in the demand for power in order to
struggle with outward conditions Realism seemed to announce itself in
its most brutal form. Every form of Idealism had then to ask itself
how it was going to struggle for life with the new doctrine.

I will give a short account of the position which leading thinkers
in the principal countries of Europe took up in regard to this question.

a. \textsc{Herbert Spencer} was the philosopher whose mind was best prepared by
his own previous thinking to admit the theory of Darwin to a place in
his conception of the world. His criticism of the arguments which had
been
% --- p. 210 ---
\setcounter{footnote}{0}%
\opage{210}put forward against the hypothesis of Lamarck, showed that Spencer,
as a young man, was an adherent to the evolution idea. In his
``Social Statics'' (1850) he applied this idea to human life and moral
civilization. In 1852 he wrote an essay on the Development Hypothesis,
in which he definitely stated his belief that differences of
species, like individual differences, were the
result of development. In the first edition of his Psychology (1855)
he took a step which put him in opposition to the older English school
(from Locke to Mill), in that he acknowledged an organic basis for the mental development of the individual, a basis which had come into being in the course of
generations through the inheritance of acquired qualities, so that the nature and functions of the individual are only
to be understood through its connection with the life of the whole species.
In 1857, in his essay on ``Progress'', he propounded the law of
differentiation as a general law of evolution, verified by examples from
all regions of experience, the evolution of species being only one of
these examples. On the effect which the appearance of Darwin's ``Origin of
Species'' had on his mind he writes in his Autobiography: ``Up to
that time I held that the sole cause of organic evolution is the
inheritance of functionally-produced modifications. The `Origin of
Species' made it clear to me that I was wrong, and that the larger part
of the facts cannot be due to any such cause. To have the theory of
organic evolution justified was of course to get further support for
that theory of evolution at large with which all my conceptions were
bound up.'' --- Instead of the metaphorical expression ``natural selection,''
Spencer introduced the term ``survival of the fittest,'' which found favor with
Darwin as well as with Wallace.

In working out his ideas of evolution, Spencer found that
differentiation was not the only form of evolution. In its simplest form
evolution is mainly a concentration, previously scattered elements being
integrated and losing independent movement. Differentiation is only
forthcoming when minor wholes arise within a greater whole. And the
highest form of evolution is reached when there is a
% --- p. 211 ---
\setcounter{footnote}{0}%
\opage{211}harmony between
concentration and differentiation, a harmony which Spencer calls
determination and which he defines as a moving equilibrium. At the same
time this definition enables him to illustrate the expression ``survival
of the fittest.'' ``Every living organism,'' he says, ``exhibits such a moving
equilibrium---a balanced set of functions constituting its life; and the
overthrow of this balanced set of functions or moving equilibrium is
what we call death. Some individuals in a species are so constituted
that their moving equilibria are less easily overthrown than those of
other individuals; and these are the fittest which survive, or, in Mr
Darwin's language, they are the select which nature preserves.''
--- Not only in the domain of organic life, but in all domains,
the summit of evolution is, according to Spencer, characterized by such
a harmony---by a moving equilibrium.

Spencer's analysis of the concept of evolution, based on a great variety
of examples, has made this concept clearer and more definite than
before. It contains the three elements; integration, differentiation and
harmony. It is true that a concept which is to be valid for all
domains of experience must have an abstract character, and between its forms in the
several domains there is, strictly speaking, only a relation of analogy, e.g.
between psychical and physical evolution, and between organic and inorganic evolution. But
this is no serious objection, because general concepts do not express
more than analogies between the phenomena which they represent. In defining evolution
(in the simplest form) as integration of matter and dissipation of
movement, Spencer takes his leading terms from the material world. But as he---not always quite consistently\footnote[1]{Cf. my letter to Spencer, 1876, now printed in \textit{Duncan}, Life and Letters of Herbert Spencer, p. 178.}---assumed a correspondence of mind and
matter, he could very well give these terms an indirect importance for
psychical evolution. Spencer has always, in my opinion with full right,
repudiated the ascription of materialism. He is no more a
% --- p. 212 ---
\setcounter{footnote}{0}%
\opage{212}materialist than Spinoza. In his ``Principles of Psychology'' (§ 63) he
expressed himself most clearly: ``Though it seems easier to translate
so-called matter into so-called spirit, than to translate so-called
spirit into so-called matter---which latter is indeed wholly
impossible---yet no translation can carry us beyond our symbols.'' --- These
words lead us naturally to a group of thinkers whose starting-point
was psychical evolution. But we have still one aspect of Spencer's
philosophy to mention.

Spencer founded his laws of evolution on comparative experience, but
he was convinced that they could be deduced from the law of the
conservation of energy. Such a deduction is, perhaps, possible for the
more elementary forms of evolution, integration and differentiation; but
it is not possible for the highest form, which is a
harmony of these two forms. Spencer can no more deduce
the necessity of harmonious totalities, ``moving equilibria,'' than Hegel could guarantee the ``higher unities''
in which all contradictions should be reconciled. In Spencer's hands the
theory of evolution acquired a more decidedly optimistic character
than in Darwin's. --- I shall deal later with the relation of Darwin
to the opposition of optimism and pessimism.

b. While the starting-point of Spencer was taken in outer experience, was biological or cosmological,
psychical evolution being conceived as in analogy with physical, a group
of thinkers---in Germany \textsc{Wundt}, in France \textsc{Fouillée}, in Italy
\textsc{Ardigò}---took, each in his own manner, their starting-point in psychical
evolution as an original fact and as a type of all evolution; the
hypothesis of Darwin comes in here too as a corroboration and as a special
example. These thinkers find the continuity of evolution
most prominent in psychical evolution, and this is for them a
motive to demand a corresponding continuity in the material, especially
in the organic domain.

To Wundt and Fouillée the concept of will is prominent. They see the
type of all evolution in the transformation of the life of will from
blind impulse to conscious choice; the theories of Lamarck and Darwin
% --- p. 213 ---
\setcounter{footnote}{0}%
\opage{213}are used to support the view that there is in nature a tendency to
evolution in steady reciprocity with external conditions. The struggle
for life is here only a secondary fact. Its apparent prominence is
explained by the circumstance that the influence of external conditions
is easily made out, while inner conditions can be inferred only
from their effects and do not immediately catch the eye. --- For Ardigò the evolution of thought was the
starting-point and the type. The history of every scientific hypothesis
shows a progress from the indefinite (indistinto) to the definite
(distinto), and such a progress is characteristic of all evolution, as Ardigò
seeks to show in a series of works. The opposition between indistinto and
distinto corresponds to Spencer's opposition between concentration and
differentiation. The hypothesis of the origin of differences of species
from more simple forms is then a special example of the general law of
evolution.

In the views of Wundt and Fouillée we find the fundamental idea of
idealism: psychical phenomena as expressions of the innermost nature of
existence. They differ from the older (speculative) Idealism in the great stress which
they lay on evolution as a real, historical process which is going on
through steady conflict with external conditions. The Romantic dread
of reality is broken. It is beyond doubt that Darwin's emphasis on the
struggle for life as a necessary condition of evolution has been a very
important factor in carrying philosophy back to reality from the heaven
of pure ideas.
--- The philosophy of Ardigò, on the other side, appears
as a continuation and deepening of positivism, though the Italian
thinker arrived at his point of view independently of French-English
positivism. The idea of continuous evolution is here maintained in
opposition to Comte's and Mill's philosophy of discontinuity. From Wundt
and Fouillée Ardigò differs in conceiving psychical evolution not as an
immediate revelation of the innermost nature of existence, but only as a
single, though the most accessible example, of evolution.

c. The French philosophers \textsc{Boutroux} and \textsc{Bergson} lay
greater stress on the difference between psychical and material
% --- p. 214 ---
\setcounter{footnote}{0}%
\opage{214}evolution than the thinkers hitherto mentioned do. What to immediate
intuition, purified of all reflection, stands as a continuous evolution proceeding through qualitative changes,
becomes for the outer apprehension, bound to differences of time and space, on which natural science builds, always broken up
into fragments, which are brought into connection with one another in a mechanical way. To Bergson there lies at the basis
of all evolution a spontaneous striving forward (élan de vie), which to our ever fragmentary
observation and analytic reflection appears as an outward manifold of
elements and processes. The very concept of matter in its scientific form
is the result of the breaking asunder of the immediately given which all scientific
reflection must perform. --- In this conception, which dominates Bergson's philosophy, the strongest opposition between
inner and outer conditions of evolution is expressed: in the domain of
internal conditions immediate (spontaneous) development of qualitative forms,---in the
domain of external conditions discontinuity and mechanical combination.

We see, then, that the theory of evolution has influenced philosophy
in a variety of forms. It has made idealistic thinkers revise their
relation to the real world, and it has led positivistic thinkers to find a
closer connection between the facts on which they base their views. And it
has made us all open our eyes for new possibilities---through
the prima facie inexplicable variations which are the
condition of all evolution. This last point is one of peculiar interest.
Deeper than speculative philosophy and mechanical science saw in the
days of their triumph, we catch sight of new streams, whose sources and
laws we have still to discover. Most sharply does this appear in
\textsc{Hugo de Vries}'s theory of mutation, which is only a stronger accentuation of a main
point in Darwinism, the inexplicability of variations. It is interesting to see that an analogous
problem comes into the foreground in physics through the discovery
of radioactive phenomena, and in psychology through the assumption
of psychical new formations (as held by Boutroux, William James and
Bergson). From this point of view, Darwin's ideas, as well as the
% --- p. 215 ---
\setcounter{footnote}{0}%
\opage{215}analogous
ideas in the physical and psychological domains, incite us to renewed examination of the first
principles of our knowledge, their foundation and their value. On the other hand,
his theory of the struggle for existence challenges us to examine the
conditions and discuss the outlook as to the persistence of the human forms of life
and of human values. It is not enough to
hope (or fear) the rising of new forms; we have also to investigate the
possibility of upholding the forms and ideals which have hitherto been
the bases of human life. Darwin has here given his age the most earnest
and most impressive lesson. This side of Darwin's theory is of peculiar
interest to some special philosophical problems to which I now pass.

\begin{center}IV.\end{center}


Among philosophical problems the problem of knowledge has more and more come to
occupy a foremost place. It is natural, then, to ask how
Darwin and the hypothesis whose most eminent representative he is, stand
to this problem.

Darwin started an hypothesis. But every hypothesis is won by inference
from certain presuppositions, and every inference is based on
the general principles of human thought. The evolution hypothesis
presupposes, then, human thought and its principles. And not only
the abstract logical principles, e.g. that the self-contradictory cannot be true, are thus presupposed. For the evolution
hypothesis purports to be not only a formal arrangement of phenomena,
but to express also the law of a real process. It supposes, then,
that the real data---all that in our knowledge which we do not produce
ourselves, but which we in the main simply receive---are subjected
to laws which are at least analogous to the logical relations of our
thoughts; in other words, it assumes the validity of the principle of
causality. If organic species could arise without cause there would
be no use in framing hypotheses. Only if we assume the principle of
causality, is there a problem to solve.

Though Darwinism has had a great influence on
% --- p. 216 ---
\setcounter{footnote}{0}%
\opage{216}philosophy considered as
a striving after a scientific view of the world, yet here is a point of
view---the epistemological---where philosophy is not only independent of natural science but
reaches beyond any result of natural science.

Perhaps the following consideration will be urged against this. The
powers and functions of organic beings only persist (perhaps also
only arise) when they correspond sufficiently to the conditions under
which the struggle of life is to go on. Human thought itself is, then, a
variation (or a mutation) which has been able to persist and to
survive. Is not, then, the problem of knowledge solved by the evolution
hypothesis? \textsc{Spencer} had given an affirmative answer to this question
before the appearance of ``The Origin of Species''. For the individual, he
said, there is an a priori, original, basis for all mental
life; but in the whole species all powers have arisen and developed in reciprocity with
external conditions. Knowledge is here considered from the practical
point of view, as a weapon in the struggle for life, as an organ
which has been continuously in use for countless generations. In recent years the
economic or pragmatic epistemology, as developed by \textsc{Avenarius} and \textsc{Mach}
in Germany, and by \textsc{James} in America, points in the same direction.
Science, it is said, only maintains those principles and presuppositions
which are necessary to the simplest and clearest orientation in
the world of experience. All assumptions which cannot be applied to
experience and to practical work, will successively be eliminated as superfluous.

In these views a striking and important application is made of the idea
of struggle for life to the development of human thought. It is true that thought must,
as all other things in the world, struggle for life. But this whole
consideration belongs to the biology or psychology of knowledge, not to the theory of
knowledge, whose task is not to examine the actual, historical origin and development of knowledge, but to discuss its validity. Every
% --- p. 217 ---
\setcounter{footnote}{0}%
\opage{217}hypothesis to explain the historical
origin and development of knowledge is itself an \textit{object} for the theory of
knowledge, because it works with the forms and principles
of human thought. We cannot go further back than these forms and
principles. It is the aim of epistemology to ascertain which these forms and principles are that constitute the ultimate presuppositions of all proof, and for
which therefore no further reason can be given. --- In my ``Psychology'' (VII C, 3), and already
in my book ``Den engelske Filosofi i vor Tid'' (1874), I have
criticized the evolutionistic treatment of the problem of
knowledge.

But there is another side of the problem which is, perhaps, of more
importance and which epistemology generally overlooks. If new variations
can arise, not only in organic but perhaps also in inorganic nature, new
tasks are placed before the human mind. The question is, then, if it has
forms in which there is room for the new matter? We are here touching
a presupposition which the great master of epistemology did not bring to
light. Kant supposed confidently that no other matter of knowledge could
stream forth from the dark source which he called ``the thing-in-itself,''
than such as could be synthesized in our existing forms of knowledge.
He mentions the possibility of other forms than the human, and warns us
against the dogmatic assumption that the human conception of existence
should be absolutely adequate. But he seems to be quite sure that the
``thing-in-itself'' works constantly, and consequently always gives us only
what our powers can master. This assumption was a consequence of Kant's
rationalistic tendency, but one for which no warrant can be given.
Evolution and rational connection stand in a constant opposition to each other. For the new, a place must first be sought in the world of thought, and can it always be found?
Evolutionism and systematism are opposing tendencies which can never
be absolutely harmonized one with the other. Evolution may at any time
break the forms which the system constructs.

Darwin himself felt a great difference in looking at variation as an
evolutionist and as a
% --- p. 218 ---
\setcounter{footnote}{0}%
\opage{218}systematist. When he was working at his evolution
theory, he was very glad to find variations; but they were a hindrance
to him when he worked as a systematist, in preparing his work on
Cirripedia. He says in a letter: ``I had thought the same parts of the
same species more resemble (than they do anyhow in Cirripedia) objects
cast in the same mould. Systematic work would be easy were it not
for this confounded variation, which, however, is pleasant to me as a
speculatist, though odious to me as a systematist.'' --- Darwin could indeed be angry with variations even as
an evolutionist; but then only because he could not explain their origin, not
because he could not classify them. ``If,'' he says in another letter, ``as I must think, external
conditions produce little direct effect, what the devil determines each
particular variation?'' ---

What Darwin experienced here in
his particular domain holds good of all knowledge. All knowledge is
systematic, in so far as it strives to put phenomena in quite definite
relations, one to another. But the systematization can never be
complete. And here Darwin has contributed much to widen the world for
us. He has shown us forces and tendencies in nature which make absolute
systems impossible, at the same time that they give us new objects and
problems. There is still a place for what Lessing called ``the unceasing
striving after truth,'' while ``absolute truth'' is unattainable so long as life and experience are going on. ---

There is here a special remark to be made. As we have seen above, recent
research has shown that natural selection or the struggle for life does not
explain the origin of species. \textsc{Hugo de Vries} distinguishes between partial
and embryonal variations, or between variations and mutations, only the
last-named being heritable, and therefore of importance for the origin
of new species. But the existence of variations is not only of interest
for the problem of the origin of species; it has also a more general
interest. An individual
% --- p. 219 ---
\setcounter{footnote}{0}%
\opage{219}does not lose its importance for knowledge,
because its qualities are not heritable. In higher
beings at least, individual peculiarities will become more and more
independent objects of interest. Knowledge takes account of the
biographies not only of species, but also of individuals\footnote[1]{The new science of Ecology (the theory of plant communities, their economy and their relation to their surroundings) occupies an intermediate position between the biography of species and the biography of individuals. Cf. \textit{Warming: Om Planterigets Livsformer} (University program 1908) p. 76, and \textit{Drude's} and \textit{Robinson's} lectures on Ecology at the Congress in St. Louis. (Congress of Arts and Science). St. Louis 1904 V p. 177---206.}. It seeks to
find the law of development of each single individual. As Leibniz said long ago,
individuality consists precisely in the law of the changes of a being (la loi de changement fait l'individualité de chaque
substance). Here is a world
which is almost new for science, which till now has mainly occupied
itself with general laws and forms. But general laws and forms are ultimately only means
to understand the individual phenomena, in whose nature and history
a manifold of laws and forms always cooperate\footnote[2]{See on this my work \textit{Den menneskelige Tanke, dens Former og dens Opgaver} p. 222---226 and p. 338---343. (Later note.)}. The importance of this
remark will appear in the sequel.

\begin{center}V.\end{center}


To many people the Darwinian theory of natural selection or struggle
for existence seemed to change the whole conception of life, and
particularly all the conditions on which the validity of ethical ideas
depends. If only that has persistence which can be adapted to a given
condition, what will then be the fate of our ideals, of our standards
of good and evil? Blind force seems to reign, and the only thing that
counts seems to be the most heedless use of power. Darwinism, it was
said, has proclaimed brutality. No other difference seems permanent save
that between the sound, powerful and happy on the one
% --- p. 220 ---
\setcounter{footnote}{0}%
\opage{220}side, the sick,
feeble and unhappy on the other; and every attempt to alleviate this
difference seems to lead to general enervation. Some of those who
interpreted Darwinism in this manner felt an aesthetic delight in
contemplating the heedlessness and energy of the great struggle for
existence and anticipated the realization of a higher human type as the
outcome of it: so \textsc{Nietzsche} (at least in some of his writings). Others, such as \textsc{Dühring} and \textsc{Kropotkin}, recognizing the
same consequences in Darwinism, regarded precisely this as one of the strongest
objections against it.

This interpretation of Darwinism was frequent in the interval between
the two main works of Darwin---``The Origin of Species'' (1859) and ``The Descent
of Man'' (1871). But even during this interval it was evident to an attentive
reader that Darwin himself did not found his standard of good and evil
on the features of the life of nature he had emphasized so strongly. He
did not justify the ways along which nature reached its ends; he only
pointed them out. The ``real'' was not to him, as to Hegel, one with
the ``rational.'' Darwin has, indeed, by his whole conception of nature,
rendered a great service to ethics in making the difference between the
life of nature and the ethical life appear in so strong a light. The
ethical problem could now be stated in a sharper form than before.

This was not the first time that the idea of the struggle for life
was put in relation to the ethical problem. In the seventeenth century
\textsc{Thomas Hobbes} gave the first impulse to the whole modern discussion of
ethical principles in his theory of the war of all against all (bellum omnium contra omnes). Men, he
taught, are in the state of nature enemies one of another, and they live
either in fright or in the glory of power. But it was not the opinion
of Hobbes that this made ethics impossible. On the contrary, he found a
standard for virtue and vice in the fact that some qualities and actions
have a tendency to bring us out of the state of war and to secure
peace, while other qualities have a contrary tendency. Later even
\textsc{Immanuel Kant's} ideal ethics had---so far as
% --- p. 221 ---
\setcounter{footnote}{0}%
\opage{221}can be seen---a
similar origin. Shortly before the publication of his fundamental work on ethics,
Kant wrote his ``Idee zu einer allgemeinen Weltgeschichte'' (1784),
where---in a way which reminds us of Hobbes, and is prophetic of
Darwin---he describes the forward-driving power of struggle in the human
world. It is here as with the struggle of the trees for light and air,
through which they compete with one another in height. The unrest and anxiety of struggle in the human world can only be allayed by an ordinance which gives everyone his full
liberty so far as it can coexist with the equal liberty of others. And such an
ordinance is precisely what is acknowledged and proclaimed in the
moral law which Kant, in the year after that essay, put forward in his ``Grundlegung zur Metaphysik der Sitten'' (1785): the rule one follows in one's action must be able to hold for everyone in the same position,---and in a more positive form: every man is to be treated as an end, never merely as a means\footnote[1]{Cf. \textit{Den nyere Filosofis Historie}\textsuperscript{2}. II 72---76.}. --- Similarly this whole line of thought may be traced in \textsc{Hegel}, though
with him it has acquired a more abstract character. And
in \textsc{Schopenhauer's} theory of the blind impulse of self-preservation, which he calls ``the will to live,'' and of its abrogation by
the ethical feeling, which is founded on sympathy, we have a
more individualistic form of the same idea.

It was, then, not entirely a foreign point of view which Darwin
introduced into ethical thought, even if we take no account of the
poetical character of the word ``struggle'' or of the more direct
adaptation, through the use and non-use of organs, which Darwin also
emphasized. And he has (in the third chapter of ``The Descent of Man'') given a special discussion of the origin
of the ethical consciousness. The characteristic expression of this
consciousness he found, just as Kant did, in the idea of an ``ought''; it was
precisely the origin of this idea which should be explained. His hypothesis
was that the ethical ``ought'' has its origin in the social and parental (maternal and paternal)
instincts, which, as well as other instincts (e.g. the
% --- p. 222 ---
\setcounter{footnote}{0}%
\opage{222}instinct of
self-preservation), lie deeper than the feelings of pleasure and pain and are to a certain degree independent of them. In many species,
not least in the human species, these instincts are fostered by natural
selection; and when the powers of memory and comparison are developed,
so that single acts can be valued according to the claims of the deep
social instinct,---then consciousness of duty and remorse are possible.
Blind instinct has developed to conscious ethical will.

Darwin, as a moral philosopher belongs to the school
that was founded by \textsc{Shaftesbury}, and was afterwards represented by
\textsc{Hutcheson}, \textsc{Hume} and \textsc{Adam Smith}, in the nineteenth century---besides Darwin---by \textsc{Comte} and \textsc{Spencer}. His merit is, first,
that he has given this tendency of thought a biological foundation, and, secondly,
that he has stamped on it a doughty character in showing that ethical
ideas and sentiments are forces which are at work in
the struggle for life.

There are still many questions to solve. In the first place, the ethical
development within the human species contains many features still unexplained.
This side of the ethical problem has in most recent times been illuminated by the works of \textsc{Westermarck} and
\textsc{Hobhouse}. But, secondly, we are confronted by the great problem whether after all
a genetic-historical theory can be of decisive importance here. To
every consequent ethical line of thought there is a standard of valuation, a
primordial value which determines the single ethical judgments as their
last presupposition. The ``rightness'' of this basis, the ``value''
of the primordial value, can as little be justified as the ``rationality'' of the
last presuppositions of a reason (the logical principles). There is here revealed a possibility of ethical
skepticism which evolutionistic ethics (as well as intuitive or
rationalistic ethics) has overlooked. No demonstration can show that the
results of the ethical development are definitive and universal. We
meet here again with the opposition of systematization and
evolution. There will, I think, always be an open question here, though
comparative ethics, of which we have so far only the first attempts, can
do much to throw light on it. ---

% --- p. 223 ---
\setcounter{footnote}{0}%
\opage{223}It would carry us too far to discuss all the philosophical works
on ethics, which have been influenced directly or indirectly by
evolutionism. I may here refer on this to C. M. \textsc{Williams'} book
\textit{A Review of the Systems of Ethics founded on the theory of Evolution}
(New York and London 1893), in which, besides Darwin, the following
authors are reviewed: Wallace, Häckel, Spencer, Fiske, Rolph, Barratt,
Stephen, Carneri, Høffding, Gizycki, Rée. As instructive works which
criticize evolutionism, may be cited \textsc{Guyau's} \textit{La Morale Anglaise
contemporaine} (Paris 1879) and \textsc{Sorley's} \textit{Ethics of Naturalism}
(1885, 2nd ed. 1904). --- I will here only mention some interesting contributions to ethical discussion
which can be found in Darwinism besides the idea of struggle for life. ---

The attention which Darwin has directed to variations has opened our
eyes to the individual differences in human nature as well as in nature generally.
There is here a fact of great importance for ethical thought, no matter
from what ultimate premise it starts. Only from a very abstract point of
view can different individuals be treated in the same manner. And yet the most
eminent ethical thinkers of the preceding period, Jeremy Bentham and Immanuel Kant,
who discussed ethical questions from very opposite standpoints, agreed
in regarding all men as equal in respect of ethical endowment. In regard
to Bentham, \textsc{Leslie Stephen} remarks (in his book \textit{English Literature and Society in the 18. century} p. 187): ``He is determined to be thoroughly
empirical, to take men as he found them. But his utilitarianism supposed
that men's views of happiness and utility were uniform and clear, and
that all that was wanted was to show them the means by which their ends
could be reached.'' And Kant supposed that every man
would find the categorical imperative in his consciousness, when he
came to sober reflection, and that all would have the same qualifications
to follow it. But if continual variations, great or small, are going on
in human nature, it is the duty of ethics to make allowance for them,
both in making claims, and in
% --- p. 224 ---
\setcounter{footnote}{0}%
\opage{224}valuing what is done. A new set of ethical
problems have their origin here. In my essay on the law of relativity in
ethics \textit{(Etiske Undersøgelser} 1891) I have treated this point more fully.

It is an interesting fact that \textsc{Stuart Mill's} book On Liberty
appeared in the same year as ``The Origin of Species''. Though Mill agreed
with Bentham about the original equality of all men's endowments, he
regarded individual differences as a necessary result of physical and
social influences, and he claimed that free play shall be allowed to
differences of character so far as is possible without injury to other
men. It is a condition of individual happiness as well as of individual and social progress that a man's
mode of action should be determined by his own character and not
by tradition and custom. --- Darwin's doctrine of the importance of variations in nature furnished a biological foundation for this
view.

And here we have reached a point of view from which the criticism,
which in recent years has often been directed against Darwin,---that
variations are not always of importance in the struggle for life,---is of
no weight. From an ethical standpoint, and particularly from the
ethical standpoint of Darwin himself, it is a duty to foster individual
variations that can be valuable, even though they can neither be of
service for physical preservation nor be physically inherited. The
distinction between variation and mutation is here without importance.
It is quite natural that biologists should be particularly interested in
such variations as can be inherited and produce new species. But in the
human world there is not only a physical, but also a mental and social
heredity. When an ideal human character has taken form, then there
is shaped a type, which through imitation and influence can become an
important factor in subsequent development, even if it cannot form a
species in the biological sense of the word. Spiritually strong
men very often succumb in the physical struggle for life; but they can
nevertheless be victorious through the typical influence they exert,
perhaps on very distant generations, if the remembrance of them is kept
alive, be it in
% --- p. 225 ---
\setcounter{footnote}{0}%
\opage{225}historical or in legendary form. Their very tragic failure
can show that a type has taken form which is maintained at all risks,
a standard of life which is adhered to in spite of the strongest
opposition. The question ``to be or not to be'' can be put from
very different levels of being. It has too often been considered a
consequence of Darwinism that this question is only to be put from
the standpoint of purely physical self-preservation. When a stage is reached, where ideal (ethical,
intellectual, aesthetic) interests are concerned, the struggle for life
is a struggle for the preservation of these interests. The giving up of them, of the
higher claims on life which they make, is then a sort of death. There is a
spiritual death as well as a physical one.

\begin{center}VI.\end{center}


The Socratic character of Darwin's mind appears in his wariness in
drawing the last consequences of his doctrine, in contrast both with the
audacious theories of so many of his followers and with the consequences
which his antagonists were busy in drawing. Though he, as we have seen,
saw from the beginning that his hypothesis would occasion ``a whole of
metaphysics,'' he was himself very reserved as to the ultimate questions,
and his answers to such questions were extorted from him.

As to the question of optimism and pessimism, Darwin held that though
pain and suffering were very often the ways by which animals were led
to pursue that course of action which is most beneficial to the species,
yet pleasurable feelings were the most habitual guides. ``We see this,'' he says in a letter, ``in
the pleasure from exertion, even occasionally from great exertion of the
body or mind, in the pleasure of our daily meals, and especially in the
pleasure derived from sociability, and from loving our families.'' But
there was to him so much suffering in the world that it was a strong
argument against the existence of an intelligent First Cause.

% --- p. 226 ---
\setcounter{footnote}{0}%
\opage{226}It seems to me that Darwin was not so clear on another question, that
of the relation between improvement and adaptation. In a letter to Lyell (25 Oct. 1859) he says:
``When you contrast natural selection and improvement, you seem always
to overlook that every step in the natural selection of each species
implies improvement in that species in relation to its condition of
life . . . Improvement implies, I suppose, each form obtaining many parts
or organs, all excellently adapted for their functions . . . All this
seems to me quite compatible with certain forms fitted for simple
conditions, remaining unaltered, or being degraded.'' But the great question is, if the conditions of life will in
the long run favor ``improvement'' in the sense of differentiation
(or, more exactly, harmony of differentiation and integration). Many beings are best
adapted to their conditions of life if they have few organs and few
necessities. And beings which have manifold organs and other necessities than the purely elementary ones will, if the standard is merely adaptation to the given conditions, fare best by being ``degraded.'' Pessimism would not only be the necessary consequence, if suffering
outweighed happiness, but also if the most elementary forms of happiness
were predominant, or if there were a constant tendency to lower the standard
of life, so that inertia or stable
equilibrium were the result. There are animals which are very highly differentiated and
active in their young state, but later lose their complex organization
and concentrate themselves on the one function of nutrition. In the
human world analogies to this sort of adaptation are not wanting. Young
idealists very often end as old Philistines. There is, then, really a definite difference between progress
and adaptation.

Another question of great importance in respect to human evolution
is, whether there will be always a possibility for the existence of
an impulse to progress, an impulse to make great claims on life, to be
active and to alter the conditions of life instead of adapting to them
in a passive manner. Many people do not develop because they have
too few
% --- p. 227 ---
\setcounter{footnote}{0}%
\opage{227}necessities, and because they have no power to imagine other
conditions of life than those under which they live.

In his remark
on the pleasure from exertion Darwin has a point of contact with
the practical idealism of former times---with the ideas of Lessing and
Goethe, of Condorcet and Fichte. The continual striving which was
the condition of salvation to Faust's soul, is also the condition
of salvation to mankind. Salvation here means adaptation to the highest, most comprehensive and richest conditions of life. There is a holy fire which we ought to keep
burning, if adaptation is really to be improvement. If---as I have
tried to show in my Philosophy of Religion---the innermost core of all
religion is faith in the persistence of value in the world,---and if
the highest values express themselves in an Excelsior!---then the
capital point is, that this cry should always be heard and followed. We
have here a corollary of the theory of evolution in its application to
human life. ---

When Darwin declared himself an agnostic, i.e. when he held that no well-founded answer could be given to the question of the origin and final end of existence, this was not only because he could not
harmonize the large amount of suffering in the world with the idea of a
God as its first cause, but also for theoretical reasons. If we admit
a first cause, the mind still craves to know whence it came and how it
arose. And, further, it was clear to him that, however great the connection and purposiveness one may find in the world, it was just as impossible to derive it from an end as to assume chance or blind force as its ground. Darwin's position is similar to that which Kant took in the ``Kritik der Urtheilskraft''. For human thought there are only two possibilities: design---or blind force. But neither mechanism
nor teleology can give any answer to ultimate questions.
The universe, and especially the organic life in it, can neither be
explained as a mere combination of absolute elements nor as the effect
of a constructing thought. Darwin had learned by experience in his special domain how difficult it is to set up an Either-Or, because we cannot know all the possibilities beforehand. He came then to the
% --- p. 228 ---
\setcounter{footnote}{0}%
\opage{228}conclusion, which Spinoza and Kant, each in his own way, had already reached, that the oppositions and distinctions which our experience
and our thought present, cannot safely be regarded as valid for existence in itself.

Darwin reminds us of Kant also in that he found his last stronghold in the conviction
that man has something to do, even if he cannot solve all enigmas. In a letter to a Dutch student who wished to know his theological opinions, he wrote (1873): ``The
safest conclusion seems to me that the whole subject is beyond the scope
of man's intellect; but man can do his duty.''

Is this the last word of human thought? Does not the possibility,
that man can do his duty as the relatively highly developed being that he is, necessarily suppose that the conditions of life allow of
ethical striving, so that there is a certain harmony
between cosmic order and human ideals?
--- Darwin himself has sought to show how
the consciousness of duty can arise as a natural result of evolution.
There are, then, lines of evolution which have their end in ethical
idealism, in a kingdom of values, which must struggle for life as all
things in the world must do, but a kingdom which has its firm foundation
in reality.

\end{document}
